% Utilisation d’un logiciel de traitement de texte — Informatique 3ème — Cours signé OnBuch+
% Programme officiel MINESEC. Compiler avec : tectonic N05-utilisation-logiciel-traitement-texte.tex
\newcommand{\DOCMATIERE}{Informatique}
\newcommand{\DOCNIVEAU}{3ème}
\newcommand{\DOCMODULE}{Informatique – classe de 3ème}
\newcommand{\DOCLECON}{N05}
\newcommand{\DOCTITRE}{Utilisation d’un logiciel de traitement de texte}
% OnBuch+ — préambule « cours signés » ENRICHI (compilé avec tectonic / XeLaTeX)
% Système visuel « Pop » d'OnBuch+ : fond crème (jamais blanc), bordures encre
% épaisses, ombres « dures » décalées, titres Archivo Black en capitales.
% Version MATHÉMATIQUES : graphes (pgfplots/tikz), boîtes pédagogiques
% (définition, propriété, méthode, attention, le savais-tu, exercice résolu,
% exercice, TP, bilan), page de garde et signature OnBuch+ ancrée en bas.
%
% Macros attendues dans main.tex AVANT \input{preamble} :
%   \DOCMATIERE  \DOCNIVEAU  \DOCMODULE  \DOCLECON (numéro)  \DOCTITRE
\documentclass[11pt]{article}

\usepackage{fontspec}
\usepackage[french]{babel}
\usepackage[a4paper,margin=2cm,top=3.4cm,bottom=2.8cm,headheight=1.4cm,footskip=1.3cm]{geometry}
\usepackage{amsmath,amssymb,amsfonts}
\usepackage{mathtools} % \xmapsto, \coloneqq... souvent utilisés par les modèles
\usepackage{cancel} % simplifications barrées (bilans de forces)
% \abs{z} (module, valeur absolue) et \norm{u} (norme), utilisés par les modèles
\providecommand{\abs}{}\let\abs\relax\DeclarePairedDelimiter{\abs}{\lvert}{\rvert}
\providecommand{\norm}{}\let\norm\relax\DeclarePairedDelimiter{\norm}{\lVert}{\rVert}
\usepackage[table]{xcolor} % [table] : fournit aussi \rowcolors (alternance de lignes)
\usepackage{pagecolor}
\usepackage{tikz}
\usetikzlibrary{arrows.meta,positioning,calc,shapes.geometric,shapes.misc,shapes.arrows,shapes.symbols,decorations.pathmorphing,decorations.pathreplacing,decorations.markings,patterns,shadows,mindmap,trees,fit,backgrounds,matrix,intersections,angles,quotes}
\usepackage{pgfplots}
\usepackage[version=4]{mhchem} % équations nucléaires \ce{^{235}_{92}U}
\usepackage[european,siunitx]{circuitikz} % schémas électriques
\pgfplotsset{compat=1.18}
\usepgfplotslibrary{fillbetween,groupplots} % aires entre courbes (calcul intégral)
\usepackage{enumitem}
\usepackage{fancyhdr}
% [most] charge toutes les bibliothèques tcolorbox (listings, minted, raster,
% documentation...) et gonfle inutilement la mémoire TeX sur un document long
% (~300 boîtes) : on ne charge que ce qui est réellement utilisé.
\usepackage[skins,breakable]{tcolorbox}
\usepackage{graphicx}
\usepackage{colortbl}
\usepackage{booktabs}
\usepackage{tabularx}
\usepackage{diagbox} % case d'en-tête diagonale des tableaux à double entrée
% Colonnes extensibles centrée / gauche / droite (C, L, R), souvent utilisées par les modèles
\newcolumntype{C}{>{\centering\arraybackslash}X}
\newcolumntype{L}{>{\raggedright\arraybackslash}X}
\newcolumntype{R}{>{\raggedleft\arraybackslash}X}
\usepackage{array}
\usepackage{multirow}
\usepackage{multicol}
\usepackage{siunitx}
% parse-numbers=false : accepte aussi \SI{2,5 \times 10^{-3}}{...}
\sisetup{locale=FR,output-decimal-marker={,},group-minimum-digits=4,parse-numbers=false}
\DeclareSIUnit\ohm{\ensuremath{\symup{\Omega}}} % U+2126 absent de la police
\DeclareSIUnit\division{div} % réglages d'oscilloscope (ms/div, V/div)
\DeclareSIUnit\tr{tr} % tours (vitesse de rotation, ex. \SI{1500}{\tr\per\minute})
\DeclareSIUnit\bit{bit}
\DeclareSIUnit\byte{o} % octet
\DeclareSIUnit\inch{po} % pouce
\DeclareSIUnit\ppp{ppp} % points par pouce (résolution d'image)
\usepackage{qrcode}
% Blocs de code source (PHP, C, HTML, JavaScript, SQL) — spécifique à la
% matière Informatique : listings + habillage aux couleurs OnBuch+ (voir le
% style \lstset ci-dessous, à côté des autres boîtes pédagogiques).
\usepackage{listings}
% \degree : commande fréquemment utilisée par les modèles pour un angle ou une
% température en dehors de \SI{}{} (non fournie par les paquets chargés).
\providecommand{\degree}{^{\circ}}
% \qed (fin de démonstration), utilisé par les modèles sans amsthm
\providecommand{\qed}{\hfill$\square$}
% \barre : double barre des valeurs interdites dans un tableau de variations (tkz-tab)
\providecommand{\barre}{\ensuremath{\|}}
% environnement proof (amsthm non chargé) utilisé par les modèles
\ifdefined\proof\else\newenvironment{proof}[1][Démonstration]{\par\noindent\textit{#1.}\ }{\qed\par}\fi
\usepackage{lastpage}
\usepackage{hyperref}
\hypersetup{hidelinks}

% ============================================================
% Palette « Pop » OnBuch+ — mêmes hex que la classe P (plus_ui.dart)
% ============================================================
\definecolor{poporange}{HTML}{F59321}
\definecolor{popdark}{HTML}{E07A0C}
\definecolor{popcream}{HTML}{FFF6E8}
\definecolor{popink}{HTML}{1C1714}
\definecolor{poppurple}{HTML}{7A5AE0}
\definecolor{popgreen}{HTML}{2E9E6A}
\definecolor{poppink}{HTML}{E0457B}
\definecolor{popblue}{HTML}{2F7DE1}
\definecolor{popgold}{HTML}{C98A12}
\definecolor{popmuted}{HTML}{8A7F76}
\definecolor{popline}{HTML}{EADFD2}
\definecolor{popgray}{HTML}{8A7F76}
\colorlet{popgrey}{popgray}
% Alias tolérants (le modèle nomme parfois les couleurs différemment)
\colorlet{popwhite}{white}
\colorlet{popblack}{popink}
\colorlet{popred}{poppink}
\colorlet{popyellow}{popgold}
% Teintes claires pour fonds de tableaux / zones
\colorlet{popblueL}{popblue!10!white}
\colorlet{poporangeL}{poporange!14!white}
\colorlet{popgreenL}{popgreen!12!white}
\colorlet{poppurpleL}{poppurple!12!white}
\colorlet{poppinkL}{poppink!12!white}
\colorlet{popgoldL}{popgold!14!white}

\pagecolor{popcream}

% ============================================================
% Typographie
% ============================================================
\defaultfontfeatures{Ligatures=TeX}
\setmainfont{PlusJakartaSans-Medium.ttf}[Path=../../../fonts/,BoldFont=PlusJakartaSans-Bold.ttf]
% Maths en sans-serif (Fira Math) avec chiffres et lettres droites de Plus
% Jakarta, pour que les formules se fondent dans le texte.
\usepackage[math-style=ISO,bold-style=ISO]{unicode-math}
\setmathfont{FiraMath-Regular.otf}[Path=../../../fonts/]
\setmathfont{PlusJakartaSans-Medium.ttf}[Path=../../../fonts/,range={up/num,up/{Latin,latin},\mathup}]
\usepackage{icomma} % virgule décimale sans espace en mode math
% Caractères mathématiques Unicode absents des polices de texte (Plus Jakarta,
% Archivo Black) : rendus par leur équivalent mathématique, en texte comme en
% formule — sinon ils s'affichent en carré vide (ex. « Arithmétique dans ℤ »).
\usepackage{newunicodechar}
\newunicodechar{ℝ}{\ensuremath{\mathbb{R}}}
\newunicodechar{ℤ}{\ensuremath{\mathbb{Z}}}
\newunicodechar{ℂ}{\ensuremath{\mathbb{C}}}
\newunicodechar{ℕ}{\ensuremath{\mathbb{N}}}
\newunicodechar{ℚ}{\ensuremath{\mathbb{Q}}}
\newunicodechar{𝔻}{\ensuremath{\mathbb{D}}}
\newunicodechar{α}{\ensuremath{\alpha}}
\newunicodechar{β}{\ensuremath{\beta}}
\newunicodechar{γ}{\ensuremath{\gamma}}
\newunicodechar{δ}{\ensuremath{\delta}}
\newunicodechar{ε}{\ensuremath{\varepsilon}}
\newunicodechar{θ}{\ensuremath{\theta}}
\newunicodechar{λ}{\ensuremath{\lambda}}
\newunicodechar{μ}{\ensuremath{\mu}}
\newunicodechar{σ}{\ensuremath{\sigma}}
\newunicodechar{τ}{\ensuremath{\tau}}
\newunicodechar{φ}{\ensuremath{\varphi}}
\newunicodechar{ψ}{\ensuremath{\psi}}
\newunicodechar{ω}{\ensuremath{\omega}}
\newunicodechar{Σ}{\ensuremath{\Sigma}}
% majuscules produites par \MakeUppercase dans les titres (α-aminés -> Α)
\newunicodechar{Α}{\ensuremath{\alpha}}
\newunicodechar{Β}{\ensuremath{\beta}}
\newunicodechar{Γ}{\ensuremath{\Gamma}}
\newunicodechar{Δ}{\ensuremath{\Delta}}
\newunicodechar{Ε}{\ensuremath{\varepsilon}}
\newunicodechar{Θ}{\ensuremath{\Theta}}
\newunicodechar{Λ}{\ensuremath{\Lambda}}
\newunicodechar{Μ}{\ensuremath{\mu}}
\newunicodechar{Τ}{\ensuremath{\tau}}
\newunicodechar{Φ}{\ensuremath{\Phi}}
\newunicodechar{Ψ}{\ensuremath{\Psi}}
\newunicodechar{Ω}{\ensuremath{\Omega}}
\newunicodechar{↦}{\ensuremath{\mapsto}}
\newunicodechar{∈}{\ensuremath{\in}}
\newunicodechar{∉}{\ensuremath{\notin}}
\newunicodechar{∧}{\ensuremath{\wedge}}
\newunicodechar{⇔}{\ensuremath{\Leftrightarrow}}
\newunicodechar{⟺}{\ensuremath{\Longleftrightarrow}}
\newunicodechar{⇒}{\ensuremath{\Rightarrow}}
\newunicodechar{⟹}{\ensuremath{\Longrightarrow}}
\newunicodechar{∘}{\ensuremath{\circ}}
\newunicodechar{∪}{\ensuremath{\cup}}
\newunicodechar{∩}{\ensuremath{\cap}}
\newunicodechar{≡}{\ensuremath{\equiv}}
\newunicodechar{⊂}{\ensuremath{\subset}}
\newunicodechar{⊥}{\ensuremath{\perp}}
\newunicodechar{∃}{\ensuremath{\exists}}
\newunicodechar{∀}{\ensuremath{\forall}}
\newunicodechar{∛}{\ensuremath{\sqrt[3]{\,}}}
\newunicodechar{∜}{\ensuremath{\sqrt[4]{\,}}}
\newunicodechar{⃗}{\ensuremath{{}^{\to}}}
\newunicodechar{ⁿ}{\ifmmode^{n}\else\textsuperscript{\ensuremath{n}}\fi}
\newunicodechar{ˣ}{\ifmmode^{x}\else\textsuperscript{\ensuremath{x}}\fi}
\newunicodechar{ᵇ}{\ifmmode^{b}\else\textsuperscript{\ensuremath{b}}\fi}
\newunicodechar{ʳ}{\ifmmode^{r}\else\textsuperscript{\ensuremath{r}}\fi}
\newunicodechar{ᵘ}{\ifmmode^{u}\else\textsuperscript{\ensuremath{u}}\fi}
\newunicodechar{ʸ}{\ifmmode^{y}\else\textsuperscript{\ensuremath{y}}\fi}
\newunicodechar{ᵃ}{\ifmmode^{a}\else\textsuperscript{\ensuremath{a}}\fi}
\newunicodechar{ᵉ}{\ifmmode^{e}\else\textsuperscript{\ensuremath{e}}\fi}
\newunicodechar{ᵏ}{\ifmmode^{k}\else\textsuperscript{\ensuremath{k}}\fi}
\newunicodechar{ᵖ}{\ifmmode^{p}\else\textsuperscript{\ensuremath{p}}\fi}
\newunicodechar{ᴺ}{\ifmmode^{N}\else\textsuperscript{\ensuremath{N}}\fi}
\newunicodechar{ᶜ}{\ifmmode^{c}\else\textsuperscript{\ensuremath{c}}\fi}
\newunicodechar{ʰ}{\ifmmode^{h}\else\textsuperscript{\ensuremath{h}}\fi}
\newunicodechar{ᵠ}{\ifmmode^{q}\else\textsuperscript{\ensuremath{q}}\fi}
\newunicodechar{ₙ}{\ifmmode_{n}\else\textsubscript{\ensuremath{n}}\fi}
\newunicodechar{ₐ}{\ifmmode_{a}\else\textsubscript{\ensuremath{a}}\fi}
\newunicodechar{ᵢ}{\ifmmode_{i}\else\textsubscript{\ensuremath{i}}\fi}
\newunicodechar{ᵣ}{\ifmmode_{r}\else\textsubscript{\ensuremath{r}}\fi}
\newunicodechar{ₖ}{\ifmmode_{k}\else\textsubscript{\ensuremath{k}}\fi}
\newunicodechar{ᵦ}{\ifmmode_{\beta}\else\textsubscript{\ensuremath{\beta}}\fi}
\newunicodechar{ₓ}{\ifmmode_{x}\else\textsubscript{\ensuremath{x}}\fi}
\newfontfamily\popdisplay{ArchivoBlack-Regular.ttf}[Path=../../../fonts/]
\newfontfamily\poplogo{SpaceGrotesk-Bold.ttf}[Path=../../../fonts/]
\newfontfamily\popsemi{PlusJakartaSans-SemiBold.ttf}[Path=../../../fonts/]
\newfontfamily\popxbold{PlusJakartaSans-ExtraBold.ttf}[Path=../../../fonts/]
\newfontfamily\popxboldls{PlusJakartaSans-ExtraBold.ttf}[Path=../../../fonts/,LetterSpace=3.0]

\setlist[enumerate]{leftmargin=1.7em,itemsep=5pt,topsep=4pt}
\setlist[itemize]{leftmargin=1.7em,itemsep=4pt,topsep=4pt,label={\color{poporange}\textbullet}}
\setlength{\parindent}{0pt}
\setlength{\parskip}{5pt}
\renewcommand{\arraystretch}{1.25}

% pgfplots : style Pop par défaut
\pgfplotsset{
  every axis/.append style={axis line style={popink,line width=1pt},tick style={popink},
    grid style={popline,line width=0.5pt},label style={font=\small},tick label style={font=\footnotesize}},
  popcurve/.style={poporange,line width=1.8pt,no markers,smooth},
}
% Même style disponible dans un \draw TikZ simple (hors pgfplots)
\tikzset{popcurve/.style={poporange,line width=1.8pt}}
\tikzset{popgrid/.style={popline,very thin}}
\colorlet{popcurve}{poporange} % le nom du style est parfois utilisé comme couleur

% ============================================================
% Pill / squiggle
% ============================================================
\newcommand{\poppill}[3]{%
  \tikz[baseline=-0.55ex]\node[draw=popink,line width=1.2pt,rounded corners=2.6pt,%
    fill=#2,inner xsep=6.5pt,inner ysep=3pt]{%
    \popxboldls\fontsize{7}{7}\selectfont\color{#3}\MakeUppercase{#1}};%
}
\newcommand{\popsquiggle}[1]{%
  \tikz[baseline=-0.4ex]{
    \draw[#1,line width=2.6pt,line cap=round]
      plot [smooth] coordinates {(0,0.09) (0.34,0.01) (0.68,0.21) (1.02,0.01) (1.36,0.10)};
  }%
}
% Mot-clé surligné (marqueur orange sous le texte)
\newcommand{\cle}[1]{{\popxbold\color{popdark}#1}}

% ============================================================
% En-tête / pied
% ============================================================
\pagestyle{fancy}
\fancyhf{}
\renewcommand{\headrulewidth}{0pt}
\renewcommand{\footrulewidth}{0pt}
\lhead{\raisebox{-0.15ex}{\poplogo\fontsize{13}{13}\selectfont\color{popink}OnBuch\color{poporange}+}}
\rhead{\poppill{\DOCMATIERE\ \textbullet\ \DOCNIVEAU\ \textbullet\ Leçon \DOCLECON}{white}{popink}}
\lfoot{\popxbold\fontsize{7.6}{7.6}\selectfont\color{popmuted}\MakeUppercase{Cours signé OnBuch+}\ \textbullet\ {\popsemi\DOCTITRE}}
\rfoot{\tikz[baseline=-0.6ex]\node[rounded corners=8.5pt,fill=popink,minimum height=17pt,minimum width=17pt,inner xsep=5pt,inner sep=0]{\popxbold\fontsize{8}{8}\selectfont\color{popcream}\thepage\,/\,\pageref*{LastPage}};}

% ============================================================
% Titres
% ============================================================
\newcommand{\chaptitre}[2]{%
  \begin{tcolorbox}[enhanced,colback=poporange,colframe=popink,boxrule=2.6pt,arc=11pt,
      drop shadow=popink,left=15pt,right=15pt,top=12pt,bottom=11pt,before skip=3pt,after skip=18pt]
    \poppill{#2}{popink}{popcream}\par\vspace{9pt}
    {\raggedright\hyphenpenalty=10000\exhyphenpenalty=10000\popdisplay\fontsize{22}{24}\selectfont\color{popcream}\MakeUppercase{#1}\par}
  \end{tcolorbox}
}
% Section numérotée : pastille orange avec le numéro + titre Archivo
\newcounter{popsec}
\newcommand{\coursec}[1]{%
  \stepcounter{popsec}\setcounter{popsub}{0}%
  \par\vspace{16pt}\noindent
  \begin{minipage}{\linewidth}
  \tikz[baseline=-0.7ex]\node[draw=popink,line width=1.6pt,fill=poporange,rounded corners=4pt,minimum size=22pt,inner sep=0]{\popdisplay\fontsize{12}{12}\selectfont\color{popcream}\thepopsec};%
  \hspace{8pt}{\popdisplay\fontsize{15}{17}\selectfont\color{popink}\MakeUppercase{#1}}\par
  \vspace{4pt}\noindent{\color{poporange}\rule{36pt}{3.6pt}}
  \end{minipage}\par\nopagebreak\vspace{8pt}
}
\newcounter{popsub}
\newcommand{\courssub}[1]{%
  \stepcounter{popsub}%
  \par\vspace{9pt}\noindent{\popxbold\fontsize{11.5}{13}\selectfont\color{popink}{\color{poporange}\thepopsec.\thepopsub}\hspace{6pt}#1}\par\nopagebreak\vspace{4pt}
}

% ============================================================
% Boîtes pédagogiques — même recette : fond blanc, bordure épaisse colorée,
% ombre dure encre, badge plein. Titre optionnel [#1].
% ============================================================
\tcbset{popbase/.style={breakable,enhanced,colback=white,boxrule=2.3pt,arc=9pt,
  shadow={3pt}{-3pt}{0pt}{popink},left=14pt,right=14pt,top=11pt,bottom=11pt,before skip=11pt,after skip=15pt,
  fontupper=\popsemi\fontsize{10.6}{14.5}\selectfont\color{popink},
  fontlower=\popsemi\fontsize{10.6}{14.5}\selectfont\color{popink}}}
% Coche dessinée (le glyphe ✓ n'existe pas dans les polices du document)
\newcommand{\popcheck}[1]{\tikz[baseline=-0.5ex]\draw[#1,line width=1.6pt,line cap=round,line join=round](-0.13,0.0)--(-0.04,-0.09)--(0.13,0.1);}
\providecommand{\coche}{\popcheck{popgreen}}
\providecommand{\resultat}[1]{\par\smallskip\noindent\fcolorbox{popgreen}{popgreenL}{\parbox{\dimexpr\linewidth-2\fboxsep-2\fboxrule\relax}{\textbf{Résultat.} #1}}\par\smallskip}
\newcommand{\popboxhead}[3]{\poppill{#1}{#2}{white}\ifx\relax#3\relax\else\hspace{6pt}{\popxbold\fontsize{10.6}{12}\selectfont\color{popink}#3}\fi\par\vspace{7pt}}

\newtcolorbox{aretenir}[1][]{popbase,colframe=popgreen,before upper={\popboxhead{À retenir}{popgreen}{#1}}}
\newtcolorbox{exemplebox}[1][]{popbase,colframe=poporange,before upper={\popboxhead{Exemple}{poporange}{#1}}}
\newtcolorbox{definition}[1][]{popbase,colframe=popblue,before upper={\popboxhead{Définition}{popblue}{#1}}}
\newtcolorbox{propriete}[1][]{popbase,colframe=poppurple,before upper={\popboxhead{Propriété}{poppurple}{#1}}}
\newtcolorbox{methode}[1][]{popbase,colframe=popgold,before upper={\popboxhead{Méthode}{popgold}{#1}}}
\newtcolorbox{attention}[1][]{popbase,colframe=poppink,before upper={\popboxhead{Attention}{poppink}{#1}}}
\newtcolorbox{experience}[1][]{popbase,colframe=popblue,colback=popblueL,before upper={\popboxhead{Expérience}{popblue}{#1}}}
% « Le savais-tu ? » : carte violette pleine (contexte, culture, Cameroun)
\newtcolorbox{savaistu}[1][]{popbase,colback=poppurple,colframe=popink,
  fontupper=\popsemi\fontsize{10.4}{14.2}\selectfont\color{white},
  before upper={\poppill{Le savais-tu\,?}{white}{poppurple}\ifx\relax#1\relax\else\hspace{6pt}{\popxbold\color{white}#1}\fi\par\vspace{7pt}}}
% Exercice résolu : énoncé en haut, \tcblower puis la solution sur fond crème
\newcounter{popexo}\newcounter{popexr}
\newtcolorbox{exoresolu}[1][]{popbase,colframe=poporange,colbacklower=popcream,
  segmentation style={popink,line width=1.2pt,dashed},
  before upper={\stepcounter{popexr}\popboxhead{Exercice résolu \thepopexr}{poporange}{#1}},
  before lower={\poppill{Solution}{popink}{popcream}\par\vspace{6pt}}}
% Exercice (non résolu en place) — la correction est dans la partie Corrigés
\newtcolorbox{exercice}[1][]{popbase,colframe=popink,
  before upper={\stepcounter{popexo}\popboxhead{Exercice \thepopexo}{popink}{#1}}}
% Corrigé
\newtcolorbox{corrige}[1][]{popbase,colframe=popgreen,colback=popgreenL,
  before upper={\popboxhead{Corrigé}{popgreen}{#1}}}
% Objectifs / prérequis / bilan
\newtcolorbox{objectifs}{popbase,colframe=popink,colback=poporangeL,
  before upper={\popboxhead{Objectifs de la leçon}{popink}{}}}
\newtcolorbox{prerequis}{popbase,colframe=popink,colback=popblueL,
  before upper={\popboxhead{Prérequis}{popblue}{}}}
\newtcolorbox{bilan}{popbase,colframe=popink,colback=white,boxrule=2.8pt,
  before upper={\popboxhead{Fiche bilan}{poporange}{L'essentiel à connaître pour l'examen}}}
% Figure encadrée avec légende
\newtcolorbox{popfigure}[1][]{enhanced,breakable=false,colback=white,colframe=popline,boxrule=1.4pt,arc=8pt,
  left=10pt,right=10pt,top=10pt,bottom=8pt,before skip=10pt,after skip=12pt,halign=center,
  lower separated=false,halign lower=center,
  fontlower=\popsemi\fontsize{9}{11.5}\selectfont\color{popmuted},
  #1}
\newcommand{\legende}[1]{\par\vspace{6pt}{\popsemi\fontsize{9}{11.5}\selectfont\color{popmuted}#1}}

% ============================================================
% Bloc de code source — \begin{lstlisting}[language=Php]...\end{lstlisting}
% (ou language=C, HTML, JavaScript, SQL ; option omise = pseudo-code brut).
% Même recette visuelle que les figures (\popfigure) : cadre encre épais,
% fond clair (popcream), légende possible juste après via \legende{...}
% comme pour une figure (listings ne sait pas arrondir ses coins : on garde
% un cadre rectangulaire, cohérent avec les autres tableaux du document).
% CONTRAT : les modèles ne doivent utiliser QUE \begin{lstlisting}[...] tel
% quel (pas de \begin{tcblisting}, pas de \verb pour un bloc multi-lignes).
% ============================================================
\lstdefinestyle{poplst}{
  backgroundcolor=\color{popcream},
  basicstyle=\popsemi\fontsize{8.6}{11}\selectfont\color{popink},
  keywordstyle=\bfseries\color{popblue},
  commentstyle=\itshape\color{popmuted},
  stringstyle=\color{popgreen},
  numberstyle=\tiny\color{popmuted},
  identifierstyle=\color{popink},
  frame=l,
  rulecolor=\color{popink},
  framerule=1.6pt,
  framesep=7pt,
  framexleftmargin=7pt,
  framexrightmargin=7pt,
  xleftmargin=2pt,
  xrightmargin=2pt,
  breaklines=true,
  breakatwhitespace=false,
  showstringspaces=false,
  showspaces=false,
  showtabs=false,
  tabsize=2,
  columns=flexible,
  keepspaces=true,
  numbers=none,
  aboveskip=10pt,
  belowskip=10pt,
  captionpos=b,
}
\lstset{style=poplst, language=PHP}
% La distribution TeX embarquée ne fournit pas le fichier de langue
% JavaScript de listings (lstlang*.dat) : on la redéfinit nous-mêmes
% pour que \begin{lstlisting}[language=JavaScript] fonctionne.
\lstdefinelanguage{JavaScript}{
  keywords={break,case,catch,class,const,continue,debugger,default,delete,do,
    else,export,extends,finally,for,function,if,import,in,instanceof,let,new,
    return,static,super,switch,this,throw,try,typeof,var,void,while,with,yield,
    async,await,of,get,set},
  keywordstyle=\bfseries\color{popblue},
  ndkeywords={true,false,null,undefined,NaN,Infinity},
  ndkeywordstyle=\bfseries\color{poporange},
  identifierstyle=\color{popink},
  sensitive=true,
  comment=[l]{//},
  morecomment=[s]{/*}{*/},
  string=[b]",
  morestring=[b]',
  morestring=[b]`,
}
% Même limitation pour CSS et les fichiers de configuration (apacheconf,
% nginx, ini...) : ils ne sont pas fournis. CSS et un style générique de
% type "config" (commentaires # ou ;, pas de mots-clés) couvrent les besoins.
\lstdefinelanguage{CSS}{
  keywords={color,background,background-color,margin,padding,border,
    border-radius,font,font-size,font-weight,font-family,width,height,
    display,position,top,left,right,bottom,float,clear,text-align,
    line-height,list-style,cursor,overflow,box-shadow,transition,
    flex,grid,align-items,justify-content,z-index},
  keywordstyle=\bfseries\color{popblue},
  identifierstyle=\color{popink},
  sensitive=false,
  comment=[s]{/*}{*/},
  string=[b]",
  morestring=[b]',
}
\lstdefinelanguage{apacheconf}{
  sensitive=false,
  morecomment=[l]{\#},
}
\lstdefinelanguage{Apache}{
  sensitive=false,
  morecomment=[l]{\#},
}
\lstdefinelanguage{text}{sensitive=false}
\lstdefinelanguage{powershell}{
  keywords={New-Object,New-SmbShare,Get-Acl,Set-Acl,Get-ChildItem,Set-Content,
    Get-Content,Write-Host,ForEach-Object,Where-Object,Select-Object,
    Test-Path,Remove-Item,Copy-Item,Move-Item,Start-Process,Stop-Process,
    If,Else,ElseIf,ForEach,While,Function,Return,Try,Catch},
  keywordstyle=\bfseries\color{popblue},
  identifierstyle=\color{popink},
  sensitive=false,
  comment=[l]{\#},
  string=[b]",
  morestring=[b]',
}
\lstdefinelanguage{ini}{
  sensitive=false,
  morecomment=[l]{;},
  morecomment=[l]{\#},
}

% ============================================================
% Signature OnBuch+
% ============================================================
\newcommand{\poplogoinline}[1]{{\poplogo\fontsize{#1}{#1}\selectfont\color{popink}OnBuch\color{poporange}+}}
\newcommand{\signatureonbuch}{%
  \begin{tcolorbox}[enhanced,colback=popink,colframe=popink,boxrule=2.2pt,arc=10pt,
      drop shadow=poporange,left=14pt,right=14pt,top=10pt,bottom=10pt,before skip=0pt,after skip=0pt]
    \tikz[baseline=-0.7ex]\node[circle,fill=popgreen,draw=popcream,line width=1.3pt,minimum size=22pt,inner sep=0]{\popcheck{white}};%
    \hspace{9pt}\begin{minipage}[c]{0.62\linewidth}
      {\popxbold\fontsize{12}{13}\selectfont\color{popcream}Signé\ \textbullet\ {\poplogo OnBuch\color{poporange}+}}\par\vspace{2pt}
      {\raggedright\popsemi\fontsize{8}{10.5}\selectfont\color{popline}Ludovic A.\\\MakeUppercase{Conforme au programme officiel MINESEC}\par}
    \end{minipage}\hfill
    {\poplogo\fontsize{22}{22}\selectfont\color{popcream}OnBuch\color{poporange}+}
  \end{tcolorbox}%
}

% ============================================================
% Page de garde
% #1 titre · #2 matière · #3 niveau · #4 module · #5 n° leçon · #6 durée ·
% #7 description / objectifs (texte ou itemize)
% ============================================================
\newcommand{\pagegarde}[7]{%
  \thispagestyle{empty}
  \newgeometry{a4paper,margin=1.7cm,top=1.6cm,bottom=1.4cm}%
  \begin{tikzpicture}[remember picture,overlay]
    \fill[poporange!18] ([xshift=-3.2cm,yshift=-3.6cm]current page.north east) circle (4.6cm);
    \fill[poppurple!14] ([xshift=2.4cm,yshift=7.6cm]current page.south west) circle (3.8cm);
  \end{tikzpicture}%
  {\poplogo\fontsize{30}{30}\selectfont\color{popink}OnBuch\color{poporange}+}\hfill
  \raisebox{0.6ex}{\poppill{Cours signé \textbullet\ Premium}{popink}{popcream}}\par
  \vspace{1pt}\popsquiggle{poppurple}\par
  \vspace{8pt}
  \begin{tcolorbox}[enhanced,colback=poporange,colframe=popink,boxrule=2.8pt,arc=13pt,
      drop shadow=popink,left=18pt,right=18pt,top=12pt,bottom=13pt,after skip=10pt]
    \poppill{#2 \textbullet\ #3}{popink}{popcream}\hspace{5pt}\poppill{#4}{white}{popink}\par\vspace{8pt}
    {\popdisplay\fontsize{14}{14}\selectfont\color{popink}LEÇON #5}\par\vspace{4pt}
    {\raggedright\hyphenpenalty=10000\exhyphenpenalty=10000\popdisplay\fontsize{25}{27}\selectfont\color{popcream}\MakeUppercase{#1}\par}
    \vspace{8pt}
    {\popxbold\fontsize{9.5}{9.5}\selectfont\color{popink}\MakeUppercase{Durée conseillée : #6}}
  \end{tcolorbox}
  \begin{tcolorbox}[enhanced,colback=white,colframe=popink,boxrule=2.2pt,arc=10pt,
      drop shadow=popink,left=16pt,right=16pt,top=10pt,bottom=10pt,after skip=8pt]
    \poppill{Dans cette leçon}{poppurple}{white}\par\vspace{5pt}
    \popsemi\fontsize{9.3}{12.6}\selectfont\color{popink}#7
  \end{tcolorbox}
  \noindent\begin{minipage}[t]{0.487\linewidth}
    \begin{tcolorbox}[enhanced,colback=popgreen,colframe=popink,boxrule=2.2pt,arc=10pt,
        drop shadow=popink,left=11pt,right=11pt,top=10pt,bottom=10pt,height=3.1cm,valign=center]
      \tcbox[colback=white,colframe=popink,boxrule=1.3pt,arc=5pt,boxsep=0pt,nobeforeafter,
        left=3pt,right=3pt,top=3pt,bottom=3pt,valign=center]{\qrcode[height=1.9cm]{https://play.google.com/store/apps/details?id=com.luvvix.onbuch&hl=fr}}%
      \hspace{8pt}\begin{minipage}[c]{0.5\linewidth}
        {\popdisplay\fontsize{11}{12.5}\selectfont\color{white}\MakeUppercase{Télécharge OnBuch}}\par\vspace{4pt}
        {\raggedright\popsemi\fontsize{8.4}{11}\selectfont\color{white}Gratuit \textbullet\ Google Play\\Tuteur IA, annales, concours, résultats\par}
      \end{minipage}
    \end{tcolorbox}
  \end{minipage}\hfill
  \begin{minipage}[t]{0.487\linewidth}
    \begin{tcolorbox}[enhanced,colback=poppurple,colframe=popink,boxrule=2.2pt,arc=10pt,
        drop shadow=popink,left=11pt,right=11pt,top=10pt,bottom=10pt,height=3.1cm,valign=center]
      \tikz[baseline=-0.7ex]\node[circle,fill=white,draw=popink,line width=1.3pt,minimum size=1.6cm,inner sep=0]{\popdisplay\fontsize{24}{24}\selectfont\color{poppurple}+};%
      \hspace{8pt}\begin{minipage}[c]{0.6\linewidth}
        {\popdisplay\fontsize{11}{12.5}\selectfont\color{white}\MakeUppercase{Passe à OnBuch+}}\par\vspace{4pt}
        {\raggedright\popsemi\fontsize{8.4}{11}\selectfont\color{white}Tous les cours signés, Tuteur IA illimité, fiches PDF — onglet OnBuch+ de l'app\par}
      \end{minipage}
    \end{tcolorbox}
  \end{minipage}
  \vspace{8pt}
  \popcontenu
  \vfill
  \signatureonbuch
  \restoregeometry
  \newpage
}

% Rangée « Ce cours contient » (page de garde)
\newcommand{\poptile}[3]{% #1 couleur #2 gros texte #3 légende
  \begin{tcolorbox}[enhanced,colback=#1,colframe=popink,boxrule=2pt,arc=9pt,shadow={2.5pt}{-2.5pt}{0pt}{popink},
      left=6pt,right=6pt,top=9pt,bottom=9pt,halign=center,valign=center,height=1.75cm,width=\linewidth,nobeforeafter]
    {\popdisplay\fontsize{12.5}{13}\selectfont\color{white}\MakeUppercase{#2}}\par\vspace{3pt}
    {\centering\popsemi\fontsize{7.8}{9.5}\selectfont\color{white}#3\par}
  \end{tcolorbox}}
\newcommand{\popcontenu}{%
  \par\noindent{\popxboldls\fontsize{7.5}{7.5}\selectfont\color{popmuted}CE COURS SIGNÉ CONTIENT}\par\vspace{7pt}
  \noindent\begin{minipage}[t]{0.235\linewidth}\poptile{popblue}{Cours}{complet, illustré, pas à pas}\end{minipage}\hfill
  \begin{minipage}[t]{0.235\linewidth}\poptile{poporange}{Méthodes}{et exercices résolus}\end{minipage}\hfill
  \begin{minipage}[t]{0.235\linewidth}\poptile{poppink}{Exercices}{type examen + corrigés}\end{minipage}\hfill
  \begin{minipage}[t]{0.235\linewidth}\poptile{popgreen}{Fiche bilan}{l'essentiel à retenir}\end{minipage}\par}

% Sommaire
\newtcolorbox{sommaire}{enhanced,colback=white,colframe=popink,boxrule=2.2pt,arc=10pt,
  drop shadow=popink,left=18pt,right=18pt,top=13pt,bottom=13pt,before skip=6pt,after skip=20pt,
  before upper={\poppill{Sommaire}{poppurple}{white}\par\vspace{9pt}\popsemi\fontsize{10.6}{14.5}\selectfont\color{popink}}}

% Signature de fin de cours — poussée tout en bas de la dernière page
\newcommand{\findecours}{%
  \par\vfill\nopagebreak
  \begin{center}\popsquiggle{poporange}\end{center}\vspace{4pt}
  \signatureonbuch
}
\AtBeginDocument{\def\llbracket{\mathopen{[\mkern-3.5mu[}}\def\rrbracket{\mathclose{]\mkern-3.5mu]}}} % intervalles d'entiers (glyphe ⟦ absent de la police)
% Glyphes absents de Fira Math : redéfinis à partir de glyphes disponibles
\AtBeginDocument{%
  \def\setminus{\mathbin{\backslash}}\let\smallsetminus\setminus
  \def\vdots{\mathord{\vbox{\baselineskip3.2pt\lineskiplimit0pt\kern4pt\hbox{.}\hbox{.}\hbox{.}}}}%
  \def\ddots{\mathinner{\mkern1mu\raise6.5pt\hbox{.}\mkern2mu\raise3.6pt\hbox{.}\mkern2mu\raise0.7pt\hbox{.}\mkern1mu}}%
  \def\triangle{\mathord{\scalebox{0.9}{$\symup{\Delta}$}}}%
  \def\nabla{\mathord{\raisebox{\depth}{\scalebox{1}[-1]{$\symup{\Delta}$}}}}%
  \def\checkmark{\popcheck{popdark}}%
  \def\star{\mathbin{\ast}}\def\bigstar{\mathord{\ast}}%
  \def\diamond{\mathbin{\scalebox{0.6}{$\Diamond$}}}%
  \def\bigcup{\mathop{\mathchoice{\raisebox{-0.3em}{\scalebox{1.7}{$\cup$}}}{\scalebox{1.25}{$\cup$}}{\cup}{\cup}}\displaylimits}%
  \def\bigcap{\mathop{\mathchoice{\raisebox{-0.3em}{\scalebox{1.7}{$\cap$}}}{\scalebox{1.25}{$\cap$}}{\cap}{\cap}}\displaylimits}%
  \def\lesseqqgtr{\mathrel{\lessgtr}}\def\bigoplus{\mathop{\oplus}\displaylimits}%
}
\providecommand{\ssi}{si et seulement si} % abréviation fréquente
\AtBeginDocument{\providecommand{\wideparen}{\overparen}} % arc de cercle

%% FIGFIX-BEGIN v5 — correctifs globaux des figures (docs/onbuch-generation/tools/figfix)
\usepackage{adjustbox}\usepackage{etoolbox}\tcbuselibrary{raster}
% toute figure TikZ/pgfplots plus large que la ligne est réduite à la largeur disponible
\BeforeBeginEnvironment{tikzpicture}{\begin{adjustbox}{max width=\linewidth}}
\AfterEndEnvironment{tikzpicture}{\end{adjustbox}}
% légendes pgfplots lisibles par-dessus les courbes
\pgfplotsset{every axis/.append style={legend style={fill=white,fill opacity=0.92,text opacity=1,draw=black!15,inner sep=3pt}}}
% étiquettes d'arêtes (syntaxe edge["..."]) avec fond blanc
\tikzset{every edge quotes/.append style={fill=white,inner sep=1.2pt}}
%% FIGFIX-END
\begin{document}

\pagegarde{Utilisation d’un logiciel de traitement de texte}{Informatique}{3ème}{Informatique – classe de 3ème}{N05}{3 séances (fiche de progression harmonisée)}{Cette leçon apprend à exploiter pleinement un logiciel de traitement de texte : mise en forme des caractères et des paragraphes, mise en page du document, insertion d'objets (images, tableaux, formes) et génération automatique d'une table des matières. L'élève produit à la fin un document complet et professionnel, compétence clé pour le BEPC et la vie courante.
\vspace{4pt}
\begin{multicols}{2}\small
\begin{itemize}
\item À la fin de cette leçon, je sais appliquer une mise en forme aux caractères (police, taille, style, couleur) d'un texte.
\item À la fin de cette leçon, je sais appliquer une mise en forme aux paragraphes (alignement, retrait, interligne, puces et numérotation).
\item À la fin de cette leçon, je sais régler la mise en page d'un document (orientation, marges, taille du papier, en-tête et pied de page, numérotation des pages).
\item À la fin de cette leçon, je sais insérer et habiller des objets dans un document : images, tableaux, formes, sauts de page et de section.
\item À la fin de cette leçon, je sais appliquer des styles de titres pour structurer un document.
\item À la fin de cette leçon, je sais générer et mettre à jour automatiquement une table des matières.
\end{itemize}\end{multicols}}

\begin{sommaire}
\begin{multicols}{2}
\begin{enumerate}
\item Rappels : interface et manipulation de base
\item Mise en forme des caractères et des paragraphes
\item Mise en page du document
\item Insertion d'objets dans le document
\item Styles de titres et table des matières automatique
\item Synthèse : produire un document complet
\item Activité d'intégration
\item Méthodes et exercices résolus
\item Exercices
\item Corrigés des exercices
\item Fiche bilan
\end{enumerate}
\end{multicols}
\end{sommaire}

\begin{objectifs}
\begin{itemize}
\item À la fin de cette leçon, je sais appliquer une mise en forme aux caractères (police, taille, style, couleur) d'un texte.
\item À la fin de cette leçon, je sais appliquer une mise en forme aux paragraphes (alignement, retrait, interligne, puces et numérotation).
\item À la fin de cette leçon, je sais régler la mise en page d'un document (orientation, marges, taille du papier, en-tête et pied de page, numérotation des pages).
\item À la fin de cette leçon, je sais insérer et habiller des objets dans un document : images, tableaux, formes, sauts de page et de section.
\item À la fin de cette leçon, je sais appliquer des styles de titres pour structurer un document.
\item À la fin de cette leçon, je sais générer et mettre à jour automatiquement une table des matières.
\item À la fin de cette leçon, je sais produire un document complet répondant à une situation concrète de la vie courante (lettre, rapport, affiche).
\item À la fin de cette leçon, je sais rédiger et justifier la démarche suivie pour réaliser un document demandé.
\end{itemize}
\end{objectifs}

\begin{prerequis}
\begin{itemize}
\item Notion de fichier, dossier et enregistrement d'un document (classe de 4ème).
\item Lancement d'un logiciel et découverte de l'interface d'un traitement de texte (ruban, barre d'outils, zone de saisie).
\item Saisie et correction simple d'un texte au clavier (copier, couper, coller).
\item Notion de périphériques d'entrée et de sortie (clavier, souris, imprimante).
\end{itemize}
\end{prerequis}

\newpage

\coursec{Rappels : interface et manipulation de base}

Imagine un instant la scène : nous sommes au collège de Bafoussam, en plein cœur de l'année scolaire. Le principal appelle son secrétaire dans son bureau et lui confie une mission de taille : « \emph{Prépare le rapport annuel de l'établissement. Il doit faire une vingtaine de pages, avec une page de garde soignée, des photos de la fête du sport, un tableau détaillé des effectifs par classe et, bien sûr, une table des matières. Je te le veux pour demain matin ! } »

Le secrétaire, courageux mais pressé, ouvre son logiciel et se met au travail. Il tape le texte \og au kilomètre \fg{}, appuie sur \texttt{Entrée} dix fois pour faire de la place, espace manuellement son titre au milieu de la page avec la barre d'espace, insère ses photos en les faisant glisser n'importe comment, dessine son tableau à coups de tirets \texttt{-----} et de barres \texttt{|}, et écrit sa table des matières \emph{à la main}, page par page. Résultat ? Le document est bancal. Pire : le principal revient une heure plus tard : \og \emph{Change le titre, ajoute une classe en 3ème, et met à jour les numéros de page de la table des matières ! } \fg{} Le secrétaire pâlit : tout est à refaire. La mise en page explose, les images se chevauchent, les numéros de page ne correspondent plus à rien.

\medskip
\noindent \textbf{Problématique :} Comment un logiciel de traitement de texte moderne (Microsoft Word, LibreOffice Writer, Google Docs) permet-il de structurer un document long, d'automatiser les tâches répétitives (numérotation, table des matières) et de modifier le contenu sans casser la présentation, le tout en un temps record ?

\medskip
\noindent Pour répondre à cette question, nous devons d'abord maîtriser l'outil : son interface, ses commandes de base et sa logique de gestion des fichiers. C'est l'objet de cette première section.

\courssub{Qu'est-ce qu'un logiciel de traitement de texte ?}

\begin{definition}[Logiciel de traitement de texte]
Un \cle{logiciel de traitement de texte} est une application informatique conçue pour la \textbf{création}, la \textbf{modification}, la \textbf{mise en forme} et l'\textbf{impression} de documents textuels. Contrairement à un simple éditeur de texte (comme le Bloc-notes de Windows), il gère la \cle{mise en page} (marges, colonnes, en-têtes/pieds de page), la \cle{mise en forme} (polices, couleurs, styles), l'insertion d'objets multimédias (images, tableaux, graphiques) et des fonctions avancées (publipostage, table des matières automatique, suivi des modifications).
\end{definition}

\begin{savaistu}[Le saviez-vous ? : LibreOffice Writer, la puissance du libre]
\cle{LibreOffice Writer} est l'alternative libre, gratuite et open-source à Microsoft Word. Développé par \emph{The Document Foundation}, il est installé par défaut sur de nombreuses distributions Linux (Ubuntu, Debian, Fedora) et disponible pour Windows et macOS.
\begin{itemize}
    \item \textbf{Aucune licence payante :} vous pouvez l'installer sur 1, 10 ou 1000 ordinateurs au collège de Bafoussam sans débourser un franc CFA.
    \item \textbf{Format natif ouvert (ODF) :} il enregistre par défaut en \texttt{.odt} (OpenDocument Text), un standard international (ISO/IEC 26300) qui garantit la pérennité de vos documents, même dans 20 ans.
    \item \textbf{Compatibilité :} il ouvre et enregistre parfaitement les fichiers \texttt{.docx} de Word (avec quelques différences mineures sur les mises en page très complexes).
\end{itemize}
Dans le cadre du programme MINESEC, la maîtrise de Writer est un atout majeur : elle assure l'autonomie numérique des établissements camerounais face aux coûts des licences propriétaires.
\end{savaistu}

Les trois références du marché que vous rencontrerez le plus souvent sont :
\begin{itemize}
    \item \textbf{Microsoft Word} (payant, suite Microsoft 365 / Office) : standard de facto en entreprise et administration.
    \item \textbf{LibreOffice Writer} (gratuit, open-source) : standard dans l'éducation publique et les administrations soucieuses de souveraineté numérique.
    \item \textbf{Google Docs} (gratuit, en ligne, collaboratif en temps réel) : idéal pour le travail d'équipe à distance, nécessite une connexion Internet.
\end{itemize}

\courssub{Découverte de l'interface : le poste de pilotage}

Que vous soyez sur Word ou Writer, l'interface suit le même paradigme : le \textbf{Ruban} (introduit par Office 2007). Comprendre son organisation, c'est gagner des heures de recherche dans les menus.

\begin{popfigure}
\begin{tikzpicture}[
    node distance=0.8cm and 1cm,
    box/.style={rectangle, draw=popdark, thick, rounded corners=4pt, fill=popcream, text width=3.2cm, align=center, minimum height=1.8cm, font=\small},
    arrow/.style={-Stealth, thick, popblue},
    labelbox/.style={rectangle, draw=poporange, thick, rounded corners=3pt, fill=poporangeL, font=\tiny, align=center, inner sep=2pt}
]
% Fenêtre principale
\node[rectangle, draw=popdark, very thick, rounded corners=6pt, minimum width=14cm, minimum height=9cm, fill=white] (win) {};
% Barre de titre
\node[rectangle, fill=popblue, draw=popblue, minimum width=13.6cm, minimum height=0.7cm, rounded corners=4pt, anchor=north] (titlebar) at (win.north) {};
\node[white, font=\bfseries\small] at (titlebar) {\texttt{Rapport\_Annuel\_Collège\_Bafoussam} - LibreOffice Writer};
% Ruban - Onglets
\node[rectangle, fill=popmuted, draw=popline, minimum width=13.6cm, minimum height=0.6cm, anchor=north] (ribbon_bg) at ($(titlebar.south) + (0,-0.1)$) {};
\node[labelbox, anchor=west] (tab1) at ($(ribbon_bg.north west) + (0.3,-0.1)$) {Fichier};
\node[labelbox, anchor=west, fill=popblueL, draw=popblue] (tab2) at ($(tab1.east) + (0.2,0)$) {\textbf{Accueil}};
\node[labelbox, anchor=west] (tab3) at ($(tab2.east) + (0.2,0)$) {Insertion};
\node[labelbox, anchor=west] (tab4) at ($(tab3.east) + (0.2,0)$) {Mise en page};
\node[labelbox, anchor=west] (tab5) at ($(tab4.east) + (0.2,0)$) {Références};
% Groupes du ruban (onglet Accueil actif)
\node[rectangle, draw=popline, fill=white, minimum width=2.5cm, minimum height=1.8cm, anchor=north west] (grp1) at ($(ribbon_bg.south west) + (0.3,-0.15)$) {};
\node[font=\tiny\bfseries, text=popdark] at (grp1.north) {Presse-papiers};
\node[font=\tiny, text=popdark] at (grp1.center) {Couper \textbar{} Copier \textbar{} Coller \textbar{} Format};
\node[rectangle, draw=popline, fill=white, minimum width=3.5cm, minimum height=1.8cm, anchor=north west] (grp2) at ($(grp1.north east) + (0.2,0)$) {};
\node[font=\tiny\bfseries, text=popdark] at (grp2.north) {Police};
\node[font=\tiny, text=popdark] at (grp2.center) {Calibri 11 \textbar{} \textbf{G} \textit{I} \underline{S} \textbar{} Couleur};
\node[rectangle, draw=popline, fill=white, minimum width=3.5cm, minimum height=1.8cm, anchor=north west] (grp3) at ($(grp2.north east) + (0.2,0)$) {};
\node[font=\tiny\bfseries, text=popdark] at (grp3.north) {Paragraphe};
\node[font=\tiny, text=popdark] at (grp3.center) {Alignement \textbar{} Puces \textbar{} Numérotation \textbar{} Interligne};
\node[rectangle, draw=popline, fill=white, minimum width=3.5cm, minimum height=1.8cm, anchor=north west] (grp4) at ($(grp3.north east) + (0.2,0)$) {};
\node[font=\tiny\bfseries, text=popdark] at (grp4.north) {Styles};
\node[font=\tiny, text=popdark] at (grp4.center) {Normal \textbar{} Titre 1 \textbar{} Titre 2 \textbar{} Titre 3};
% Règles
\node[rectangle, draw=popline, fill=popmuted!18, minimum width=13.6cm, minimum height=0.4cm, anchor=north] (hruler) at ($(grp1.south) + (0,-0.3)$) {};
\node[font=\tiny, rotate=90, anchor=south] at (hruler.west) {Règle verticale};
\node[font=\tiny, anchor=south] at (hruler.north) {Règle horizontale (marges, tabulations, indentations)};
% Zone de saisie
\node[rectangle, draw=popline, fill=white, minimum width=13.6cm, minimum height=4.5cm, anchor=north] (edit) at ($(hruler.south) + (0,-0.1)$) {};
\node[text=popmuted, font=\itshape\small] at (edit.center) {Zone de saisie (page blanche) : c'est ici que vous tapez votre texte.\newline Curseur clignotant = position d'insertion.};
% Barre d'état
\node[rectangle, fill=popmuted, draw=popline, minimum width=13.6cm, minimum height=0.5cm, anchor=south] (status) at (win.south) {};
\node[font=\small] at (status.west) {Page 1/1 \textbar{} Mots: 0 \textbar{} Langue: Français (Cameroun)};
\node[font=\small] at (status.east) {100\% \textbar{} [Zoom - +] \textbar{} Vue: Page Web / Plan / Brouillon};
% Flèches d'annotation
\draw[arrow] ($(tab2.north) + (0,0.3)$) -- ++(0,0.5) node[above, font=\small\bfseries, text=popblue] {Onglets du Ruban};
\draw[arrow] ($(grp2.east) + (0.5,0)$) -- ++(1.5,0) node[right, font=\small, text=popdark, align=left] {Groupes de commandes\\contextuels};
\draw[arrow] ($(hruler.north) + (0,0.2)$) -- ++(0,0.6) node[above, font=\small\bfseries, text=popdark] {Règles (repères visuels)};
\draw[arrow] ($(edit.west) + (-1.5, 1)$) -- ++(-1,0) node[left, font=\small\bfseries, text=popdark, align=right] {Zone de saisie\\(surface de travail)};
\draw[arrow] ($(status.west) + (-1,0)$) -- ++(-1.5,0) node[left, font=\small\bfseries, text=popdark, align=right] {Barre d'état\\infos contextuelles};
\end{tikzpicture}
\legende{Interface type de LibreOffice Writer / Microsoft Word : identification des zones fonctionnelles.}
\end{popfigure}

\noindent Décortiquons les cinq zones essentielles annotées sur le schéma :

\begin{enumerate}
    \item \textbf{La barre de titre} (haut de fenêtre) : affiche le nom du fichier (\texttt{Rapport\_Annuel\_Collège\_Bafoussam.odt}) et le nom de l'application. Si le fichier n'est pas encore enregistré, on voit \og Sans titre 1 \fg{}.
    \item \textbf{Le Ruban et ses Onglets} : c'est la \textbf{boîte à outils} principale. Chaque onglet regroupe des commandes par \emph{thème} :
        \begin{itemize}
            \item \cle{Accueil} : les actions les plus fréquentes (presse-papiers, police, paragraphe, \textbf{styles}).
            \item \cle{Insertion} : tout ce qui n'est pas du texte brut (images, tableaux, graphiques, en-têtes, sauts de page).
            \item \cle{Mise en page} : structure globale du document (marges, orientation, colonnes, filigrane).
            \item \cle{Références} : automatismes pour documents longs (table des matières, notes de bas de page, citations, index).
        \end{itemize}
    \item \textbf{Les Règles} (horizontale et verticale) : repères visuels pour régler les \cle{marges}, les \cle{indentations} (retraits de paragraphe) et les \cle{tabulations}. Un double-clic sur la règle ouvre la boîte de dialogue \og Paragraphe \fg{} ou \og Tabulations \fg{}.
    \item \textbf{La Zone de saisie} : la page virtuelle. Le \cle{curseur d'insertion} (trait vertical clignotant) indique où le prochain caractère apparaîtra. La sélection apparaît en surbrillance (généralement bleu).
    \item \textbf{La Barre d'état} (bas de fenêtre) : informations en temps réel : numéro de page / total, nombre de mots/caractères, langue du correcteur orthographique, niveau de zoom, mode d'affichage.
\end{enumerate}

\begin{attention}[Piège classique : les onglets contextuels]
Les onglets \textbf{ne sont pas tous visibles en permanence} ! Si vous cliquez sur une \emph{image}, un onglet \og Format d'image \fg{} (Word) ou \og Image \fg{} (Writer) apparaît \emph{au-dessus} du ruban, souvent coloré différemment. Si vous cliquez dans un \emph{tableau}, les onglets \og Création \fg{} et \og Disposition \fg{} apparaissent. \textbf{Règle d'or :} sélectionnez l'objet \emph{avant} de chercher sa commande.
\end{attention}

\courssub{Gestion des fichiers : créer, ouvrir, enregistrer, exporter}

Un document n'existe vraiment que lorsqu'il est \textbf{enregistré sur un support de stockage} (disque dur, clé USB, cloud). La mémoire vive (RAM) est volatile : une coupure de courant (fréquente au Cameroun !) efface tout ce qui n'est pas sauvegardé.

\begin{methode}[Enregistrer un document sous un nom et un emplacement précis (Enregistrer sous)]
Cette méthode s'applique \textbf{la première fois} que vous sauvegardez un fichier, ou lorsque vous voulez créer une \textbf{copie différente} (autre nom, autre dossier, autre format).
\begin{enumerate}
    \item Cliquez sur l'onglet \textbf{Fichier} (coin supérieur gauche) → \textbf{Enregistrer sous} (ou raccourci \texttt{Ctrl + Maj + S}).
    \item Dans la boîte de dialogue, choisissez l'\textbf{emplacement} (dossier) : \texttt{Documents}, \texttt{Bureau}, ou votre \texttt{Clé USB (E:)}. \emph{Astuce : créez un dossier \texttt{Informatique\_3eme} pour ranger vos TP.}
    \item Dans le champ \textbf{Nom du fichier}, tapez un nom explicite \textbf{sans caractères spéciaux} : \texttt{Rapport\_Annuel\_Collège\_Bafoussam\_v1}. Évitez : \texttt{Rapport final (1).docx}.
    \item Vérifiez le \textbf{Type de fichier} (Format) :
        \begin{itemize}
            \item \texttt{ODF Text Document (.odt)} : format natif LibreOffice, recommandé pour l'archivage.
            \item \texttt{Microsoft Word (.docx)} : format standard pour l'échange avec des utilisateurs Word.
            \item \texttt{PDF (.pdf)} : format \emph{figé} pour impression ou diffusion (non modifiable facilement).
        \end{itemize}
    \item Cliquez sur \textbf{Enregistrer}.
\end{enumerate}
\end{methode}

\begin{propriete}[Raccourci vital : Enregistrer régulièrement (Ctrl + S)]
La commande \textbf{Enregistrer} (icône disquette ou \texttt{Ctrl + S}) écrase la version précédente sur le disque \textbf{sans redemander le nom ni l'emplacement}.
\begin{itemize}
    \item \textbf{Fréquence :} toutes les 2 à 3 minutes, ou après chaque paragraphe important.
    \item \textbf{Indicateur visuel :} l'icône de disquette est souvent grisée si le fichier est déjà à jour ; elle redevient active (colorée) dès qu'une modification est faite.
    \item \textbf{Sauvegarde automatique :} Word/Writer proposent une récupération automatique (toutes les 10 min par défaut), mais \textbf{ne comptez pas dessus} : \texttt{Ctrl + S} doit devenir un réflexe musculaire.
\end{itemize}
\end{propriete}

\begin{attention}[Erreur fréquente : Confondre \og Enregistrer \fg{} et \og Enregistrer sous \fg{}]
\begin{itemize}
    \item \textbf{Scénario catastrophe :} Vous ouvrez \texttt{Modèle\_Rapport.odt}, vous le modifiez entièrement pour le collège de Bafoussam, vous faites \texttt{Fichier → Enregistrer}. \textbf{Résultat :} Votre modèle vierge est écrasé ! Il contient maintenant les données de Bafoussam. Vous avez perdu le modèle.
    \item \textbf{Bonne pratique :} Ouvrez le modèle → immédiatement \texttt{Fichier → Enregistrer sous} → nommez \texttt{Rapport\_Bafoussam\_2026.odt}. Travaillez sur cette copie. Le modèle original reste intact.
    \item \textbf{Mnémotechnique :} \og \textbf{Sous} = \textbf{S}ecours (nouveau fichier) / \og \textbf{Enregistrer} = \textbf{M}ise à jour (même fichier).
\end{itemize}
\end{attention}

\begin{exoresolu}[Gestion de versions pour le rapport du principal]
\textbf{Énoncé :} Le secrétaire doit produire le rapport annuel. Il part d'un modèle \texttt{Modele\_Rapport.odt}. Il doit produire une version pour le principal (\texttt{v1}), puis une version corrigée après relecture (\texttt{v2}), et enfin une version PDF pour le site web de l'école.
\textbf{Solution détaillée :}
\begin{enumerate}
    \item \texttt{Fichier → Ouvrir → Modele\_Rapport.odt}.
    \item \texttt{Fichier → Enregistrer sous → Dossier: Rapports\_2026 → Nom: Rapport\_Annuel\_v1.odt → Enregistrer}.
    \item Travail sur \texttt{v1} (saisie, mise en forme). \texttt{Ctrl+S} régulier.
    \item Relecture du principal → modifications demandées.
    \item \texttt{Fichier → Enregistrer sous → Nom: Rapport\_Annuel\_v2.odt}. (On garde \texttt{v1} comme archive/historique).
    \item Modifications sur \texttt{v2}. \texttt{Ctrl+S}.
    \item \texttt{Fichier → Exporter vers → Exporter au format PDF → Nom: Rapport\_Annuel\_Final.pdf}. Ce PDF est figé, prêt pour le site web.
\end{enumerate}
\textbf{Arborescence finale sur la clé USB :}
\begin{lstlisting}[language=bash]
/Rapports_2026/
  |-- Modele_Rapport.odt          (Original intact)
  |-- Rapport_Annuel_v1.odt       (Version 1, archive)
  |-- Rapport_Annuel_v2.odt       (Version de travail finale)
  \-- Rapport_Annuel_Final.pdf    (Version diffusée, non modifiable)
\end{lstlisting}
\end{exoresolu}

\courssub{Manipulation du texte : sélection, presse-papiers, annuler/rétablir}

Avant de formater, il faut \textbf{désigner} ce que l'on veut modifier. La sélection est l'action fondamentale.

\begin{propriete}[Techniques de sélection rapides (souris et clavier)]
Maîtriser ces raccourcis change la vitesse de travail du tout au tout.
\begin{center}
\begin{tabularx}{\linewidth}{|l|X|X|}
\hline
\rowcolor{popblueL}
\textbf{Action} & \textbf{Souris} & \textbf{Clavier (curseur au départ)} \\
\hline
Un \textbf{mot} & Double-clic sur le mot & \texttt{Ctrl + Maj + Flèche Droite/Gauche} \\
Une \textbf{ligne} & Clic dans la marge gauche (zone de sélection) & \texttt{Fin} puis \texttt{Maj + Origine} (ou \texttt{Maj + Flèche Bas}) \\
Un \textbf{paragraphe} & Triple-clic dans le paragraphe & \texttt{Ctrl + Maj + Flèche Bas/Haut} \\
\textbf{Tout le document} & \texttt{Ctrl + Clic} dans la marge gauche & \textbf{\texttt{Ctrl + A}} (le plus important à retenir !) \\
\hline
\end{tabularx}
\end{center}
\emph{Note :} Maintenir \texttt{Maj} (Shift) enfoncé pendant qu'on déplace le curseur (flèches, \texttt{Origine}, \texttt{Fin}, \texttt{Ctrl+Fin}) étend la sélection. C'est la \textbf{règle universelle} de sélection.
\end{propriete}

\begin{definition}[Le Presse-papiers : mémoire tampon d'échange]
Le \cle{presse-papiers} est une zone de la mémoire vive (RAM) qui stocke \textbf{un seul élément} (texte, image, fichier) à la fois (sauf historique activé dans Windows 10/11 ou extensions). Trois commandes de base :
\begin{itemize}
    \item \textbf{Couper} (\texttt{Ctrl + X}) : \emph{Supprime} la sélection du document et la place dans le presse-papiers. \og Je déplace. \fg{}
    \item \textbf{Copier} (\texttt{Ctrl + C}) : \emph{Duplique} la sélection dans le presse-papiers sans la retirer du document. \og Je clone. \fg{}
    \item \textbf{Coller} (\texttt{Ctrl + V}) : \emph{Insère} le contenu du presse-papiers à la position du curseur. \og Je dépose. \fg{}
\end{itemize}
\end{definition}

\begin{experience}[TP Guidés : Manipulation de base sur un texte existant]
\textbf{Objectif :} Maîtriser l'ouverture, la sélection, le déplacement de texte et la sauvegarde.
\textbf{Matériel :} Poste de travail avec LibreOffice Writer (ou Word). Fichier \texttt{Texte\_Brute.odt} fourni (contenant 3 paragraphes mélangés sur l'histoire du Cameroun).
\textbf{Manipulation :}
\begin{enumerate}
    \item \textbf{Ouvrir :} \texttt{Fichier → Ouvrir} → sélectionner \texttt{Texte\_Brute.odt}.
    \item \textbf{Sauvegarder immédiatement :} \texttt{Fichier → Enregistrer sous} → Dossier \texttt{TP\_Traitement\_Texte} → Nom \texttt{Mon\_Rapport\_Cameroun.odt}.
    \item \textbf{Sélectionner le 2ème paragraphe :} Triple-clic dedans. Observer la surbrillance.
    \item \textbf{Déplacer ce paragraphe en premier :} \texttt{Ctrl + X} (Couper). Placer curseur au tout début (avant 1er mot). \texttt{Ctrl + V} (Coller).
    \item \textbf{Dupliquer le dernier paragraphe à la fin :} Sélectionner dernier paragraphe (triple-clic). \texttt{Ctrl + C}. Aller à la fin (\texttt{Ctrl + Fin}). \texttt{Entrée} (nouveau paragraphe). \texttt{Ctrl + V}.
    \item \textbf{Erreur volontaire :} Supprimer un mot important (\texttt{Suppr}). Observer : le mot a disparu.
    \item \textbf{Annuler :} \texttt{Ctrl + Z} (ou icône flèche courbe gauche). Le mot réapparaît.
    \item \textbf{Rétablir :} \texttt{Ctrl + Y} (ou flèche courbe droite). Le mot disparaît à nouveau.
    \item \textbf{Enregistrer :} \texttt{Ctrl + S}. Fermer.
\end{enumerate}
\textbf{Observations attendues :} Le document \texttt{Mon\_Rapport\_Cameroun.odt} a les paragraphes réorganisés. L'original \texttt{Texte\_Brute.odt} est inchangé. L'historique Annuler/Rétablir permet de naviguer dans le temps.
\textbf{Interprétation :} \texttt{Ctrl+X}/\texttt{Ctrl+V} = Restructuration rapide (déplacer des blocs). \texttt{Ctrl+C}/\texttt{Ctrl+V} = Réutilisation (cloner). \texttt{Ctrl+Z} = Filet de sécurité immédiat.
\end{experience}

\begin{aretenir}[Récapitulatif des raccourcis indispensables (Niveau 3ème - BEPC)]
Ces raccourcis sont \textbf{universels} (Word, Writer, Google Docs, Navigateur, Explorateur de fichiers...). Apprenez-les par cœur.
\begin{center}
\begin{tabularx}{\linewidth}{|l|X|}
\hline
\rowcolor{popblueL}
\textbf{Raccourci} & \textbf{Action} \\
\hline
\texttt{Ctrl + N} & Nouveau document \\
\texttt{Ctrl + O} & Ouvrir un document existant \\
\texttt{Ctrl + S} & \textbf{Enregistrer (Sauver)} — \emph{Faites-le tout le temps !} \\
\texttt{Ctrl + Maj + S} & Enregistrer sous (nouveau nom/emplacement) \\
\texttt{Ctrl + P} & Imprimer / Aperçu avant impression \\
\texttt{Ctrl + Z} & Annuler (Undo) \\
\texttt{Ctrl + Y} & Rétablir (Redo) \\
\texttt{Ctrl + A} & \textbf{Sélectionner Tout} (All) \\
\texttt{Ctrl + X} & Couper \\
\texttt{Ctrl + C} & Copier \\
\texttt{Ctrl + V} & Coller \\
\texttt{Ctrl + F} & Rechercher (Find) \\
\texttt{Ctrl + H} & Remplacer (Replace) \\
\hline
\end{tabularx}
\end{center}
\end{aretenir}

\courssub{Synthèse de la section}

Nous avons posé les fondations. Vous savez maintenant :
\begin{itemize}
    \item Identifier les \textbf{quatre zones majeures} de l'interface (Ruban, Règles, Zone de saisie, Barre d'état) et le rôle des \textbf{onglets} (Accueil, Insertion, Mise en page, Références).
    \item Différencier \textbf{Enregistrer} (mise à jour silencieuse) et \textbf{Enregistrer sous} (création d'une nouvelle occurrence), et choisir le bon format (\texttt{.odt} pour travailler, \texttt{.pdf} pour diffuser).
    \item \textbf{Sélectionner} efficacement (mot, ligne, paragraphe, tout) grâce à la souris et surtout au clavier (\texttt{Ctrl+A}, \texttt{Maj+Flèches}).
    \item Utiliser le \textbf{presse-papiers} (\texttt{Ctrl+X/C/V}) pour restructurer votre texte et le duo \textbf{Annuler/Rétablir} (\texttt{Ctrl+Z/Y}) comme filet de sécurité.
\end{itemize}

Dans la prochaine section, nous allons transformer ce texte brut en un document lisible et professionnel : \textbf{la mise en forme des caractères et des paragraphes} (polices, alignements, interlignes, puces, bordures). C'est là que le rapport du principal va commencer à ressembler à quelque chose !

\coursec{Mise en forme des caractères et des paragraphes}

Dans la section précédente, tu as appris à te repérer dans l'interface d'un logiciel de traitement de texte (Word, Writer) : le ruban, les onglets, la zone de saisie, la règle. Tu sais maintenant taper du texte, le sélectionner, le copier, le déplacer et enregistrer ton document. Mais un texte tapé « brut » ressemble à un brouillon : tout est écrit avec la même police, la même taille, sans aucune hiérarchie. Or, dans la vie courante, aucun document sérieux n'est présenté ainsi : pense à un bulletin de notes, à une affiche de la fête de la jeunesse ou à une lettre de demande de stage adressée au principal du collège. Chacun de ces documents utilise des \textbf{caractères} et des \textbf{paragraphes} soigneusement mis en forme pour être lisibles et agréables à lire. C'est précisément l'objet de cette section : apprendre à transformer un texte brut en un document clair et professionnel.

\begin{definition}[Mise en forme et mise en page]
La \cle{mise en forme} est l'ensemble des modifications d'aspect appliquées aux \textbf{caractères} (police, taille, gras, italique, couleur\ldots) et aux \textbf{paragraphes} (alignement, retraits, espacements, puces, bordures\ldots). La \cle{mise en page}, que nous étudierons à la section suivante, concerne le document entier : orientation de la feuille, marges, en-têtes et pieds de page, numérotation des pages.
\end{definition}

\courssub{La mise en forme des caractères}

\textbf{Intuition.} Un caractère, c'est une lettre, un chiffre ou un symbole. Quand tu écris « Interrogation de Mathématiques » sur une affiche, tu veux que ce titre soit grand, épais et bien visible ; la date, elle, peut être plus petite. La mise en forme des caractères permet exactement cela : donner à chaque portion de texte l'apparence qui convient à son rôle.

\begin{definition}[Attributs d'un caractère]
Les principaux attributs de mise en forme d'un caractère sont :
\begin{itemize}
\item la \cle{police} : le dessin des lettres (exemples : Arial, Times New Roman, Calibri) ;
\item la \cle{taille} : la hauteur des caractères, exprimée en \textbf{points} (exemples : 11, 12, 16) ;
\item le \cle{gras} (\textbf{B}) : lettres épaissies, pour attirer l'attention ;
\item l'\cle{italique} (\emph{I}) : lettres penchées, pour les titres d'œuvres ou les mots étrangers ;
\item le \cle{souligné} : un trait sous le texte ;
\item la \cle{couleur} de police et le \cle{surlignage} (comme un stabilo sur la copie) ;
\item l'\cle{exposant} (comme dans $m^2$) et l'\cle{indice} (comme dans $H_2O$).
\end{itemize}
\end{definition}

\begin{propriete}[Règle d'or : sélectionner d'abord]
Dans un traitement de texte, toute mise en forme s'applique à la \textbf{sélection}. Il faut donc toujours : (1) sélectionner le texte à modifier, (2) cliquer sur le bouton ou choisir le réglage voulu dans l'onglet \texttt{Accueil}. Si rien n'est sélectionné, la mise en forme s'applique au texte que tu vas taper ensuite.
\end{propriete}

\begin{exemplebox}[Un exemple chiffré complet]
Pour la page de garde d'un exposé, on choisit par exemple :
\begin{itemize}
\item le titre « LES VOLCANS DU CAMEROUN » en police \textbf{Arial}, taille \textbf{16}, \textbf{gras}, \textbf{centré} ;
\item le corps du texte en \textbf{Times New Roman}, taille \textbf{12}, \textbf{justifié}, avec un \textbf{interligne de 1,5} ;
\item la formule de l'eau écrite $H_2O$ grâce au bouton \textbf{indice}, et une aire de 25~m$^2$ grâce au bouton \textbf{exposant}.
\end{itemize}
\end{exemplebox}

\begin{attention}[Piège : le gras partout ne sert à rien]
Une erreur fréquente consiste à mettre en gras, en couleur et en souligné de très longs passages. Résultat : plus rien ne ressort et le document devient illisible. La mise en forme doit \emph{guider l'œil} : réserve le gras aux mots vraiment importants et limite-toi à deux ou trois couleurs dans tout le document.
\end{attention}

\begin{propriete}[Espace insécable : protéger la ponctuation]
En français, on ne sépare pas un mot de la ponctuation double qui le suit (\texttt{; : ? ! « »}). Utilise une \cle{espace insécable} (raccourci \texttt{Ctrl + Maj + Espace} sur Word / \texttt{Ctrl + Espace} sur Writer) pour éviter qu'un deux-points ou un point d'interrogation ne se retrouve seul en début de ligne.
\end{propriete}

\courssub{L'alignement des paragraphes}

\textbf{Intuition.} Un paragraphe est un bloc de texte terminé par la touche \texttt{Entrée}. Selon le document, on ne place pas le texte de la même façon sur la ligne : un titre est souvent centré, une date en haut d'une lettre est placée à droite, et le texte d'un livre occupe toute la largeur de la page. C'est le rôle de l'\textbf{alignement}.

\begin{definition}[Les quatre alignements]
Il existe quatre alignements de paragraphe :
\begin{itemize}
\item \cle{aligné à gauche} : le texte part du bord gauche, le bord droit est irrégulier (alignement par défaut) ;
\item \cle{centré} : chaque ligne est centrée entre les deux marges (titres, invitations) ;
\item \cle{aligné à droite} : le texte touche le bord droit (date et lieu dans une lettre) ;
\item \cle{justifié} : le texte touche à la fois le bord gauche et le bord droit, le logiciel ajustant les espaces entre les mots (livres, journaux).
\end{itemize}
\end{definition}

\begin{popfigure}
\begin{center}
\begin{tikzpicture}[font=\small]
% cadre commun : largeur 10cm
\foreach \y/\txt in {
  0/{Aligné à gauche},
  -2.2/{Centré},
  -4.4/{Aligné à droite},
  -6.6/{Justifié}}{
  \draw[popline, thick] (0,\y) rectangle (10,\y-1.6);
  \node[anchor=north west, text=popdark, font=\bfseries\small] at (0,\y) {\txt};
}
% lignes simulées : gauche
\foreach \i/\w in {0/9.2, 1/8.4, 2/9.6, 3/5.1}{
  \fill[popblue!45] (0.2, -0.55-\i*0.28) rectangle (0.2+\w*0.95, -0.68-\i*0.28);}
% centré
\foreach \i/\w in {0/8.6, 1/9.0, 2/8.2, 3/4.6}{
  \fill[poporange!55] ({(10-\w*0.95)/2}, -2.75-\i*0.28) rectangle ({(10+\w*0.95)/2}, -2.88-\i*0.28);}
% droite
\foreach \i/\w in {0/9.0, 1/8.6, 2/9.4, 3/5.4}{
  \fill[popgreen!55] ({9.8-\w*0.95}, -4.95-\i*0.28) rectangle (9.8, -5.08-\i*0.28);}
% justifié
\foreach \i/\w in {0/9.6, 1/9.6, 2/9.6, 3/5.0}{
  \fill[poppurple!45] (0.2, -7.15-\i*0.28) rectangle (0.2+\w, -7.28-\i*0.28);}
\end{tikzpicture}
\end{center}
\legende{Les quatre alignements appliqués au même paragraphe : à gauche, centré, à droite et justifié.}
\end{popfigure}

\begin{methode}[Mettre en forme un paragraphe]
Pour appliquer une mise en forme à un paragraphe :
\begin{enumerate}
\item Clique n'importe où \textbf{dans} le paragraphe (inutile de tout sélectionner : les réglages de paragraphe s'appliquent au paragraphe entier où se trouve le curseur). Pour plusieurs paragraphes, sélectionne-les tous.
\item Ouvre l'onglet \texttt{Accueil}.
\item Dans le groupe \texttt{Paragraphe}, clique sur le bouton voulu : aligner à gauche, centrer, aligner à droite ou justifier.
\item Pour les réglages fins (retraits, espacements), clique sur la petite flèche en bas à droite du groupe \texttt{Paragraphe} pour ouvrir la boîte de dialogue complète.
\item Vérifie le résultat à l'écran, puis enregistre ton document.
\end{enumerate}
\end{methode}

\begin{attention}[Erreur fréquente : les faux alignements avec des espaces]
Beaucoup d'élèves « centrent » un titre en tapant plusieurs fois sur la barre d'\textbf{espace} avant le texte. C'est une très mauvaise méthode : dès que la taille des caractères, la police ou les marges changent, le titre n'est plus centré du tout. Il faut \textbf{toujours} utiliser les boutons d'alignement de l'onglet \texttt{Accueil}, qui garantissent un alignement exact et stable.
\end{attention}

\begin{savaistu}[Pourquoi les livres sont-ils justifiés ?]
L'alignement \textbf{justifié} est celui utilisé dans les livres, les journaux et les documents officiels (comme le Journal officiel de la République du Cameroun). En alignant le texte sur les deux marges, il forme des blocs réguliers qui donnent aux pages imprimées un aspect net et professionnel, et qui facilitent la lecture des longs textes.
\end{savaistu}

\courssub{Retraits et espacements}

\textbf{Intuition.} Quand tu rédiges une rédaction sur ta copie, le professeur te demande souvent de commencer chaque paragraphe « en retrait », c'est-à-dire un peu à l'intérieur de la marge, et de laisser de l'espace entre les paragraphes. Le traitement de texte fait exactement la même chose, mais de façon automatique et précise.

\begin{definition}[Retraits et espacements]
\begin{itemize}
\item Le \cle{retrait de première ligne} : la première ligne du paragraphe commence plus à droite que les autres (par exemple 1,25~cm).
\item Le \cle{retrait gauche} (ou droit) : tout le paragraphe est décalé par rapport à la marge gauche (ou droite) ; utile pour une citation.
\item L'\cle{espacement avant} et \cle{après} : l'espace vide laissé au-dessus et en dessous du paragraphe (par exemple 6~pt après).
\item L'\cle{interligne} : l'espace vertical entre les lignes d'un même paragraphe (simple, 1,5, double\ldots).
\end{itemize}
\end{definition}

\begin{popfigure}
\begin{center}
\begin{tikzpicture}[font=\small]
% marge gauche
\draw[popline, dashed] (0,0) -- (0,-4.6);
\node[anchor=east, text=popmuted, font=\footnotesize] at (0,-4.4) {marge};
% espacement avant
\draw[<->, poporange, thick] (1.1,-0.1) -- (1.1,-0.7);
\node[anchor=west, text=poporange, font=\footnotesize] at (1.25,-0.4) {espacement avant};
% paragraphe
\draw[popblue, thick] (0,-0.7) rectangle (9.5,-3.3);
% retrait première ligne
\draw[<->, popgreen, thick] (0,-0.85) -- (1.6,-0.85);
\node[text=popgreen, font=\footnotesize, anchor=south west] at (0.1,-0.82) {retrait de 1re ligne};
% lignes simulées
\fill[popblue!40] (1.6,-1.15) rectangle (9.3,-1.28);
\foreach \i in {1,2,3}{
  \fill[popblue!40] (0.2,{-1.15-\i*0.55}) rectangle (9.3,{-1.28-\i*0.55});}
\fill[popblue!40] (0.2,-3.35+0.55) rectangle (9.3,-3.48+0.55);
% interligne
\draw[<->, poppink, thick] (9.7,-1.15) -- (9.7,-1.70);
\node[anchor=west, text=poppink, font=\footnotesize] at (9.8,-1.42) {interligne};
% espacement après
\draw[<->, poporange, thick] (1.1,-3.3) -- (1.1,-4.0);
\node[anchor=west, text=poporange, font=\footnotesize] at (1.25,-3.65) {espacement après};
% paragraphe suivant
\draw[popline] (0,-4.0) rectangle (9.5,-4.5);
\end{tikzpicture}
\end{center}
\legende{Un paragraphe annoté : retrait de première ligne, interligne, espacements avant et après.}
\end{popfigure}

\begin{propriete}[Réglages usuels d'un paragraphe]
Dans la boîte de dialogue \texttt{Paragraphe} (onglet \texttt{Accueil}, groupe \texttt{Paragraphe}), on règle : le retrait de première ligne (zone \texttt{Retrait}, option \texttt{Première ligne}), les espacements \texttt{Avant} et \texttt{Après} en points, et l'\texttt{Interligne} (simple, 1,5 ligne, double\ldots). Pour une rédaction scolaire, on utilise souvent : retrait de première ligne de 1,25~cm, interligne 1,5, espacement après de 6~pt.
\end{propriete}

\begin{attention}[N'utilise jamais Entrée pour espacer les paragraphes]
Pour créer de l'espace entre deux paragraphes, ne tape pas plusieurs fois sur \texttt{Entrée} : cela crée des paragraphes vides qui perturbent le document (sauts de page imprévus, table des matières fausse). Utilise l'\textbf{espacement après} du paragraphe, réglé en points.
\end{attention}

\courssub{Listes à puces, listes numérotées et bordures}

\textbf{Intuition.} Pour énumérer le matériel d'un TP ou les étapes d'une recette, on ne rédige pas un long paragraphe : on présente les éléments les uns sous les autres, chacun précédé d'un symbole ou d'un numéro. Le traitement de texte automatise cette présentation.

\begin{definition}[Listes et bordures]
\begin{itemize}
\item Une \cle{liste à puces} présente des éléments sans ordre particulier, chacun précédé d'un symbole (point, carré, flèche\ldots).
\item Une \cle{liste numérotée} présente des éléments \textbf{ordonnés} (1, 2, 3\ldots), par exemple les étapes d'une démarche.
\item Une \cle{liste à plusieurs niveaux} combine les deux : des sous-éléments en retrait sous un élément principal (niveau 1, 1.1, 1.2\ldots).
\item La \cle{bordure} de paragraphe trace un cadre autour du paragraphe ; la \cle{trame de fond} le colore (boutons du groupe \texttt{Paragraphe}).
\end{itemize}
\end{definition}

\begin{methode}[Créer une liste à puces ou numérotée]
\begin{enumerate}
\item Sélectionne les lignes concernées (ou place le curseur et tape la première ligne).
\item Dans l'onglet \texttt{Accueil}, groupe \texttt{Paragraphe}, clique sur le bouton \texttt{Puces} ou \texttt{Numérotation}.
\item Tape chaque élément en validant par \texttt{Entrée} : la puce ou le numéro suivant apparaît automatiquement.
\item Pour créer un sous-niveau, appuie sur la touche \texttt{Tabulation} en début de ligne ; pour remonter d'un niveau, utilise \texttt{Maj + Tabulation}.
\item Pour terminer la liste, appuie deux fois sur \texttt{Entrée} ou clique à nouveau sur le bouton de liste.
\end{enumerate}
\end{methode}

\courssub{Reproduire et effacer la mise en forme ; mise en forme directe ou par styles}

\textbf{Intuition.} Tu viens de mettre un titre en Arial 16 gras bleu, et ton document contient dix autres titres identiques à traiter. Faut-il tout refaire dix fois ? Non : le bouton \textbf{Reproduire la mise en forme} (en forme de \textbf{pinceau}, dans l'onglet \texttt{Accueil}) permet de « copier » l'aspect d'un texte et de le « peindre » sur un autre. À l'inverse, le bouton \textbf{Effacer la mise en forme} (gomme avec la lettre A) supprime tous les attributs appliqués et ramène le texte à l'aspect par défaut.

\begin{methode}[Utiliser le pinceau de reproduction]
\begin{enumerate}
\item Sélectionne le texte dont la mise en forme te convient.
\item Clique une fois sur le bouton \texttt{Reproduire la mise en forme} (double-clic pour l'appliquer à plusieurs endroits).
\item « Peins » le texte cible en le sélectionnant : il adopte immédiatement la même police, taille, couleur et alignement.
\item Appuie sur \texttt{Échap} pour désactiver le pinceau.
\end{enumerate}
\end{methode}

\begin{definition}[Mise en forme directe et mise en forme par styles]
Tout ce que nous venons de voir est de la \cle{mise en forme directe} : on applique les attributs à la main, texte par texte. Elle est simple mais devient fastidieuse et source d'incohérences dans un long document. La \cle{mise en forme par styles} consiste à appliquer des styles prédéfinis (\texttt{Titre 1}, \texttt{Titre 2}, \texttt{Normal}\ldots) : le logiciel gère alors l'uniformité automatiquement et peut générer une \textbf{table des matières}. Nous l'étudierons en détail à la section 5.
\end{definition}

\begin{experience}[TP guidé : mettre en forme une convocation]
\textbf{Objectif :} produire la convocation à une réunion des parents d'élèves du collège. \textbf{Logiciel :} Word ou Writer.
\begin{enumerate}
\item Tape le titre « CONVOCATION » et le paragraphe : « Le Principal du collège convoque les parents d'élèves de 3ème à la réunion de rentrée qui se tiendra samedi à 9 h dans la salle polyvalente. »
\item Sélectionne le titre et applique : police Arial, taille 16, gras, centré.
\item Clique dans le paragraphe et applique : police Times New Roman, taille 12, alignement justifié, interligne 1,5, retrait de première ligne 1,25~cm, espacement après 6~pt.
\item Ajoute une liste à puces : « Relevé de notes », « Manuel de mathématiques », « Tenue correcte ».
\item Mets en évidence le mot « obligatoire » en gras et rouge, puis applique une trame de fond jaune clair au titre.
\item Enregistre le document sous le nom \texttt{convocation.odt} et fais vérifier par ton professeur.
\end{enumerate}
\textbf{Observation attendue :} le titre ressort clairement, le texte est justifié comme dans un livre, et la liste est automatiquement précédée de puces alignées.
\end{experience}

\begin{exoresolu}[Choisir les bons réglages]
Énoncé. Aïcha doit taper une lettre de demande de stage. Elle veut : (a) placer « Douala, le 12 janvier » contre la marge droite ; (b) écrire son objet en gras et souligné ; (c) présenter le corps de la lettre comme dans un livre, avec un alinéa à chaque paragraphe ; (d) espacer les lignes pour que le correcteur puisse annoter. Indique les réglages à appliquer.
\tcblower
Solution détaillée.
\begin{itemize}
\item (a) Cliquer dans le paragraphe de la date et choisir l'alignement \textbf{à droite} (et non des espaces !).
\item (b) Sélectionner l'objet et cliquer sur les boutons \textbf{Gras} puis \textbf{Souligné} du groupe \texttt{Police}.
\item (c) Sélectionner le corps de la lettre et appliquer l'alignement \textbf{justifié}, puis régler le \textbf{retrait de première ligne} à 1,25~cm dans la boîte de dialogue \texttt{Paragraphe}.
\item (d) Dans la même boîte de dialogue, choisir l'\textbf{interligne 1,5} (ou double) pour laisser de la place aux annotations.
\end{itemize}
\end{exoresolu}

\begin{exercice}[À toi de jouer]
\begin{enumerate}
\item Cite les quatre alignements de paragraphe et donne un exemple d'usage pour chacun.
\item Explique pourquoi il est interdit de centrer un titre avec la barre d'espace.
\item Quelle différence y a-t-il entre l'interligne et l'espacement après un paragraphe ?
\item Décris la démarche pour transformer cinq lignes de texte en liste numérotée à deux niveaux.
\end{enumerate}
\end{exercice}

\begin{corrige}[Exercice : à toi de jouer]
\begin{enumerate}
\item Gauche (texte courant de courrier électronique), centré (titre, invitation), droite (date et lieu d'une lettre), justifié (corps d'un livre ou d'un journal).
\item Les espaces ne garantissent pas la position : si la police, la taille ou les marges changent, le titre se décale. Le bouton \texttt{Centrer} centre exactement et durablement.
\item L'interligne est l'espace \emph{entre les lignes} d'un même paragraphe ; l'espacement après est l'espace \emph{sous} le paragraphe, avant le suivant.
\item Sélectionner les cinq lignes, cliquer sur \texttt{Numérotation} dans l'onglet \texttt{Accueil}, puis placer le curseur au début des lignes de sous-niveau et appuyer sur \texttt{Tabulation}.
\end{enumerate}
\end{corrige}

\begin{aretenir}[L'essentiel de la section]
\begin{itemize}
\item On met en forme les \textbf{caractères} (police, taille, gras, italique, souligné, couleur, exposant, indice) après les avoir \textbf{sélectionnés}.
\item Les quatre alignements sont : gauche, centré, droite et \textbf{justifié} (celui des livres).
\item Retraits, interligne et espacements se règlent dans la boîte de dialogue \texttt{Paragraphe} — jamais avec des espaces ni des \texttt{Entrée} répétées.
\item Les listes à puces et numérotées se créent automatiquement ; la tabulation gère les niveaux.
\item Le pinceau reproduit une mise en forme ; les \textbf{styles}, que nous verrons à la section 5, la rendent cohérente dans tout le document.
\end{itemize}
\end{aretenir}

\coursec{Mise en page du document}

Après avoir soigné l'apparence de votre texte — police, taille, couleur, alignement, interligne — il est temps de préparer la \emph{page} elle-même. Un beau paragraphe mal positionné sur la feuille, des marges trop étroites qui disparaissent à la reliure, ou un numéro de page manquant rendent un document difficile à lire, voire invalide pour un examen ou une administration. La \textbf{mise en page} définit le cadre physique de votre travail : c'est l'architecture du document avant même d'être imprimé.

\textbf{Orientation et format du papier}\par

\textbf{Orientation : portrait et paysage}\par

Par défaut, une page est en \textbf{orientation portrait} : la hauteur (\SI{29,7}{cm}) est supérieure à la largeur (\SI{21}{cm}). C'est le standard pour les lettres, les rapports, les dissertations. L'\textbf{orientation paysage} inverse les dimensions (largeur \SI{29,7}{cm}, hauteur \SI{21}{cm}) ; elle s'impose pour les grands tableaux, les schémas larges, les graphiques ou les certificats.

\begin{popfigure}
\begin{tikzpicture}[scale=0.35, every node/.style={font=\small}]
    % Portrait
    \draw[popblue, thick] (0,0) rectangle (21,29.7);
    \node[popblue, align=center] at (10.5, 14.85) {Portrait\\21 × 29,7 cm};
    \draw[<->, popmuted] (0,-1) -- (21,-1) node[midway, below] {21 cm};
    \draw[<->, popmuted] (-1,0) -- (-1,29.7) node[midway, left] {29,7 cm};
    
    % Paysage
    \begin{scope}[xshift=30cm]
        \draw[poporange, thick] (0,0) rectangle (29.7,21);
        \node[poporange, align=center] at (14.85, 10.5) {Paysage\\29,7 × 21 cm};
        \draw[<->, popmuted] (0,-1) -- (29.7,-1) node[midway, below] {29,7 cm};
        \draw[<->, popmuted] (-1,0) -- (-1,21) node[midway, left] {21 cm};
    \end{scope}
    
    % Flèche de rotation
    \draw[popgreen, thick, ->] (23, 14.85) arc (90:-90:3) node[midway, above, popgreen] {Rotation 90°};
\end{tikzpicture}
\legende{Comparatif des orientations portrait et paysage (format A4).}
\end{popfigure}

\begin{methode}[Changer l'orientation d'une page ou d'une section]
    \begin{enumerate}
        \item Onglet \texttt{Mise en page} (ou \texttt{Disposition} selon les versions).
        \item Groupe \texttt{Paramètres de page} $\rightarrow$ \texttt{Orientation}.
        \item Choisir \texttt{Portrait} ou \texttt{Paysage}.
        \item \textbf{Astuce :} Pour n'appliquer le paysage qu'à une seule page (ex. un grand tableau au milieu d'un rapport), placez le curseur au début de la page concernée, choisissez \texttt{Sauts} $\rightarrow$ \texttt{Page suivante} (Saut de section), changez l'orientation, puis remettez en portrait avec un nouveau saut de section à la fin.
        \item \textit{Note LibreOffice Writer :} \texttt{Format} $\rightarrow$ \texttt{Page} $\rightarrow$ onglet \texttt{Page} $\rightarrow$ \texttt{Orientation}.
    \end{enumerate}
\end{methode}

\textbf{Format du papier : A4, A5, Lettre}\par

Le \textbf{format} détermine les dimensions physiques de la feuille. Au Cameroun, comme dans la majeure partie du monde (norme ISO 216), le \textbf{format A4 (\SI{21}{\centi\meter} × \SI{29,7}{\centi\meter})} est la référence absolue pour l'administration, l'école et l'entreprise. Le format A5 (\SI{14,8}{\centi\meter} × \SI{21}{\centi\meter}) correspond à une demi-A4, utilisé pour les livrets ou flyers. Le format \emph{Letter} (\SI{21,59}{\centi\meter} × \SI{27,94}{\centi\meter}) est la norme nord-américaine ; évitez-le sauf demande expresse d'un correspondant étranger.

\begin{savaistu}[Le format A4 au Cameroun]
    Le format A4 est la norme légale pour les actes administratifs, les bulletins de notes, les fiches de paie, les convocations d'examen (BEPC, Probatoire, Baccalauréat) et la correspondance officielle. Une marge gauche de \SI{2,5}{\centi\meter} (ou \SI{3}{\centi\meter} pour reliure) est souvent exigée pour permettre la perforation et le classement dans les dossiers à anneaux des administrations (MINESEC, MINFOPRA, mairies, banques). Un document au format Letter ou avec des marges incorrectes peut être rejeté par un service de secrétariat.
\end{savaistu}

\begin{methode}[Régler le format du papier]
    \begin{enumerate}
        \item Onglet \texttt{Mise en page} $\rightarrow$ groupe \texttt{Paramètres de page} $\rightarrow$ \texttt{Taille}.
        \item Sélectionner \texttt{A4} (souvent premier de la liste).
        \item Si absent : \texttt{Autres formats de papier} $\rightarrow$ onglet \texttt{Papier} $\rightarrow$ \texttt{Format} : A4 (\SI{21}{\centi\meter} × \SI{29,7}{\centi\meter}).
        \item \textit{Note LibreOffice Writer :} \texttt{Format} $\rightarrow$ \texttt{Page} $\rightarrow$ onglet \texttt{Page} $\rightarrow$ \texttt{Format de papier} : A4.
    \end{enumerate}
\end{methode}

\textbf{Réglage des marges}\par

Les \textbf{marges} sont les espaces blancs réservés entre le texte et les bords physiques de la feuille. Elles assurent la lisibilité, permettent la reliure/perforation et évitent que l'imprimante ne coupe du contenu (zone non imprimable).

\begin{popfigure}
\begin{tikzpicture}[scale=0.45, every node/.style={font=\small}]
    % Feuille A4 Portrait
    \draw[popdark, thick] (0,0) rectangle (21,29.7);
    
    % Marges (exemple 2.5 cm)
    \def\m{2.5}
    \draw[popblue!50, dashed] (\m,\m) rectangle (21-\m, 29.7-\m);
    
    % Cotes marges
    \draw[<->, popblue] (0, -0.5) -- (\m, -0.5) node[midway, below] {Marge gauche : 2,5 cm};
    \draw[<->, popblue] (21-\m, -0.5) -- (21, -0.5) node[midway, below] {Marge droite : 2,5 cm};
    \draw[<->, popblue] (-0.5, 0) -- (-0.5, \m) node[midway, left] {Marge haute : 2,5 cm};
    \draw[<->, popblue] (-0.5, 29.7-\m) -- (-0.5, 29.7) node[midway, left] {Marge basse : 2,5 cm};
    
    % Zone en-tête / pied de page
    \draw[poporange, densely dotted] (\m, 29.7-1.5) -- (21-\m, 29.7-1.5);
    \node[poporange, above] at (10.5, 29.7-1.5) {Zone d'en-tête (ex. 1,5 cm depuis le haut)};
    
    \draw[poporange, densely dotted] (\m, 1.5) -- (21-\m, 1.5);
    \node[poporange, below] at (10.5, 1.5) {Zone de pied de page (ex. 1,5 cm depuis le bas)};
    
    % Corps de texte
    \node[align=center, text width=14cm, popmuted] at (10.5, 15) {Corps du document\\ (zone de saisie du texte)};
    
    % Légendes
    \node[popblue, rotate=90] at (-1.5, 14.85) {Hauteur 29,7 cm};
    \node[popblue] at (10.5, -1.8) {Largeur 21 cm};
\end{tikzpicture}
\legende{Schéma d'une page A4 portrait avec marges de 2,5 cm, zones d'en-tête et de pied de page repérées.}
\end{popfigure}

\begin{definition}[Marges]
    Les \textbf{marges} définissent la zone imprimable (cadre de texte) à l'intérieur de la page. On distingue :
    \begin{itemize}
        \item \textbf{Marge supérieure (haute)} et \textbf{inférieure (basse)} : espaces verticaux.
        \item \textbf{Marge gauche} et \textbf{droite} : espaces horizontaux.
        \item \textbf{Gouttière} : espace supplémentaire ajouté à la marge intérieure (gauche pour pages impaires, droite pour paires) pour la reliure.
    \end{itemize}
\end{definition}

\begin{exemplebox}[Réglage standard pour un rapport scolaire (A4 Portrait)]
    \begin{itemize}
        \item \textbf{Haute} : \SI{2,5}{\centi\meter} (ou \SI{3}{\centi\meter} si en-tête chargé).
        \item \textbf{Basse} : \SI{2,5}{\centi\meter}.
        \item \textbf{Gauche} : \SI{3}{\centi\meter} (prévoir perforation/reliure à gauche).
        \item \textbf{Droite} : \SI{2}{\centi\meter}.
        \item \textbf{Gouttière} : \SI{0,5}{\centi\meter} (position : Gauche).
    \end{itemize}
    Ces valeurs garantissent une présentation aérée et conforme aux exigences du jury du BEPC ou de l'administration camerounaise.
\end{exemplebox}

\begin{methode}[Régler les marges (prédéfinies ou personnalisées)]
    \begin{enumerate}
        \item Onglet \texttt{Mise en page} $\rightarrow$ groupe \texttt{Paramètres de page} $\rightarrow$ \texttt{Marges}.
        \item Choisir un préréglé (\texttt{Normal}, \texttt{Étroit}, \texttt{Modéré}, \texttt{Large}, \texttt{Miroir} pour recto-verso).
        \item Pour des valeurs exactes : \texttt{Marges personnalisées...} (bas de la liste) $\rightarrow$ onglet \texttt{Marges} $\rightarrow$ saisir les valeurs en cm (Haut, Bas, Gauche, Droite, Gouttière).
        \item \textbf{Appliquer à :} \texttt{Tout le document} (par défaut) ou \texttt{À partir de ce point} (après un saut de section).
        \item Valider par \texttt{OK}.
        \item \textit{Note LibreOffice Writer :} \texttt{Format} $\rightarrow$ \texttt{Page} $\rightarrow$ onglet \texttt{Page} $\rightarrow$ saisir les marges.
    \end{enumerate}
\end{methode}

\begin{attention}[Erreur fréquente : marges trop petites]
    Ne descendez pas en dessous de \textbf{\SI{1,5}{\centi\meter}} (sauf marge droite \SI{1}{\centi\meter}) : la zone non imprimable de la plupart des imprimantes (surtout les modèles d'entrée de gamme dans les cybercafés ou écoles) « mange » le texte. Résultat : bas de page coupé, numéros de page manquants, tableau tronqué. Testez toujours avec \texttt{Aperçu avant impression} avant d'imprimer 50 exemplaires.
\end{attention}

\textbf{En-tête et pied de page}\par

L'\textbf{en-tête} (zone haute) et le \textbf{pied de page} (zone basse) sont des zones répétées sur \emph{toutes} les pages (ou selon sections) en dehors du corps de texte principal. Ils portent l'identité du document.

\begin{definition}[En-tête et Pied de page]
    \begin{itemize}
        \item \textbf{En-tête :} zone située dans la marge haute. Contient typiquement : titre du document, nom de l'établissement (ex. \emph{Lycée Bilingue de Yaoundé}), logo, date, nom de l'auteur.
        \item \textbf{Pied de page :} zone située dans la marge basse. Contient typiquement : \textbf{numéro de page}, mention de confidentialité, chemin du fichier, nom de l'établissement.
    \end{itemize}
    Ces zones sont liées aux \textbf{styles de page} ; une modification sur une page se répercute sur toutes les pages du même style/section.
\end{definition}

\begin{methode}[Insérer un en-tête avec nom de l'établissement et date]
    \begin{enumerate}
        \item Onglet \texttt{Insertion} $\rightarrow$ groupe \texttt{En-tête et pied de page} $\rightarrow$ \texttt{En-tête} $\rightarrow$ choisir un modèle (ex. \texttt{Blank} / \texttt{Vide} ou \texttt{Alphabet}).
        \item Le curseur clignote dans la zone d'en-tête (ruban \texttt{Outils d'en-tête et de pied de page} / \texttt{Conception} s'affiche).
        \item Saisir le texte : \texttt{Lycée Bilingue de Yaoundé \textbar\ Rapport de Stage 2025}.
        \item Pour la date auto : onglet \texttt{Conception} $\rightarrow$ \texttt{Date et heure} $\rightarrow$ choisir format (ex. \texttt{14 mars 2025}) $\rightarrow$ cocher \texttt{Mettre à jour automatiquement} $\rightarrow$ \texttt{OK}.
        \item Fermer : \texttt{Fermer l'en-tête et le pied de page} (ou double-clic dans le corps du document).
        \item \textit{Note LibreOffice Writer :} \texttt{Insertion} $\rightarrow$ \texttt{En-tête et pied de page} $\rightarrow$ \texttt{En-tête} $\rightarrow$ \texttt{Style par défaut}.
    \end{enumerate}
\end{methode}

\textbf{Numérotation automatique des pages}\par

La numérotation manuelle (taper « 1 », « 2 »...) est une \textbf{erreur majeure} : à la moindre insertion de paragraphe, tout décale. La numérotation \textbf{automatique} est un champ dynamique qui se met à jour seul.

\begin{methode}[Insérer un numéro de page automatique]
    \begin{enumerate}
        \item Onglet \texttt{Insertion} $\rightarrow$ groupe \texttt{En-tête et pied de page} $\rightarrow$ \texttt{Numéro de page}.
        \item Choisir la position : \texttt{Haut de page} (en-tête), \texttt{Bas de page} (pied de page - \textbf{recommandé pour rapports}), \texttt{Marges de page} (latéral), \texttt{Position actuelle}.
        \item Choisir le style (alignement gauche, centre, droite ; avec ou sans trait de séparation, format « Page 1 sur 3 »).
        \item \textbf{Option « Page 1 sur X » :} insère le numéro courant ET le total des pages (champ \texttt{NUMPAGES}).
        \item Pour \textbf{ne pas numéroter la page de garde} : dans l'onglet \texttt{Conception} (actif dans l'en-tête/pied), cocher \texttt{Première page différente}. Le numéro 1 apparaîtra sur la page 2.
        \item Fermer l'en-tête/pied de page.
        \item \textit{Note LibreOffice Writer :} \texttt{Insertion} $\rightarrow$ \texttt{Numéro de page} (curseur dans pied de page) $\rightarrow$ \texttt{Format} pour « Page 1 sur N ».
    \end{enumerate}
\end{methode}

\begin{aretenir}[Numérotation et sections]
    Par défaut, la numérotation est continue sur tout le document. Si vous insérez un \textbf{saut de section} (page suivante), vous pouvez \textbf{redémarrer la numérotation} (ex. chiffres romains i, ii, iii pour la table des matières, puis arabes 1, 2, 3 pour le corps) via \texttt{Numéro de page} $\rightarrow$ \texttt{Format de numéros de page} $\rightarrow$ \texttt{Commencer à : 1} et choisir le format.
\end{aretenir}

\textbf{Sauts de page et sauts de section}\par

Pour forcer le passage à la page suivante, n'utilisez \textbf{jamais} la touche \texttt{Entrée} répétée. Le texte bougerait à la moindre modification.

\begin{definition}[Saut de page vs Saut de section]
    \begin{itemize}
        \item \textbf{Saut de page (\texttt{Ctrl + Entrée}) :} force le curseur au début de la page suivante. \emph{Mêmes mise en page, marges, en-tête/pied, numérotation.}
        \item \textbf{Saut de section (Page suivante) :} crée une nouvelle \textbf{section}. Permet de changer : orientation (portrait/paysage), marges, colonnes, en-tête/pied de page \emph{différents} (délier du précédent via \texttt{Lier au précédent}), numérotation redémarrée.
    \end{itemize}
\end{definition}

\begin{attention}[Le piège de la touche Entrée]
    \textbf{Ne faites jamais :} \texttt{Entrée} \texttt{Entrée} \texttt{Entrée}... pour « descendre » le titre du chapitre suivant en haut de page.
    \begin{itemize}
        \item Si vous ajoutez une ligne au chapitre 1, le titre du chapitre 2 « tombe » au milieu de la page 2.
        \item L'aperçu avant impression montre un espace blanc géant variable.
        \item \textbf{Solution :} Placez le curseur \emph{avant} le titre du chapitre 2 $\rightarrow$ \texttt{Mise en page} $\rightarrow$ \texttt{Sauts} $\rightarrow$ \texttt{Page suivante} (Saut de page). Le titre reste collé en haut de page quoi qu'il arrive.
    \end{itemize}
\end{attention}

\begin{experience}[TP guidé : Document mixte Portrait / Paysage]
    \textbf{Objectif :} Produire un rapport de 3 pages : page 1 (garde, portrait), page 2 (grand tableau, paysage), page 3 (suite texte, portrait), avec numérotation continue (page 1 non numérotée, page 2 = 1, page 3 = 2).
    
    \textbf{Matériel :} Word / LibreOffice Writer.
    
    \textbf{Manipulation :}
    \begin{enumerate}
        \item Page 1 : Saisir le titre. \texttt{Mise en page} $\rightarrow$ \texttt{Sauts} $\rightarrow$ \texttt{Page suivante} (Saut de section).
        \item Page 2 : Curseur en haut. \texttt{Mise en page} $\rightarrow$ \texttt{Orientation} $\rightarrow$ \texttt{Paysage}. Insérer le tableau large.
        \item Fin page 2 : \texttt{Mise en page} $\rightarrow$ \texttt{Sauts} $\rightarrow$ \texttt{Page suivante} (Saut de section).
        \item Page 3 : Curseur en haut. \texttt{Mise en page} $\rightarrow$ \texttt{Orientation} $\rightarrow$ \texttt{Portrait}. Saisir la suite.
        \item Pied de page page 2 : \texttt{Insertion} $\rightarrow$ \texttt{Numéro de page} $\rightarrow$ \texttt{Bas de page} $\rightarrow$ \texttt{Simple 2}. Onglet \texttt{Conception} : cocher \texttt{Lier au précédent} (désactiver) pour les sections 1 et 3 si besoin, mais ici on veut continu. \texttt{Numéro de page} $\rightarrow$ \texttt{Format} $\rightarrow$ \texttt{Commencer à : 1}.
        \item Page 1 : Pied de page $\rightarrow$ cocher \texttt{Première page différente} (vide).
    \end{enumerate}
    
    \textbf{Observations :} Le tableau tient sur une page large. La numérotation suit l'ordre logique (1, 2) et non l'ordre physique (2, 3).
    
    \textbf{Interprétation :} Le saut de section isole la mise en page (orientation) sans casser la continuité logique du document (numérotation, styles).
\end{experience}

\textbf{Aperçu avant impression et impression}\par

Avant toute impression, l'\textbf{Aperçu avant impression} (raccourci \texttt{Ctrl + P} ou \texttt{Fichier} $\rightarrow$ \texttt{Imprimer}) est \textbf{obligatoire}. C'est la dernière vérification : mise en page réelle, marges effectives, en-têtes/pieds visibles, numérotation correcte, pas de veuves/orphelines (dernière ligne d'un paragraphe seule en haut de page, ou première ligne seule en bas).

\begin{methode}[Vérifier et imprimer correctement]
    \begin{enumerate}
        \item \texttt{Ctrl + P} (ou \texttt{Fichier} $\rightarrow$ \texttt{Imprimer}).
        \item \textbf{Volet gauche (Paramètres) :}
            \begin{itemize}
                \item \texttt{Imprimer toutes les pages} / \texttt{Pages} (ex. 1, 3, 5-7) / \texttt{Sélection}.
                \item \texttt{Recto verso} : \texttt{Imprimer sur les deux faces - Retourner les feuilles sur le bord long} (standard livret) ou \texttt{bord court} (calendrier). Si imprimante non recto-verso auto : \texttt{Imprimer manuellement en recto verso}.
                \item \texttt{Pages par feuille} : 1 (standard), 2 (économie papier, lisible si police $\ge$ 10 pt).
                \item \textbf{Marges} : vérifier « Normales » ou « Étroites » selon votre réglage.
            \end{itemize}
        \item \textbf{Volet droit (Aperçu) :} Feuilleter avec les flèches. Zoomer. Vérifier : en-têtes/pieds, numéros, tableaux non coupés, images visibles.
        \item \textbf{Imprimante :} Sélectionner la bonne (ex. \texttt{HP LaserJet 1020}, \texttt{EPSON L3150 Series}, \texttt{Microsoft Print to PDF} pour créer un PDF).
        \item Cliquer \texttt{Imprimer}.
    \end{enumerate}
\end{methode}

\begin{exoresolu}[Correction d'un document mal mis en page]
    \textbf{Énoncé :} Un élève a rédigé son rapport de stage (15 pages). Il a utilisé la touche Entrée pour sauter les pages. Les marges sont par défaut (2,54 cm). L'en-tête contient son nom seulement. Le pied de page est vide. L'enseignant exige : marge gauche 3 cm (reliure), en-tête « Lycée de Bafoussam — Rapport de Stage 2025 », pied de page numérotation centrée « Page X sur Y », page de garde non numérotée.
    
    \textbf{Solution détaillée :}
    \begin{enumerate}
        \item \textbf{Marges :} \texttt{Mise en page} $\rightarrow$ \texttt{Marges} $\rightarrow$ \texttt{Marges personnalisées} $\rightarrow$ Gauche : 3 cm, Droite : 2 cm, Haut/Bas : 2,5 cm. \texttt{Appliquer à : Tout le document}. \texttt{OK}.
        \item \textbf{Nettoyer les sauts manuels :} Afficher les marques de formatage (\texttt{Ctrl + *} ou bouton \texttt{¶}). Supprimer les séries de \texttt{Entrée} (paragraphe vides) entre les chapitres. Remplacer par \texttt{Mise en page} $\rightarrow$ \texttt{Sauts} $\rightarrow$ \texttt{Page suivante} avant chaque titre de chapitre (Style \texttt{Titre 1}).
        \item \textbf{En-tête :} \texttt{Insertion} $\rightarrow$ \texttt{En-tête} $\rightarrow$ \texttt{Vide}. Taper \texttt{Lycée de Bafoussam — Rapport de Stage 2025}. \texttt{Conception} $\rightarrow$ cocher \texttt{Première page différente}.
        \item \textbf{Pied de page / Numérotation :} \texttt{Insertion} $\rightarrow$ \texttt{Pied de page} $\rightarrow$ \texttt{Vide}. \texttt{Conception} $\rightarrow$ \texttt{Numéro de page} $\rightarrow$ \texttt{Bas de page} $\rightarrow$ \texttt{Page X sur Y} (centré). Comme \texttt{Première page différente} est coché, la page 1 reste sans numéro. La page 2 affiche « Page 2 sur 15 ».
        \item \textbf{Corriger le numéro de départ (optionnel) :} Si on veut que la page 2 affiche « Page 1 sur 14 » : Pied de page page 2 $\rightarrow$ \texttt{Numéro de page} $\rightarrow$ \texttt{Format de numéros de page} $\rightarrow$ \texttt{Commencer à : 1}.
        \item \textbf{Aperçu :} \texttt{Ctrl + P} $\rightarrow$ vérifier pages 1 à 3. Imprimer en PDF pour valider.
    \end{enumerate}
\end{exoresolu}

\begin{exemplebox}[Raccourcis clavier utiles pour la mise en page]
    \begin{center}
    \begin{tabularx}{\linewidth}{|l|X|}
        \hline
        \rowcolor{popblueL}
        \textbf{Action} & \textbf{Raccourci (Word / Writer)} \\
        \hline
        Saut de page & \texttt{Ctrl + Entrée} \\
        \hline
        Ouvrir boîte de dialogue Mise en page (marges, papier, disposition) & \texttt{Alt + P, M} (puis flèches) / \texttt{Format} $\rightarrow$ \texttt{Page} (Writer) \\
        \hline
        Aperçu avant impression & \texttt{Ctrl + P} \\
        \hline
        Afficher/Masquer marques de formatage (¶) & \texttt{Ctrl + *} (pavé numérique) ou \texttt{Ctrl + Maj + 8} \\
        \hline
        Aller à l'en-tête / pied de page & Double-clic zone haute/basse \\
        \hline
    \end{tabularx}
    \end{center}
\end{exemplebox}

\textbf{Synthèse de la section}\par

La mise en page structure l'espace physique de votre document. Maîtrisez ces points clés pour le BEPC et vos futurs dossiers :
\begin{itemize}
    \item \textbf{Format A4 Portrait} par défaut ; \textbf{Paysage} seulement si nécessaire (tableau large), via \textbf{Saut de section}.
    \item \textbf{Marges} : Gauche 3 cm (reliure), autres 2--2,5 cm. Jamais $<$ 1,5 cm.
    \item \textbf{En-tête / Pied de page} : identité du document + \textbf{numérotation automatique} (champ dynamique, jamais saisi à la main).
    \item \textbf{Saut de page (\texttt{Ctrl+Entrée})} pour changer de page ; \textbf{Saut de section} pour changer de mise en page (orientation, marges, numérotation).
    \item \textbf{Aperçu avant impression (\texttt{Ctrl+P})} : \textbf{dernier contrôle obligatoire} avant d'imprimer ou d'exporter en PDF.
\end{itemize}

Dans la prochaine section, nous verrons comment insérer des \textbf{objets} (images, tableaux, formes, graphiques) pour enrichir vos documents.

\coursec{Insertion d'objets dans le document}

Après avoir structuré la page (marges, orientation, en-têtes et pieds de page), nous devons maintenant l'habiter. Un document professionnel — qu'il s'agisse d'un rapport d'exposé au collège de Nkolbisson, d'une liste de fournitures pour la rentrée, ou d'une affiche pour la kermesse de l'établissement — ne se limite pas à du texte aligné. Il a besoin d'images pour illustrer, de tableaux pour organiser des données chiffrées, de formes pour attirer l'œil, et de liens pour naviguer. Cette section te donne les clés pour insérer et maîtriser ces objets, sans que ton document ne devienne ingérable.

\courssub{Insertion et manipulation d'images}

Une image vaut mille mots, mais seulement si elle est bien placée. Dans un traitement de texte (LibreOffice Writer, Microsoft Word), l'insertion se fait via l'onglet \textbf{Insertion} \textgreater{} \textbf{Image} (ou \textbf{Illustrations} \textgreater{} \textbf{Images}). Tu peux choisir un fichier sur ton disque (\texttt{.jpg}, \texttt{.png}, \texttt{.gif}), ou coller une capture d'écran (touche \texttt{Impr. Écran} puis \texttt{Ctrl+V}).

\begin{experience}[Insérer et ajuster une photo de classe]
\textbf{Objectif :} Insérer une photo, la redimensionner sans la déformer, et rogner les bords inutiles.

\textbf{Matériel :} Un ordinateur, un traitement de texte, un fichier image \texttt{photo\_classe.jpg}.
\begin{enumerate}
    \item Place le curseur à l'endroit désiré (par exemple après le titre « Photo de la 3ème A »).
    \item Onglet \textbf{Insertion} \textgreater{} \textbf{Image} \textgreater{} \textbf{À partir d'un fichier} \textgreater{} sélectionne \texttt{photo\_classe.jpg} \textgreater{} \textbf{Ouvrir}.
    \item \textbf{Redimensionner :} clique sur l'image. Des poignées (petits carrés) apparaissent aux coins et au milieu des côtés. \textbf{Tire une poignée d'angle} en maintenant la touche \texttt{Maj} enfoncée pour conserver les proportions (le rapport largeur/hauteur). Relâche la souris avant la touche.
    \item \textbf{Rogner (recadrer) :} sélectionne l'image \textgreater{} onglet \textbf{Format} \textgreater{} \textbf{Rogner}. Des poignées noires épaisses remplacent les poignées blanches. Tire-les vers l'intérieur pour masquer les zones indésirables (le mur au fond, le bord du tableau). Clique en dehors de l'image pour valider.
\end{enumerate}
\textbf{Observation :} le rognage ne supprime pas les pixels du fichier source, il les masque seulement. Tu peux donc les récupérer plus tard en rouvrant l'outil Rogner et en écartant les poignées.
\end{experience}

\begin{attention}[L'erreur du déplacement « sauvage »]
Si tu cliques-glisses une image venant d'être insérée (par défaut \textbf{Alignée sur le texte}), elle se comporte comme un gros caractère : elle pousse les lignes, crée de grands blancs verticaux, et refuse de se placer exactement où tu veux (par exemple à gauche d'un paragraphe). \textbf{Solution :} change immédiatement son mode d'habillage (voir ci-dessous) en \textbf{Carré} avant de la déplacer.
\end{attention}

\courssub{L'habillage du texte autour d'une image : le secret de la mise en page}

C'est la notion \textbf{clé} pour passer du « devoir maison » au document professionnel. Par défaut, l'image est \emph{alignée sur le texte} : elle reste sur la ligne de base, comme une lettre géante. Pour que le texte coule \emph{autour} de l'image, il faut changer sa propriété \textbf{Habillage}.

\begin{definition}[Habillage du texte]
L'\cle{habillage} définit la relation géométrique entre le cadre de l'image (son rectangle englobant) et les lignes de texte environnantes. Il détermine si le texte passe au-dessus, en dessous, à côté, ou par-dessus l'image.
\end{definition}

\begin{popfigure}
\begin{tikzpicture}[scale=0.9, every node/.style={font=\footnotesize}]
    \tikzset{
        img/.style={draw=popdark, thick, fill=popcream, minimum width=1.6cm, minimum height=1.2cm, rounded corners=2pt},
        txt/.style={draw=popblue, fill=popblueL, minimum height=0.45cm, rounded corners=2pt, inner sep=2pt},
        lbl/.style={font=\bfseries\small, text=popdark}
    }
    % Colonne 1 : Aligné sur le texte
    \node[lbl] at (0,1.1) {Aligné sur le texte};
    \node[txt, minimum width=3.6cm] (a1) at (0,0.55) {Texte avant l'image...};
    \node[img] (a2) at (0,-0.35) {\textbf{IMAGE}};
    \node[txt, minimum width=3.6cm] (a3) at (0,-1.35) {Texte après : grand espace.};
    % Colonne 2 : Carré
    \node[lbl] at (5.6,1.1) {Carré};
    \node[img] (b1) at (6.4,-0.35) {\textbf{IMAGE}};
    \node[txt, minimum width=1.9cm] (b2) at (4.35,-0.05) {Le texte};
    \node[txt, minimum width=1.9cm] (b3) at (4.35,-0.65) {contourne};
    \node[txt, minimum width=3.6cm] (b4) at (5.6,-1.35) {puis reprend normalement.};
    % Colonne 3 : Derrière le texte
    \node[lbl] at (11.2,1.1) {Derrière le texte};
    \node[img, fill=poppinkL] (c1) at (11.2,-0.35) {};
    \node[txt, fill=popgreenL, draw=popgreen, minimum width=3.6cm] (c2) at (11.2,-0.05) {Le texte s'écrit par-dessus.};
    \node[txt, fill=popgreenL, draw=popgreen, minimum width=3.6cm] (c3) at (11.2,-0.65) {L'image sert de fond.};
    \node[txt, minimum width=3.6cm] (c4) at (11.2,-1.35) {Attention à la lisibilité !};
\end{tikzpicture}
\legende{Les trois modes d'habillage principaux : Aligné sur le texte (par défaut), Carré (le texte contourne l'image), Derrière le texte (image en fond, comme un filigrane).}
\end{popfigure}

\begin{propriete}[Modes d'habillage courants]
\begin{itemize}
    \item \textbf{Aligné sur le texte :} l'image est traitée comme un caractère. Elle ne flotte pas. Utile pour une petite icône dans une phrase.
    \item \textbf{Carré :} le texte entoure le rectangle de l'image. C'est le mode le plus utilisé pour les photos dans un rapport.
    \item \textbf{Rapproché :} le texte suit le contour \emph{réel} de l'image (si celle-ci possède des zones transparentes, comme un logo PNG). Idéal pour les formes irrégulières.
    \item \textbf{Derrière le texte :} l'image passe en arrière-plan. Sert pour les \textbf{filigranes} (par exemple le mot « CONFIDENTIEL » ou le logo de l'école en fond de page).
    \item \textbf{Devant le texte :} l'image recouvre le texte. Utile pour poser une annotation graphique (flèche, bulle) sur une capture d'écran.
\end{itemize}
\end{propriete}

\begin{methode}[Changer l'habillage d'une image]
\begin{enumerate}
    \item Clique sur l'image pour la sélectionner.
    \item Dans Word, clique sur le petit bouton \textbf{Options de mise en page} qui apparaît en haut à droite de l'image ; dans LibreOffice Writer, fais un clic droit \textgreater{} \textbf{Habillage}.
    \item Choisis \textbf{Carré} (recommandé pour commencer).
    \item Déplace l'image : le texte se réorganise instantanément autour d'elle.
    \item \emph{Pour aller plus loin :} ouvre les paramètres d'habillage pour régler la distance entre l'image et le texte (par exemple 0,3 cm de chaque côté).
\end{enumerate}
\end{methode}

\begin{savaistu}[Le filigrane officiel]
Au Cameroun, les bulletins de notes officiels, les diplômes du BEPC et les actes d'état civil portent souvent un filigrane (par exemple « RÉPUBLIQUE DU CAMEROUN » ou des armoiries) placé \textbf{derrière le texte}. Cela garantit l'authenticité visuelle du document sans gêner la lecture. Tu peux en créer un toi-même : onglet \textbf{Mise en page} (ou \textbf{Format} \textgreater{} \textbf{Page} dans Writer) \textgreater{} \textbf{Filigrane} \textgreater{} \textbf{Filigrane personnalisé} \textgreater{} Image ou Texte.
\end{savaistu}

\courssub{Insertion et mise en forme d'un tableau}

Le tableau est l'outil \emph{incontournable} pour présenter des données structurées : listes d'élèves avec moyennes, emplois du temps, budgets, états de stock. Contrairement à un tableau de tableur (Calc, Excel), le tableau du traitement de texte est essentiellement statique, mais il s'intègre parfaitement dans le flux du document.

\begin{methode}[Créer un tableau standard (exemple : effectifs par classe)]
\begin{enumerate}
    \item Place le curseur à l'endroit voulu.
    \item Onglet \textbf{Insertion} \textgreater{} \textbf{Tableau}.
    \item \textbf{Méthode visuelle :} survole la grille pour sélectionner \textbf{4 colonnes $\times$ 5 lignes} (1 ligne d'en-tête + 4 classes), puis clique.
    \item \textbf{Méthode précise :} \textbf{Insérer un tableau...} \textgreater{} saisis \texttt{Nombre de colonnes : 4}, \texttt{Nombre de lignes : 5} \textgreater{} \textbf{OK}.
    \item Remplis la ligne 1 (en-têtes) : \texttt{Classe}, \texttt{Garçons}, \texttt{Filles}, \texttt{Total}.
    \item Remplis les données ligne par ligne (touche \texttt{Tab} pour passer à la cellule suivante).
\end{enumerate}
\end{methode}

\begin{exemplebox}[Tableau des effectifs — Collège de Biyem-Assi]
Voici le rendu visé : bordures visibles, ligne d'en-tête colorée, nombres centrés.
\begin{center}
\begin{tabularx}{\linewidth}{|>{\bfseries}X|>{\centering}X|>{\centering}X|>{\centering\arraybackslash\bfseries}X|}
\hline
\rowcolor{popblueL}
\textbf{Classe} & \textbf{Garçons} & \textbf{Filles} & \textbf{Total} \\
\hline
3ème A & 28 & 32 & 60 \\
\hline
3ème B & 25 & 30 & 55 \\
\hline
3ème C & 30 & 28 & 58 \\
\hline
3ème D & 22 & 26 & 48 \\
\hline
\end{tabularx}
\end{center}
\textbf{Astuce :} sélectionne la ligne d'en-tête \textgreater{} clic droit \textgreater{} \textbf{Propriétés du tableau} \textgreater{} onglet \textbf{Ligne} \textgreater{} coche \textbf{Répéter en haut de chaque page comme ligne d'en-tête}. C'est indispensable pour les tableaux qui s'étendent sur plusieurs pages.
\end{exemplebox}

\begin{popfigure}
\begin{tikzpicture}[scale=0.85, every node/.style={font=\small}]
    \tikzset{
        cell/.style={draw=popdark, minimum height=0.8cm, inner sep=2pt, align=center},
        header/.style={cell, fill=popblueL, font=\bfseries, minimum width=1.9cm},
        merged/.style={cell, fill=poporangeL, font=\bfseries, minimum width=8.2cm},
        normal/.style={cell, fill=white, minimum width=1.9cm},
        total/.style={cell, fill=popgreenL, font=\bfseries, minimum width=1.9cm}
    }
    % Ligne 1 fusionnée
    \node[merged] (h) at (0,0) {EFFECTIFS DU COLLÈGE — ANNÉE 2025-2026};
    % Ligne 2 : en-têtes
    \node[header] (c1) at (-3.15,-0.9) {Classe};
    \node[header] (c2) at (-1.05,-0.9) {Garçons};
    \node[header] (c3) at (1.05,-0.9) {Filles};
    \node[header] (c4) at (3.15,-0.9) {Total};
    % Ligne 3
    \node[normal] at (-3.15,-1.7) {3ème A};
    \node[normal] at (-1.05,-1.7) {28};
    \node[normal] at (1.05,-1.7) {32};
    \node[total] at (3.15,-1.7) {60};
    % Ligne 4
    \node[normal] at (-3.15,-2.5) {3ème B};
    \node[normal] at (-1.05,-2.5) {25};
    \node[normal] at (1.05,-2.5) {30};
    \node[total] at (3.15,-2.5) {55};
    % Annotations
    \draw[poporange, thick, ->] (-4.6,-0.45) -- (-4.6,-0.05) node[midway, left, font=\tiny, align=right] {Cellules\\fusionnées};
    \draw[popblue, thick, ->] (4.6,-0.9) -- (4.15,-0.9) node[left=0pt, font=\tiny] {};
    \node[font=\tiny, text=popblue, anchor=west] at (4.7,-0.9) {Ligne d'en-tête};
    \node[font=\tiny, text=popgreen!60!black, anchor=west] at (4.7,-1.7) {Colonne Total};
    \draw[popgreen!60!black, thick, ->] (4.65,-1.7) -- (4.15,-1.7);
\end{tikzpicture}
\legende{Structure d'un tableau 4 colonnes $\times$ 4 lignes, avec fusion de la première ligne pour le titre général.}
\end{popfigure}

\textbf{Fusion et fractionnement de cellules.}\par
La \cle{fusion} combine plusieurs cellules adjacentes en une seule grande cellule. C'est indispensable pour les titres de tableau (comme ci-dessus) ou pour regrouper des catégories (par exemple une ligne « Total général » sur 4 colonnes).
\begin{itemize}
    \item \textbf{Fusionner :} sélectionne les cellules \textgreater{} clic droit \textgreater{} \textbf{Fusionner les cellules}.
    \item \textbf{Fractionner :} l'opération inverse. Place le curseur dans la cellule fusionnée \textgreater{} clic droit \textgreater{} \textbf{Fractionner les cellules} \textgreater{} choisis le nombre de lignes ou de colonnes.
\end{itemize}

\begin{exoresolu}[Créer le tableau du conseil de classe]
\textbf{Énoncé :} Le principal te demande de préparer le tableau de saisie pour le conseil de classe de 3ème A (45 élèves). Le tableau doit contenir les colonnes : \texttt{Nom}, \texttt{Prénom}, \texttt{Moyenne} (sur 20), \texttt{Rang}, \texttt{Décision} (Admis / Ajourné). Prévois une ligne d'en-tête et 45 lignes de saisie, avec un style de tableau à lignes alternées.
\tcblower
\textbf{Solution détaillée :}
\begin{enumerate}
    \item \textbf{Insertion} \textgreater{} \textbf{Tableau} \textgreater{} \textbf{Insérer un tableau...} \textgreater{} Colonnes : \texttt{5}, Lignes : \texttt{46} (1 en-tête + 45 élèves) \textgreater{} \textbf{OK}.
    \item Remplir l'en-tête (ligne 1) : \texttt{Nom} | \texttt{Prénom} | \texttt{Moyenne} | \texttt{Rang} | \texttt{Décision}.
    \item Sélectionner tout le tableau (clic sur la petite croix en haut à gauche du tableau).
    \item Onglet \textbf{Création de tableau} (ou \textbf{Conception}) \textgreater{} galerie de styles \textgreater{} choisir un style à lignes alternées.
    \item \textbf{Alignement :} sélectionne les colonnes « Moyenne » et « Rang » \textgreater{} onglet \textbf{Disposition} \textgreater{} \textbf{Centrer}.
    \item \textbf{Largeur :} double-clique sur la bordure droite de la colonne « Nom » pour l'ajuster automatiquement au contenu.
    \item \textbf{Option importante :} sélectionne la ligne d'en-tête \textgreater{} \textbf{Propriétés du tableau} \textgreater{} \textbf{Ligne} \textgreater{} coche \textbf{Répéter comme ligne d'en-tête} : ainsi l'en-tête réapparaît en haut de la page 2.
\end{enumerate}
\textbf{Résultat :} un tableau propre et lisible, prêt à être imprimé en deux exemplaires (un pour le jury, un pour l'affichage).
\end{exoresolu}

\begin{attention}[Le tableau d'un traitement de texte n'est pas un tableur]
Dans Word ou Writer, on peut saisir une petite formule dans une cellule (onglet \textbf{Disposition} \textgreater{} \textbf{Formule}, par exemple une somme des cellules au-dessus), mais le résultat \textbf{ne se met pas à jour automatiquement} quand on modifie une valeur : il faut sélectionner le champ et appuyer sur \texttt{F9} pour l'actualiser. Pour de vrais calculs, on utilise un tableur (Calc, Excel). \textbf{Au niveau 3ème, on saisit dans le tableau des résultats déjà calculés.}
\end{attention}

\courssub{Formes, zones de texte et WordArt : la touche visuelle}

Pour une affiche (par exemple « Journée Portes Ouvertes » ou « Inscriptions BEPC 2026 »), le texte linéaire ne suffit pas. Les \textbf{formes} (rectangles, flèches, étoiles, bulles) et le \textbf{WordArt} (texte stylisé avec effets d'ombre, de relief ou de couleur) transforment une page en support de communication.

\begin{methode}[Créer une affiche simple avec formes et WordArt]
\begin{enumerate}
    \item \textbf{Fond :} \textbf{Insertion} \textgreater{} \textbf{Formes} \textgreater{} \textbf{Rectangle} \textgreater{} dessine un rectangle couvrant la zone voulue. Clic droit \textgreater{} \textbf{Format de la forme} \textgreater{} \textbf{Remplissage} : couleur ou dégradé ; \textbf{Contour} : aucun.
    \item \textbf{Titre principal :} \textbf{Insertion} \textgreater{} \textbf{WordArt} \textgreater{} choisis un style. Tape : \texttt{JOURNÉE PORTES OUVERTES}. Agrandis avec les poignées d'angle et place le titre en haut, au centre.
    \item \textbf{Informations clés (zone de texte) :} \textbf{Insertion} \textgreater{} \textbf{Zone de texte} \textgreater{} dessine la zone. Tape : \texttt{Samedi 15 mars 2026, de 8 h à 16 h}. Mets en forme : police épaisse, taille 24 pt, couleur contrastée.
    \item \textbf{Accents visuels :} ajoute une \textbf{étoile} de chaque côté du titre, et une \textbf{flèche} pointant vers l'information essentielle.
    \item \textbf{Ordre des objets :} clic droit sur un objet \textgreater{} \textbf{Mettre en arrière-plan} (pour le fond) ou \textbf{Mettre au premier plan} (pour le titre).
\end{enumerate}
\end{methode}

\begin{propriete}[Zone de texte et WordArt]
\begin{itemize}
    \item \textbf{Zone de texte :} conteneur flottant que l'on peut positionner librement n'importe où sur la page (marges, entre colonnes). Idéal pour les mises en page de type affiche ou magazine.
    \item \textbf{WordArt :} objet \emph{texte} doté d'effets graphiques (déformations en arche, en vague, ombres, dégradés). Il se manipule comme une image et se redimensionne sans perte de qualité.
\end{itemize}
\end{propriete}

\courssub{Symboles, caractères spéciaux et liens hypertextes}

Un document rigoureux respecte la typographie. Écrire « Ecole » sans accent, « 5000 F » avec un espace ordinaire, ou « (c) » au lieu de « © » sont des maladresses à éviter.

\begin{definition}[Caractères spéciaux et symboles]
Les \cle{caractères spéciaux} sont des caractères non accessibles directement au clavier : lettres accentuées majuscules (\texttt{É}, \texttt{Ç}, \texttt{À}), symboles juridiques (\texttt{©}, \texttt{®}), opérateurs mathématiques (\texttt{×}, \texttt{÷}, \texttt{±}, \texttt{$\leq$}, \texttt{$\geq$}), tirets longs, et les espaces particuliers comme l'\textbf{espace insécable}. Pour la monnaie camerounaise, il n'existe pas de symbole unique : on écrit simplement \texttt{FCFA} (code officiel : \texttt{XAF}).
\end{definition}

\begin{methode}[Insérer un caractère spécial (exemple : © ou É majuscule)]
\begin{enumerate}
    \item Onglet \textbf{Insertion} \textgreater{} \textbf{Symbole} \textgreater{} \textbf{Autres symboles...}
    \item Choisis la police (généralement \texttt{(police normale)}) et le sous-ensemble : \textbf{Latin étendu} pour É, Ç, Œ ; \textbf{Opérateurs mathématiques} pour $\leq$, $\geq$, $\times$.
    \item Double-clique sur le caractère voulu \textgreater{} \textbf{Insérer} \textgreater{} \textbf{Fermer}.
    \item \textbf{Raccourcis utiles (Windows) :}
    \begin{itemize}
        \item \texttt{É} : \texttt{Alt + 144} (pavé numérique) ;
        \item \texttt{Ç} : \texttt{Alt + 128} ;
        \item \texttt{©} : \texttt{Alt + 0169} ;
        \item \textbf{Espace insécable} : \texttt{Ctrl + Maj + Espace}.
    \end{itemize}
\end{enumerate}
\end{methode}

\begin{aretenir}[L'espace insécable, réflexe du document propre]
Mets toujours une \textbf{espace insécable} (\texttt{Ctrl + Maj + Espace}) :
\begin{itemize}
    \item avant les signes doubles \texttt{: ; ! ?} et avant \texttt{\%} ;
    \item entre un nombre et son unité : \texttt{15~000~FCFA}, \texttt{12~km}, \texttt{20~\%} ;
    \item dans les noms composés : \texttt{Jean~Paul}, \texttt{Rue~de~la~Paix}.
\end{itemize}
Ainsi, « 15 000 FCFA » ne sera jamais coupé en fin de ligne : le groupe reste uni.
\end{aretenir}

\textbf{Le lien hypertexte.}\par
Le \cle{lien hypertexte} transforme un mot ou une image en bouton cliquable menant vers une adresse Internet (URL), une adresse e-mail, ou un autre endroit du document (signet).

\begin{methode}[Insérer un lien hypertexte]
\begin{enumerate}
    \item Sélectionne le texte (par exemple « Site du MINESEC ») ou l'image (un logo).
    \item \textbf{Insertion} \textgreater{} \textbf{Lien hypertexte} (ou raccourci \texttt{Ctrl + K}).
    \item \textbf{Vers un site web :} colle l'adresse complète, par exemple \texttt{https://www.minesec.gov.cm}.
    \item \textbf{Vers un e-mail :} saisis \texttt{mailto:} suivi de l'adresse, par exemple \texttt{mailto:principal@collegebiyem.cm}. Le clic ouvre le logiciel de messagerie.
    \item \textbf{Vers l'intérieur du document :} choisis un titre ou un signet créé au préalable (\textbf{Insertion} \textgreater{} \textbf{Signet}).
    \item Clique sur \textbf{OK}. Pour suivre le lien : \texttt{Ctrl + clic} sur le texte souligné.
\end{enumerate}
\end{methode}

\begin{exemplebox}[Lien interne : retour au sommaire]
En bas d'une longue page, tu peux insérer un lien « Retour au sommaire » pointant vers un signet placé au début du document. C'est le principe de navigation utilisé dans les documents numériques professionnels et les pages web.
\end{exemplebox}

\courssub{Synthèse de la section : check-list de l'insertion d'objets}

Avant de passer à la section suivante (styles de titres et table des matières), vérifie que tu sais faire tout ceci sans aide.

\begin{exercice}[Auto-évaluation — Insertion d'objets]
\begin{enumerate}
    \item Tu insères le logo de ton collège (image PNG) dans l'en-tête. Quel habillage choisir pour qu'il ne pousse pas le titre du document ?
    \item Tu crées un tableau « Notes du 1er trimestre » qui s'étend sur 3 pages. Quelle option activer sur la ligne d'en-tête ?
    \item Tu écris « Le prix est de 15 000 FCFA ». Quel raccourci utiliser pour les espaces à l'intérieur de ce groupe ?
    \item Comment insérer le symbole « $\geq$ » (supérieur ou égal) ?
    \item Cite les deux méthodes pour créer un tableau de 5 colonnes et 46 lignes.
\end{enumerate}
\end{exercice}

\begin{corrige}[Exercice n°1 -- Auto-évaluation]
\begin{enumerate}
    \item \textbf{Derrière le texte} (ou \textbf{Devant le texte} avec un positionnement précis) : dans l'en-tête, l'image ne doit pas perturber le flux du document.
    \item Clic droit sur la ligne d'en-tête \textgreater{} \textbf{Propriétés du tableau} \textgreater{} onglet \textbf{Ligne} \textgreater{} cocher \textbf{Répéter en haut de chaque page comme ligne d'en-tête}.
    \item \texttt{Ctrl + Maj + Espace} (espace insécable) entre « 15 » et « 000 », puis entre « 000 » et « FCFA » : le groupe « 15 000 FCFA » reste insécable.
    \item \textbf{Insertion} \textgreater{} \textbf{Symbole} \textgreater{} \textbf{Autres symboles...} \textgreater{} sous-ensemble \textbf{Opérateurs mathématiques} \textgreater{} double-cliquer sur $\geq$ \textgreater{} \textbf{Insérer}.
    \item Méthode visuelle : \textbf{Insertion} \textgreater{} \textbf{Tableau} \textgreater{} survoler la grille. Méthode précise : \textbf{Insertion} \textgreater{} \textbf{Tableau} \textgreater{} \textbf{Insérer un tableau...} \textgreater{} saisir 5 colonnes et 46 lignes \textgreater{} \textbf{OK}.
\end{enumerate}
\end{corrige}

Tu maîtrises maintenant l'insertion d'images (et leur habillage, notion capitale), la construction de tableaux structurés (fusion, répétition d'en-tête), l'usage des formes et du WordArt pour l'affichage, la rigueur typographique (symboles, espaces insécables) et les liens hypertextes. La section suivante te montrera comment automatiser la structure de ton document grâce aux \textbf{styles de titres} et générer une \textbf{table des matières} en un clic — la touche finale du document professionnel.

\coursec{Styles de titres et table des matières automatique}

Après avoir inséré images, tableaux et formes dans notre document (section précédente), nous devons maintenant l'organiser pour qu'il soit lisible, professionnel et facile à naviguer. Imagine un rapport de stage de 30 pages sans sommaire : ton jury perdrait un temps précieux à chercher l'introduction, la méthodologie ou la conclusion. Pire, si tu insères une page au milieu, tous les numéros de page deviennent faux. La solution professionnelle, c'est d'utiliser les \textbf{styles de titres} pour structurer le document, puis de laisser le logiciel générer et mettre à jour la \textbf{table des matières} automatiquement.

\courssub{La notion de style : séparer le fond de la forme}

\begin{definition}[Style de paragraphe]
Un \cle{style} est un ensemble nommé de caractéristiques de mise en forme (police, taille, couleur, alignement, interligne, espacement avant/après, tabulations, bordures, numérotation) que l'on applique en un clic à un paragraphe entier. Modifier le style met à jour \emph{tous} les paragraphes qui l'utilisent.
\end{definition}

Dans un traitement de texte (LibreOffice Writer, Microsoft Word), les styles sont organisés par catégories : \textbf{Styles de titre} (\texttt{Titre 1}, \texttt{Titre 2}, \texttt{Titre 3}...), \textbf{Style par défaut} (\texttt{Normal} ou \texttt{Standard}) pour le corps de texte, \textbf{Styles de liste}, \textbf{Styles de cadre}, etc. L'intérêt pédagogique et professionnel est double :
\begin{enumerate}
    \item \textbf{Cohérence visuelle :} tous les chapitres (\texttt{Titre 1}) ont exactement la même apparence sans effort de mémorisation.
    \item \textbf{Structure sémantique :} le logiciel « comprend » la hiérarchie du document (Chapitre > Section > Sous-section). C'est cette structure qui alimente la table des matières, le volet de navigation et l'accessibilité (lecteurs d'écran pour malvoyants).
\end{enumerate}

\begin{attention}[Erreur fréquente : mise en forme manuelle vs styles]
Ne sélectionne pas un titre pour le mettre en \textbf{Gras}, \texttt{Taille 16}, \texttt{Centré} manuellement. Cela crée une « fausse » structure : le logiciel ne voit qu'un paragraphe normal en gras. Conséquence : il n'apparaîtra \emph{jamais} dans la table des matières automatique, le volet de navigation restera vide, et si tu changes d'avis sur la police des titres, tu devras les corriger un par un. \textbf{Toujours utiliser les styles Titre 1, Titre 2, etc.}
\end{attention}

\courssub{Structurer son document avec les styles de titres}

La hiérarchie standard suit le plan de ton rapport :
\begin{itemize}
    \item \texttt{Titre 1} : Chapitres principaux (ex. \textbf{1. Introduction}, \textbf{2. Méthodologie}).
    \item \texttt{Titre 2} : Sections à l'intérieur d'un chapitre (ex. \textbf{2.1 Population d'étude}, \textbf{2.2 Outils de collecte}).
    \item \texttt{Titre 3} : Sous-sections (ex. \textbf{2.2.1 Questionnaire}, \textbf{2.2.2 Entretien}).
    \item \texttt{Normal} (ou \texttt{Standard}) : Corps de texte, paragraphes courants.
\end{itemize}

\begin{methode}[Appliquer un style de titre]
\begin{enumerate}
    \item Place le curseur n'importe où dans la ligne du titre (pas besoin de tout sélectionner, le style s'applique au paragraphe entier).
    \item Dans l'onglet \textbf{Accueil} (ou \textbf{Styles} dans la barre latérale), repère la galerie de styles.
    \item Clique sur \textbf{Titre 1}, \textbf{Titre 2} ou \textbf{Titre 3} selon le niveau hiérarchique.
    \item \textit{Astuce raccourci :} \texttt{Ctrl+Alt+1} pour Titre 1, \texttt{Ctrl+Alt+2} pour Titre 2, \texttt{Ctrl+Alt+3} pour Titre 3 (Word) ; \texttt{Ctrl+1}, \texttt{Ctrl+2}, \texttt{Ctrl+3} (LibreOffice Writer).
\end{enumerate}
\end{methode}

\begin{exemplebox}[Structure d'un rapport de stage en 3ème]
Voici à quoi ressemble la structure (volet de navigation) d'un rapport bien stylé :
\begin{lstlisting}[language=text, basicstyle=\ttfamily\small]
▸ 1. Introduction                          (Titre 1)
▸ 2. Présentation de l'entreprise          (Titre 1)
    ▸ 2.1 Historique                       (Titre 2)
    ▸ 2.2 Organigramme                     (Titre 2)
        ▸ 2.2.1 Direction générale         (Titre 3)
        ▸ 2.2.2 Services techniques        (Titre 3)
    ▸ 2.3 Effectifs                        (Titre 2)
▸ 3. Activités réalisées                   (Titre 1)
    ▸ 3.1 Maintenance poste de travail     (Titre 2)
    ▸ 3.2 Installation réseau              (Titre 2)
▸ 4. Conclusion                            (Titre 1)
▸ Bibliographie                            (Titre 1)
\end{lstlisting}
Remarque : la numérotation automatique (1, 2.1, 2.2.1) est gérée par les styles si on a activé une \textbf{numérotation de chapitres} (Outils > Numérotation des chapitres dans LibreOffice ; onglet Accueil > Liste multiniveau dans Word).
\end{exemplebox}

\courssub{Modifier un style : la puissance de la mise à jour globale}

\begin{propriete}[Modification en cascade]
Quand tu modifies la définition d'un style (ex. changer la police de \texttt{Titre 1} de \texttt{Arial} à \texttt{Times New Roman}, ou augmenter l'espace avant), \textbf{tous} les paragraphes portant ce style dans le document se mettent à jour instantanément. C'est le principe de la \emph{séparation du contenu et de la présentation}.
\end{propriete}

\begin{methode}[Modifier un style existant (deux méthodes équivalentes)]
\begin{enumerate}
    \item \textbf{Par le volet Styles (recommandé) :} Ouvre le volet \textbf{Styles} (\texttt{F11} dans LibreOffice, \texttt{Ctrl+Alt+Maj+S} dans Word). Clic droit sur \texttt{Titre 1} > \textbf{Modifier...}. Change les attributs (Police, Effets de police, Indents \& Espacement, Organisateur). Valide par \textbf{OK}.
    \item \textbf{Par l'exemple (mise à jour par sélection) :} Formate un titre existant exactement comme tu le souhaites (sélectionne-le, applique tes mises en forme). Sans le désélectionner, clic droit sur \texttt{Titre 1} dans la galerie > \textbf{Mettre à jour « Titre 1 » à partir de la sélection}.
\end{enumerate}
\end{methode}

\courssub{La table des matières automatique : génération et mise à jour}

\begin{definition}[Table des matières automatique]
Une \cle{table des matières} (TdM) est une liste structurée, placée généralement en début de document, qui récapitule les titres hiérarchiques (Titre 1, Titre 2, Titre 3...) avec leur numéro de page courant. Elle est \textbf{générée dynamiquement} à partir des styles de titres : si tu changes un titre ou sa page, la TdM se met à jour sur commande.
\end{definition}

\begin{propriete}[Condition nécessaire et suffisante]
\textbf{La table des matières automatique ne fonctionne que si les titres du document sont formatés avec les styles de titres (Titre 1, Titre 2, etc.).} Sans styles de titres appliqués, le menu de génération reste vide ou insère une table vide. La mise en forme manuelle (gras, taille) ne suffit pas.
\end{propriete}

\begin{methode}[Générer la table des matières en 3 clics]
\begin{enumerate}
    \item Place le curseur à l'endroit où insérer la table (généralement en page 2, après la page de garde).
    \item Onglet \textbf{Références} (Word) ou \textbf{Insertion} > \textbf{Table des matières et index} > \textbf{Table des matières, index ou bibliographie} (LibreOffice).
    \item Clique sur \textbf{Table des matières automatique 1} (ou \textbf{OK} avec les options par défaut : Niveaux = 3, Numéros de page activés).
\end{enumerate}
\end{methode}

\begin{attention}[Piège classique : oublier de mettre à jour]
Après avoir généré la table, tu continues ton travail : tu ajoutes un chapitre, tu modifies un titre, le texte déborde sur la page suivante... \textbf{La table des matières ne se met pas à jour toute seule !} Elle affiche encore l'ancien état. \textbf{Obligation :} Clic droit dans la table > \textbf{Mettre à jour les champs} (ou touche \texttt{F9}) > Choisir \textbf{Mettre à jour toute la table} > \textbf{OK}. Fais-le systématiquement avant d'imprimer ou d'exporter en PDF.
\end{attention}

\begin{experience}[TP guidé : Structurer et générer la TdM d'un rapport]
\textbf{Objectif :} Produire un mini-rapport de 3 pages avec styles hiérarchiques et table des matières fonctionnelle.
\textbf{Logiciel :} LibreOffice Writer ou Microsoft Word.
\textbf{Consignes :}
\begin{enumerate}
    \item Crée un nouveau document. Enregistre-le \texttt{Rapport\_Stage\_3eme.odt}.
    \item Page 1 : Titre principal « Rapport de stage au Lycée Bilingue de Yaoundé » (style \texttt{Titre 1} centré, taille 24). Auteur, Classe, Date (style \texttt{Normal}).
    \item Insère un \textbf{saut de page} (\texttt{Ctrl+Entrée}).
    \item Page 2 : Insère la table des matières (Méthode ci-dessus).
    \item Insère un saut de page.
    \item Page 3 : Rédige la structure suivante en appliquant les styles indiqués :
    \begin{itemize}
        \item \texttt{Titre 1} : 1. Introduction
        \item \texttt{Normal} : Deux lignes de texte factice.
        \item \texttt{Titre 1} : 2. Présentation du lycée
        \item \texttt{Titre 2} : 2.1 Historique
        \item \texttt{Normal} : Texte factice.
        \item \texttt{Titre 2} : 2.2 Effectifs
        \item \texttt{Titre 3} : 2.2.1 Élèves
        \item \texttt{Titre 3} : 2.2.2 Personnel
        \item \texttt{Titre 1} : 3. Activités menées
        \item \texttt{Titre 2} : 3.1 Maintenance
        \item \texttt{Titre 2} : 3.2 Réseau
        \item \texttt{Titre 1} : 4. Conclusion
    \end{itemize}
    \item Retourne en Page 2 (sur la table des matières). Clic droit > \textbf{Mettre à jour les champs} > \textbf{Mettre à jour toute la table}.
    \item Vérifie : les numéros de page correspondent-ils ? La hiérarchie (retraits) est-elle visible ?
    \item \textbf{Test de robustesse :} Ajoute un paragraphe long en page 3 pour décaler « 4. Conclusion » en page 4. Remets à jour la table (\texttt{F9}). Le numéro de page a-t-il changé ?
\end{enumerate}
\textbf{Observation attendue :} La table reflète fidèlement la structure et les pages. La modification de style (ex. mettre \texttt{Titre 1} en rouge) se répercute dans la table.
\end{experience}

\courssub{Illustration : Hiérarchie des styles et table des matières générée}

\begin{popfigure}
\begin{tikzpicture}[
    scale=0.9,
    every node/.style={font=\small},
    level 1/.style={sibling distance=5.5cm},
    level 2/.style={sibling distance=2.8cm},
    level 3/.style={sibling distance=1.6cm},
    level distance=1.6cm,
    edge from parent/.style={draw=popdark, thick, -{Stealth[length=2mm]}},
    titre1/.style={rectangle, rounded corners=3pt, draw=popblue, very thick, fill=popblueL, text width=4.5cm, align=center, minimum height=1.1cm},
    titre2/.style={rectangle, rounded corners=3pt, draw=poporange, thick, fill=poporangeL, text width=3.8cm, align=center, minimum height=0.9cm},
    titre3/.style={rectangle, rounded corners=3pt, draw=popgreen, thick, fill=popgreenL, text width=3.2cm, align=center, minimum height=0.8cm},
    tdm/.style={rectangle, draw=poppurple, very thick, fill=poppurpleL, text width=5cm, align=left, minimum height=6cm},
    arrow/.style={draw=popmuted, thick, -{Stealth[length=2mm]}, dashed}
]
% Arbre des styles (gauche)
\node [titre1, align=center] (root){Titre 1\\ \texttt{(Chapitre Principal)}\\\small\textit{ex: 1. Introduction}}
    child { node [titre2, align=center] (c1){Titre 2\\ \texttt{(Section)}\\\small\textit{ex: 2.1 Historique}}
        child { node [titre3, align=center] (c11){Titre 3\\ \texttt{(Sous-section)}\\\small\textit{ex: 2.1.1 Fondation}} }
        child { node [titre3, align=center] (c12){Titre 3\\ \texttt{(Sous-section)}\\\small\textit{ex: 2.1.2 Évolution}} }
    }
    child { node [titre2, align=center] (c2){Titre 2\\ \texttt{(Section)}\\\small\textit{ex: 2.2 Effectifs}}
        child { node [titre3, align=center] (c21){Titre 3\\ \texttt{(Sous-section)}\\\small\textit{ex: 2.2.1 Élèves}} }
        child { node [titre3, align=center] (c22){Titre 3\\ \texttt{(Sous-section)}\\\small\textit{ex: 2.2.2 Personnel}} }
    }
    child { node [titre1, align=center] (c3){Titre 1\\ \texttt{(Chapitre Principal)}\\\small\textit{ex: 3. Activités}} };

% Table des matières (droite)
\node [tdm, right=3.5cm of root, anchor=west, align=center] (tdmnode){
    \textbf{Table des matières}\\\vspace{2pt}
    1. Introduction ............ 3\\
    2. Présentation du lycée ............ 4\\
    \hspace{15pt}2.1 Historique ............ 4\\
    \hspace{30pt}2.1.1 Fondation ............ 4\\
    \hspace{30pt}2.1.2 Évolution ............ 5\\
    \hspace{15pt}2.2 Effectifs ............ 5\\
    \hspace{30pt}2.2.1 Élèves ............ 5\\
    \hspace{30pt}2.2.2 Personnel ............ 6\\
    3. Activités ............ 7\\
    4. Conclusion ............ 8
};

% Flèches de génération
\draw [arrow] (root.east) -- ++(1.2,0) |- (tdmnode.north west) node [midway, above, sloped, font=\tiny, text=popdark] {Génération auto};
\draw [arrow] (c1.east) -- ++(1.2,0) |- (tdmnode.west);
\draw [arrow] (c3.east) -- ++(1.2,0) |- (tdmnode.south west);

% Légendes styles
\node [below=0.3cm of root, font=\tiny\bfseries, text=popblue] {Style : Titre 1};
\node [below=0.3cm of c1, font=\tiny\bfseries, text=poporange] {Style : Titre 2};
\node [below=0.3cm of c11, font=\tiny\bfseries, text=popgreen] {Style : Titre 3};
\node [below=0.3cm of tdmnode, font=\small\bfseries, text=poppurple] {Table des matières générée (Références > Table des matières)};
\end{tikzpicture}
\legende{Relation bidirectionnelle : les styles de titres (gauche) structurent le document ; la table des matières (droite) en est la projection dynamique. Toute modification de titre ou de pagination impose une mise à jour (F9).}
\end{popfigure}

\courssub{Extension : Listes automatiques des tableaux et des illustrations}

\begin{savaistu}[Les listes d'objets : Tables des figures et des tableaux]
Dans les rapports techniques, mémoires de master ou thèses de doctorat (ex. à l'Université de Yaoundé I ou à l'ENSET), on exige souvent trois listes en début de document :
\begin{enumerate}
    \item \textbf{Table des matières} (basée sur styles Titre 1/2/3).
    \item \textbf{Table des figures / illustrations} (basée sur les \textbf{légendes} des images).
    \item \textbf{Table des tableaux} (basée sur les \textbf{légendes} des tableaux).
\end{enumerate}
Le principe est identique : tu insères une \textbf{légende} (Clic droit sur l'image/tableau > \textbf{Insérer une légende} > Catégorie « Figure » ou « Tableau ») ; puis onglet \textbf{Références} > \textbf{Insérer une table des illustrations} (Word) ou \textbf{Insertion > Table des matières et index > Table des illustrations} (LibreOffice). La mise à jour se fait aussi par \texttt{F9}.
\end{savaistu}

\begin{exoresolu}[Corriger une table des matières défaillante]
\textbf{Énoncé :} Un élève a généré sa table des matières. Elle n'affiche que : « 1. Introduction ............ 2 » alors que le document contient 5 chapitres (Titre 1) et plusieurs sections (Titre 2). Il a bien mis les titres en gras et taille 16 manuellement. Explique la cause et donne la procédure de correction complète.

\textbf{Solution détaillée :}
\begin{enumerate}
    \item \textbf{Cause :} Les titres sont formatés \emph{manuellement} (gras, taille, centrage) mais n'ont \emph{pas} le style \texttt{Titre 1} ni \texttt{Titre 2} appliqué. Le générateur de table des matières ne repère que les paragraphes portant le style \texttt{Titre 1}, \texttt{Titre 2}, etc.
    \item \textbf{Correction :}
    \begin{enumerate}
        \item Sélectionne le premier titre principal (ex. « 1. Introduction »).
        \item Applique le style \texttt{Titre 1} (Galerie styles ou \texttt{Ctrl+Alt+1}).
        \item Répète pour tous les chapitres principaux (\texttt{Titre 1}) et toutes les sections (\texttt{Titre 2}).
        \item \textit{Astuce :} Utilise le volet \textbf{Navigation} (Word) ou \textbf{Navigateur} (\texttt{F5} dans LibreOffice) pour voir quels paragraphes sont reconnus comme titres (ils apparaissent dans l'arborescence). Ceux qui manquent n'ont pas le bon style.
        \item Une fois tous les titres stylisés, place le curseur dans la table des matières existante.
        \item Clic droit > \textbf{Mettre à jour les champs} > \textbf{Mettre à jour toute la table} > \textbf{OK}.
    \end{enumerate}
    \item \textbf{Vérification :} La table affiche maintenant la structure complète avec les bons numéros de page.
\end{enumerate}
\end{exoresolu}

\courssub{Synthèse de la section}

\begin{aretenir}[Ce qu'il faut retenir sur les styles et la TdM]
\begin{itemize}
    \item Un \textbf{style} est un ensemble de formats enregistré sous un nom (\texttt{Titre 1}, \texttt{Normal}...). Il assure la cohérence et porte le sens structurel.
    \item \textbf{Toujours} utiliser \texttt{Titre 1}, \texttt{Titre 2}, \texttt{Titre 3} pour les titres (jamais mise en forme manuelle).
    \item La \textbf{table des matières automatique} (Onglet \textbf{Références}) se construit \emph{uniquement} à partir de ces styles.
    \item Après toute modification du document (texte, sauts de page, titres), \textbf{mettre à jour la table} (Clic droit > Mettre à jour les champs / \texttt{F9} > « Mettre à jour toute la table »).
    \item Extension : même principe pour \textbf{Table des figures} et \textbf{Table des tableaux} (via les légendes).
\end{itemize}
\end{aretenir}

\begin{exercice}[Entraînement : Rapport d'activité du club informatique]
Tu es secrétaire du club informatique de ton collège. Tu dois rédiger le rapport d'activité du premier trimestre (3 pages minimum).
\begin{enumerate}
    \item Crée le document. Page de garde : Titre club, « Rapport d'activités - Trimestre 1 », ton nom, date.
    \item Page 2 : Insère la table des matières automatique.
    \item Pages 3+ : Structure le rapport avec les styles :
    \begin{itemize}
        \item Titre 1 : 1. Bilan des effectifs
        \item Titre 2 : 1.1 Nouveaux inscrits ; 1.2 Anciens membres
        \item Titre 1 : 2. Séances réalisées
        \item Titre 2 : 2.1 Initiation Word (date, contenu) ; 2.2 Initiation Excel (date, contenu) ; 2.3 Maintenance salle info (date, actions)
        \item Titre 1 : 3. Projets du 2ème trimestre
        \item Titre 2 : 3.1 Concours de frappe ; 3.2 Atelier algorithmique
        \item Titre 1 : 4. Conclusion
    \end{itemize}
    \item Insère au moins une image (logo du club ou photo générique) avec une légende « Figure 1 - Logo du club ». Génère la \textbf{Table des illustrations} en page 2 (sous la table des matières).
    \item Modifie le style \texttt{Titre 1} : police \texttt{Calibri}, couleur \texttt{Bleu foncé}, taille 16, espacement avant 18 pt. Vérifie la mise à jour globale.
    \item Ajoute du texte factice pour décaler « 4. Conclusion » sur une nouvelle page. Mets à jour les deux tables (\texttt{F9}).
    \item Enregistre en \texttt{.odt} et exporte en \texttt{.pdf}.
\end{enumerate}
\end{exercice}

\coursec{Synthèse : produire un document complet}

Dans la section précédente, nous avons appris à appliquer des \cle{styles de titres} et à générer automatiquement une \cle{table des matières}. Tu possèdes maintenant, une par une, toutes les techniques du traitement de texte : mise en forme des caractères et des paragraphes, mise en page, insertion d'objets, table des matières. Mais savoir utiliser chaque outil séparément ne suffit pas : dans la vie courante (à l'école, dans une administration, dans une entreprise), on te demandera de \emph{produire un document complet}, du début à la fin. Cette dernière section t'apprend à \cle{organiser une démarche de production} : c'est exactement le savoir-faire exigé au BEPC.

\courssub{La démarche complète de production d'un document}

\begin{definition}[Produire un document]
\cle{Produire un document} avec un logiciel de traitement de texte, c'est suivre une démarche organisée en étapes successives : analyser la consigne, structurer le plan, saisir le texte, mettre en forme, insérer les objets, générer la table des matières, vérifier, puis enregistrer et exporter le document final.
\end{definition}

Pourquoi suivre des étapes dans cet ordre ? Imagine un élève qui commence par choisir les couleurs et les polices avant même d'avoir écrit son texte : il perdra du temps, car chaque modification du contenu risque de casser sa mise en forme. Le bon réflexe du professionnel est : \textbf{d'abord le fond (le contenu), ensuite la forme (la présentation), enfin la vérification}.

\begin{methode}[Les étapes de production d'un document]
\begin{enumerate}
\item \textbf{Analyser la consigne} : lire attentivement ce qui est demandé (type de document, nombre de pages, objets à insérer, destinataire).
\item \textbf{Structurer le plan} : écrire les grandes parties et les titres du document sur papier ou directement à l'écran.
\item \textbf{Saisir le texte} : taper le contenu sans se soucier encore de la présentation.
\item \textbf{Mettre en forme} : appliquer les styles de titres, la police, la taille, les alignements, les listes à puces.
\item \textbf{Insérer les objets} : images, tableaux, en-tête, pied de page avec numérotation.
\item \textbf{Générer la table des matières} : à partir des styles de titres (onglet \emph{Références}).
\item \textbf{Vérifier} : relire, corriger l'orthographe, contrôler avec l'\cle{aperçu avant impression}.
\item \textbf{Enregistrer et exporter} : sauvegarder le fichier, puis l'exporter en PDF pour la diffusion.
\end{enumerate}
\end{methode}

\begin{popfigure}
\begin{center}
\begin{tikzpicture}[
  node distance=0.55cm,
  etape/.style={rectangle, rounded corners=3pt, draw=popblue, fill=popblueL, thick,
    text width=4.6cm, align=center, minimum height=0.85cm, font=\small},
  fleche/.style={-{Stealth[length=2.5mm]}, thick, popdark}
]
\node[etape, align=center] (p1){\textbf{1. Plan}\\analyser et structurer};
\node[etape, below=of p1, align=center] (p2){\textbf{2. Saisie}\\taper le contenu};
\node[etape, below=of p2, align=center] (p3){\textbf{3. Mise en forme}\\styles, polices, alignements};
\node[etape, below=of p3, align=center] (p4){\textbf{4. Insertions}\\images, tableaux, en-tête};
\node[etape, below=of p4, align=center] (p5){\textbf{5. Table des matières}\\génération automatique};
\node[etape, below=of p5, draw=poporange, fill=poporangeL, align=center] (p6){\textbf{6. Vérification}\\orthographe, aperçu};
\node[etape, below=of p6, draw=popgreen, fill=popgreenL, align=center] (p7){\textbf{7. Export PDF}\\diffusion du document};
\draw[fleche] (p1) -- (p2);
\draw[fleche] (p2) -- (p3);
\draw[fleche] (p3) -- (p4);
\draw[fleche] (p4) -- (p5);
\draw[fleche] (p5) -- (p6);
\draw[fleche] (p6) -- (p7);
\draw[fleche, dashed, poporange] (p6.west) -- ++(-1.2,0) |- (p3.west)
  node[midway, left, font=\scriptsize\itshape, text=poporange, align=right]{erreur\\détectée};
\end{tikzpicture}
\end{center}
\legende{Organigramme de la démarche de production d'un document. Si la vérification révèle un problème, on remonte dans la démarche pour corriger, puis on revérifie.}
\end{popfigure}

Remarque la flèche en pointillés sur le schéma : la vérification n'est pas une fin en soi. Si l'aperçu révèle un problème (un titre mal stylé, une image qui déborde), on \emph{remonte} dans la démarche pour corriger, puis on revérifie. C'est le principe de l'\cle{amélioration par retour en arrière}.

\courssub{Exemple guidé : le cahier des charges d'un rapport}

En entreprise comme à l'école, on part souvent d'un \cle{cahier des charges}, c'est-à-dire la liste précise de ce que le document doit contenir.

\begin{exemplebox}[Cahier des charges d'un rapport de 5 pages]
Le principal du collège demande au club informatique un rapport d'activité de l'année. Voici le cahier des charges :
\begin{itemize}
\item \textbf{5 pages} au format A4, orientation portrait, marges de 2,5 cm ;
\item une \textbf{page de garde} avec le nom du collège, le titre et l'année scolaire ;
\item une \textbf{table des matières} automatique en page 2 ;
\item \textbf{2 images} (photo d'une séance du club, logo du collège) avec légendes ;
\item \textbf{1 tableau} récapitulatif des dépenses du club ;
\item un \textbf{en-tête} avec le nom du collège et un \textbf{pied de page} avec la numérotation des pages ;
\item police Times New Roman taille 12 pour le texte, titres en style \emph{Titre 1} et \emph{Titre 2} ;
\item diffusion finale en \textbf{PDF} par courriel au principal.
\end{itemize}
\end{exemplebox}

\begin{exoresolu}[Produire le rapport du club informatique]
Énoncé : en suivant le cahier des charges ci-dessus, décris la démarche complète de production du rapport, étape par étape.
\tcblower
\textbf{Solution détaillée.}
\begin{enumerate}
\item \textbf{Analyse} : le document est un rapport de 5 pages, destiné au principal, avec 2 images, 1 tableau, une table des matières et une sortie PDF.
\item \textbf{Plan} : page de garde ; table des matières ; 1. Présentation du club ; 2. Activités de l'année ; 3. Bilan financier ; 4. Perspectives.
\item \textbf{Saisie} : on tape tout le texte des quatre parties, sans mise en forme.
\item \textbf{Mise en forme} : on sélectionne chaque titre de partie et on lui applique le style \emph{Titre 1} ; le corps du texte est mis en Times New Roman 12, justifié.
\item \textbf{Insertions} : on insère les 2 images (menu \emph{Insertion} $>$ \emph{Images}) avec une légende chacune, le tableau des dépenses (3 colonnes : Poste, Montant prévu, Montant dépensé), l'en-tête «\ Collège bilingue de \ldots\ »\ et le pied de page avec le numéro de page.
\item \textbf{Table des matières} : on place le curseur en page 2, onglet \emph{Références} $>$ \emph{Table des matières} $>$ style automatique.
\item \textbf{Vérification} : on lance le correcteur orthographique, puis \emph{Fichier} $>$ \emph{Imprimer} pour voir l'aperçu : on contrôle que les 5 pages sont correctes, que les numéros de la table des matières correspondent aux pages réelles.
\item \textbf{Export} : \emph{Fichier} $>$ \emph{Enregistrer sous} (format .docx pour garder une version modifiable), puis \emph{Fichier} $>$ \emph{Exporter} $>$ \emph{Créer PDF} pour la version à envoyer.
\end{enumerate}
\end{exoresolu}

\courssub{Les vérifications finales avant diffusion}

Un document envoyé avec des fautes ou une présentation incohérente donne une mauvaise image de son auteur. Avant toute diffusion, effectue cette liste de contrôle :

\begin{itemize}
\item \textbf{Orthographe et grammaire} : lancer le correcteur (touche F7 dans Word) puis relire soi-même, car le correcteur ne comprend pas toujours le sens (il ne signalera pas «\ \emph{il son parti}\ »\ si la phrase est grammaticalement possible).
\item \textbf{Cohérence des styles} : tous les titres de même niveau doivent utiliser le même style ; on ne doit pas avoir un titre en gras 14 fait «\ à la main\ »\ à côté de titres en style \emph{Titre 1}.
\item \textbf{Numérotation des pages} : vérifier que les numéros apparaissent bien dans le pied de page et qu'ils se suivent.
\item \textbf{En-tête et pied de page} : présents sur toutes les pages (sauf éventuellement la page de garde).
\item \textbf{Table des matières à jour} : si du texte a été ajouté après sa génération, clic droit sur la table $>$ \emph{Mettre à jour les champs}.
\item \textbf{Aperçu avant impression} : dernière vérification visuelle, page par page.
\end{itemize}

\begin{attention}[Imprimer sans passer par l'aperçu avant impression]
L'erreur la plus fréquente (et la plus coûteuse !) est d'imprimer directement un document sans l'avoir examiné dans l'\cle{aperçu avant impression}. Résultat typique : un tableau coupé en deux pages, une image qui déborde dans la marge, une dernière page presque vide, 20 feuilles gaspillées au cybercafé. Passe \textbf{toujours} par \emph{Fichier} $>$ \emph{Imprimer} pour visualiser le rendu réel avant de cliquer sur le bouton \emph{Imprimer}.
\end{attention}

\courssub{Enregistrement et exportation en PDF}

Une fois le document terminé et vérifié, il faut le sauvegarder correctement :

\begin{enumerate}
\item \textbf{Enregistrer le fichier source} : \emph{Fichier} $>$ \emph{Enregistrer sous}, choisir un nom clair (par exemple \texttt{rapport\_club\_info\_2026.docx}) et un dossier bien rangé. Ce fichier reste modifiable.
\item \textbf{Exporter en PDF} : \emph{Fichier} $>$ \emph{Exporter} $>$ \emph{Créer un document PDF}. Le fichier PDF obtenu est la version \emph{finale}, destinée à être diffusée (imprimée, envoyée par courriel, déposée en administration).
\end{enumerate}

\begin{aretenir}[Pourquoi deux fichiers ?]
On conserve \textbf{deux versions} : le fichier de traitement de texte (.docx ou .odt), modifiable, pour les corrections futures ; et le fichier \cle{PDF}, non modifiable, pour la diffusion. On n'envoie jamais le fichier modifiable comme version officielle.
\end{aretenir}

\begin{savaistu}[Le PDF, un format universel]
Le format \cle{PDF} (\emph{Portable Document Format}), créé par la société Adobe en 1993, conserve la mise en page \emph{exactement à l'identique} sur tous les ordinateurs, téléphones et tablettes, quels que soient le système d'exploitation et les polices installées. \`A l'inverse, un fichier Word (.docx) peut s'afficher différemment d'une machine à l'autre : une police absente est remplacée, un tableau se décale, la pagination change. C'est pourquoi les administrations camerounaises (comme les sites de la DGI ou les concours de la fonction publique) exigent les dossiers en PDF.
\end{savaistu}

\courssub{Justifier sa démarche par écrit}

Au BEPC, on ne te demandera pas seulement de produire : on te demandera aussi d'\emph{expliquer} ce que tu as fait. C'est le savoir-faire «\ rédiger et justifier une démarche\ ». Une bonne justification répond à trois questions : \textbf{quoi} (ce qu'on a fait), \textbf{comment} (l'outil ou la commande utilisée), \textbf{pourquoi} (la raison du choix).

\begin{methode}[Rédiger la justification d'une démarche]
Pour chaque étape de production, rédiger une phrase en trois parties :
\begin{enumerate}
\item \textbf{Dire ce qu'on a fait} : «\ J'ai appliqué le style \emph{Titre 1} aux titres des parties\ldots\ »
\item \textbf{Dire comment} : «\ \ldots en sélectionnant le titre puis en cliquant sur le style dans le ruban \emph{Accueil}\ldots\ »
\item \textbf{Dire pourquoi} : «\ \ldots afin de générer automatiquement la table des matières et de garantir une présentation homogène. »
\end{enumerate}
Exemple complet : «\ J'ai exporté le document en PDF (\emph{Fichier} $>$ \emph{Exporter}) afin que le principal puisse le lire avec la mise en page exacte, quel que soit son ordinateur. »
\end{methode}

\begin{exercice}[\`A toi de jouer]
Le secrétaire du conseil d'administration te demande de produire le règlement intérieur de l'association des parents d'élèves : 4 pages, titres numérotés, table des matières, logo de l'association en page de garde, numérotation des pages, diffusion en PDF.
\begin{enumerate}
\item Rédige le cahier des charges du document.
\item Décris les étapes de production dans l'ordre.
\item Rédige la justification de deux choix de mise en forme (quoi, comment, pourquoi).
\end{enumerate}
\end{exercice}

\begin{corrige}[Exercice : règlement intérieur]
\textbf{1. Cahier des charges} : document de 4 pages A4 portrait ; page de garde avec le logo de l'association ; titres de chapitres en style \emph{Titre 1} et numérotés ; table des matières automatique ; numérotation des pages en pied de page ; export final en PDF.

\textbf{2. Étapes de production} : analyser la consigne ; rédiger le plan (chapitres : But de l'association, Membres, Cotisations, Réunions, Sanctions) ; saisir le texte ; appliquer les styles \emph{Titre 1} aux chapitres et mettre le corps en police 12 justifié ; insérer le logo en page de garde et la numérotation en pied de page ; générer la table des matières (\emph{Références} $>$ \emph{Table des matières}) ; corriger l'orthographe et vérifier avec l'aperçu avant impression ; enregistrer en .docx puis exporter en PDF.

\textbf{3. Justifications} (exemples) :
\begin{itemize}
\item «\ J'ai appliqué le style \emph{Titre 1} à chaque chapitre via le ruban \emph{Accueil}, afin de générer automatiquement la table des matières et d'obtenir des titres homogènes. »
\item «\ J'ai inséré la numérotation des pages via \emph{Insertion} $>$ \emph{Numéro de page}, afin que le lecteur puisse se repérer facilement et que la table des matières soit exploitable. »
\end{itemize}
\end{corrige}

\begin{aretenir}[L'essentiel de la synthèse]
\begin{itemize}
\item Produire un document suit une démarche ordonnée : \textbf{plan} $>$ \textbf{saisie} $>$ \textbf{mise en forme} $>$ \textbf{insertions} $>$ \textbf{table des matières} $>$ \textbf{vérification} $>$ \textbf{export PDF}.
\item Le fond avant la forme : on saisit d'abord, on décore ensuite.
\item On vérifie toujours avec l'aperçu avant impression avant d'imprimer ou de diffuser.
\item On garde le fichier modifiable et on diffuse le PDF, qui conserve la mise en page partout.
\item Justifier une démarche, c'est dire \textbf{ce qu'on a fait}, \textbf{comment} et \textbf{pourquoi}.
\end{itemize}
\end{aretenir}

\newpage

\coursec{Activité d'intégration}

\coursec{Synthèse : Produire un document complet}

\begin{aretenir}[Objectifs de la synthèse]
Cette activité d'intégration mobilise l'ensemble des compétences acquises dans l'unité « Traitement de texte » pour produire un document professionnel complet. À l'issue de ce travail, tu dois être capable de :
\begin{itemize}
    \item Structurer un document long à l'aide des \cle{styles de titres} (Titre 1, Titre 2) pour en automatiser la navigation et la table des matières.
    \item Maîtriser la \cle{mise en page} (marges, orientation, sauts de section, en-têtes/pieds de page différenciés, numérotation) et la \cle{mise en forme} (polices, alignement, interligne, puces/numérotation).
    \item Insérer et paramétrer des \cle{objets variés} : tableau de données trié et formaté, image avec habillage \og Carré \fg{}, sauts de page/section.
    \item Générer et mettre à jour une \cle{table des matières automatique} basée sur la hiérarchie des styles.
    \item Justifier par écrit les choix techniques effectués (accessibilité, lisibilité, conformité au cahier des charges), développant ainsi le savoir-être « rigueur » et « communication technique ».
\end{itemize}
\end{aretenir}

\begin{experience}[Situation professionnelle : Rapport annuel du Club « Génération Numérique »]
Le \textbf{Club Informatique « Génération Numérique »} du \textbf{Lycée Bilingue de Nkolbisson} (Yaoundé) termine son année scolaire 2025-2026. En tant que secrétaire(e) du club, tu es chargé(e), avec ton binôme, de rédiger le \textbf{rapport annuel d'activités} qui sera remis au Proviseur, aux partenaires (CAMTEL, MTN Foundation) et archivé au secrétariat de l'établissement.

Le document final doit respecter scrupuleusement le cahier des charges suivant, établi par le bureau du club sous la supervision de l'encadreur (M. \textsc{Fouda}) :
\end{experience}

\begin{exemplebox}[Cahier des charges détaillé]
\begin{enumerate}
    \item \textbf{Format général :} Document de \textbf{3 pages minimum} (hors page de garde), format A4, marges \textbf{Normales} (2,54 cm), police par défaut \textbf{Calibri 11 pt}, interligne \textbf{1,15}, alignement \textbf{Justifié}.
    \item \textbf{Page de garde (page 1, Section 1) :}
    \begin{itemize}
        \item Centrage vertical et horizontal.
        \item Logo du club (fichier \texttt{logo\_club.png}) inséré en haut, redimensionné à \textbf{3 cm de hauteur}, habillage \textbf{Devant le texte} pour centrage parfait.
        \item Titre principal : \og RAPPORT ANNUEL D'ACTIVITÉS 2025-2026 \fg{} — Style \textbf{Titre}, police \textbf{Arial 26 pt}, couleur \textbf{Bleu foncé (RVB 0, 51, 102)}, espacement avant/après 18 pt.
        \item Sous-titre : \og Club Informatique « Génération Numérique » — Lycée Bilingue de Nkolbisson \fg{} — Style \textbf{Sous-titre}, police \textbf{Calibri 14 pt}, italique.
        \item Infos bas de page : Année scolaire, Président (binôme), Secrétaire (toi), Encadreur (M. Fouda). Police Calibri 11 pt.
        \item \textbf{Aucun numéro de page} sur cette page.
    \end{itemize}
    \item \textbf{Sommaire automatique (page 2, début Section 2) :}
    \begin{itemize}
        \item Inséré via \texttt{Références > Table des matières > Automatique 1}.
        \item Doit lister les titres de niveau 1 et niveau 2.
        \item Titre \og SOMMAIRE \fg{} mis en forme manuelle (Arial 16, Gras, Centré) pour ne pas apparaître dans la table.
    \end{itemize}
    \item \textbf{Corps du rapport (à partir de la page 3, Section 2) :}
    \begin{itemize}
        \item \textbf{Chapitre 1 : Bilan des activités} (Style Titre 1).
        \begin{itemize}
            \item \textbf{1.1 Ateliers de formation} (Style Titre 2) : Liste à puces des 5 ateliers (ex. : \og Initiation Python \fg{}, \og Maintenance hardware \fg{}, \og Découverte Arduino \fg{}, \og Sécurité numérique \fg{}, \og PAO avec Canva \fg{}). Insérer \textbf{image illustrative} (\texttt{atelier.jpg}) à droite du paragraphe sur Python, habillage \textbf{Carré}, largeur 5 cm.
            \item \textbf{1.2 Projets menés} (Style Titre 2) : Description du projet \og Mise en réseau de la salle multimédia \fg{} (gaine, câblage RJ45, switch 24 ports TP-Link, adressage IP statique 192.168.10.0/24).
        \end{itemize}
        \item \textbf{Chapitre 2 : Gestion administrative et financière} (Style Titre 1).
        \begin{itemize}
            \item \textbf{2.1 Effectifs et bureau} (Style Titre 2) : \textbf{Tableau 6 lignes × 4 colonnes} (en-tête + 5 membres) : \textbf{Nom}, \textbf{Classe}, \textbf{Poste}, \textbf{Cotisation payée (FCFA)}. Données inventées cohérentes. Style \textbf{Moyen 1 - Accent 1}, centré, répétition ligne d'en-tête si coupure.
            \item \textbf{2.2 Budget prévisionnel vs réalisé} (Style Titre 2) : Paragraphe + second tableau simple (Recettes/Dépenses).
        \end{itemize}
        \item \textbf{Chapitre 3 : Perspectives 2026-2027} (Style Titre 1).
        \begin{itemize}
            \item \textbf{3.1 Demande de subvention} (Style Titre 2) : Lettre type au Délégué Régional MINESEC Centre.
            \item \textbf{3.2 Projets innovants} (Style Titre 2) : Liste numérotée (1. Club Robotique, 2. Compétition Inter-lycées, 3. Site web du club).
        \end{itemize}
    \end{itemize}
    \item \textbf{En-têtes / Pieds de page (Section 2 uniquement, à partir page 2) :}
    \begin{itemize}
        \item \textbf{En-tête (identique pair/impair) :} \og Lycée Bilingue de Nkolbisson — Club Informatique \fg{} (gauche, Calibri 9 pt, italique, trait de séparation bas).
        \item \textbf{Pied de page :} Numérotation \textbf{bas de page, centrée, format 1, 2, 3...}. \textbf{Numérotation commence à 1 sur la page du Sommaire (page 2 du fichier)}. Saut de section \og Page suivante \fg{} après page de garde pour dissocier les sections.
    \end{itemize}
\end{enumerate}
\end{exemplebox}

\begin{methode}[Démarche pas à pas — Consignes de réalisation]
Travaillant en binôme, réalisez les étapes dans l'ordre. Chaque étape validée par l'enseignant conditionne la suivante.
\begin{enumerate}
    \item \textbf{Préparation et mise en page de base :} Nouveau document. Configurer marges (Normales), police par défaut (Calibri 11, interligne 1,15, justifié). Insérer un \textbf{saut de section « Page suivante »} en fin de page de garde pour préparer la numérotation différenciée.
    \item \textbf{Page de garde (Section 1) :} Insérer logo, saisir titres, appliquer mise en forme spécifique (couleurs, tailles, centrage vertical via \texttt{Disposition > Mise en page > Onglet Disposition > Centrage vertical : Centré}). Vérifier absence de numéro de page.
    \item \textbf{Structure par les styles :} Sur pages suivantes, saisir titres des chapitres/sous-chapitres en appliquant \textbf{impérativement} les styles \texttt{Titre 1} et \texttt{Titre 2} (onglet \texttt{Accueil}). Modifier ces styles (clic droit > Modifier) pour uniformiser : Titre 1 (Arial 16 pt, Bleu foncé, Avant 24 pt / Après 6 pt, Niveau de plan 1), Titre 2 (Arial 13 pt, Bleu foncé, Gras, Avant 12 pt / Après 6 pt, Niveau de plan 2).
    \item \textbf{Rédaction et insertion d'objets :} Rédiger le contenu. Insérer tableau des membres (\texttt{Insertion > Tableau}), remplir, appliquer style, configurer répétition en-tête (\texttt{Disposition (Tableau) > Répéter les lignes d'en-tête}). Insérer \texttt{atelier.jpg}, habillage \texttt{Carré}, positionner à droite.
    \item \textbf{En-têtes, pieds de page et numérotation (Section 2) :} Aller en page 2. Désactiver \texttt{Lier au précédent} pour en-tête et pied de page. Configurer en-tête. Insérer numéro de page centré. Format de numéro : démarrer à \textbf{1}. Vérifier page de garde sans numéro.
    \item \textbf{Table des matières :} Curseur début page 2 (après titre SOMMAIRE). Insérer table automatique. Mettre à jour (\texttt{Mettre à jour toute la table}).
    \item \textbf{Relecture, justification et remise :} Vérifier conformité (3 pages min, styles, numérotation, objets). Rédiger un \textbf{paragraphe de justification (150 mots max)} en fin de document (après saut de page) expliquant : \og Pourquoi les styles plutôt que mise en forme manuelle ? Intérêt de l'habillage \texttt{Carré} ? Pourquoi un saut de section pour la numérotation ? \fg. Enregistrer sous \texttt{Rapport\_Club\_NomBinome.docx} et déposer sur clé USB / espace ENT.
\end{enumerate}
\end{methode}

\begin{definition}[Ressources à mobiliser — Raccourcis et outils clés]
\begin{itemize}
    \item \textbf{Raccourcis clavier :} \texttt{Ctrl+S} (Enregistrer), \texttt{Ctrl+Z} (Annuler), \texttt{Ctrl+A} (Tout sélectionner), \texttt{Ctrl+E} (Centrer), \texttt{Ctrl+J} (Justifier), \texttt{Ctrl+Entrée} (Saut de page).
    \item \textbf{Styles (Onglet Accueil) :} Galerie de styles, Volet \texttt{Styles} (\texttt{Alt+Ctrl+Maj+S}), Modification de style (\texttt{Format > Police / Paragraphe}).
    \item \textbf{Mise en page (Onglet Disposition) :} Marges, Sauts (Page, Section « Page suivante »), Orientation, Colonnes, Hyphénation.
    \item \textbf{Insertion (Onglet Insertion) :} Tableau, Images (Ce périphérique), En-tête/Pied de page, Numéro de page, Table des matières, Saut de page.
    \item \textbf{Outils de tableau (Onglets Création / Disposition contextuels) :} Styles de tableau, Fusionner/Séparer, Répéter les lignes d'en-tête, Tri, Formules (somme cotisations).
    \item \textbf{Format d'image (Onglet Format d'image contextuel) :} Habillage (En ligne, Carré, Derrière/Devant le texte), Recadrage, Corrections, Compression.
    \item \textbf{Références (Onglet Références) :} Table des matières, Mettre à jour la table, Marquer une entrée (index — hors programme mais utile).
\end{itemize}
\end{definition}

\begin{popfigure}
\begin{tikzpicture}[
    node distance=1.2cm and 1.5cm,
    box/.style={rectangle, draw=popblue, thick, fill=popblueL, rounded corners, minimum width=3cm, minimum height=1cm, align=center},
    arrow/.style={-Stealth, thick, popdark},
    sectionbox/.style={rectangle, draw=poporange, thick, fill=poporangeL, rounded corners, minimum width=4cm, minimum height=2.5cm, align=left, text width=3.8cm}
]
% Section 1
\node[sectionbox, align=center] (sec1){
    \textbf{Section 1 : Page de garde}\\
    \texttt{Lier au précédent : DÉSACTIVÉ}\\
    En-tête : Vide\\
    Pied de page : \textbf{Sans numéro}\\
    Centrage vertical : Oui
};
% Section 2
\node[sectionbox, right=of sec1, align=center] (sec2){
    \textbf{Section 2 : Corps (Sommaire + Rapport)}\\
    \texttt{Lier au précédent : DÉSACTIVÉ}\\
    En-tête : \textit{Lycée Bilingue...}\\
    Pied de page : \textbf{Numérotation à partir de 1}\\
    Centrage vertical : Non
};
% Arrow
\draw[arrow] (sec1.east) -- node[above, font=\small] {\texttt{Saut de section « Page suivante »}} (sec2.west);
% Footer detail
\node[box, below=of sec1.south, xshift=-1cm] (f1) {Page 1 : \textbf{Pas de numéro}};
\node[box, below=of sec2.south, xshift=1cm] (f2) {Page 2 : \textbf{1 (Sommaire)}};
\node[box, below=of f2] (f3) {Page 3 : \textbf{2 (Chap. 1)}};
\draw[arrow, dashed] (sec1.south) -- (f1.north);
\draw[arrow, dashed] (sec2.south) -- (f2.north);
\draw[arrow, dashed] (f2.south) -- (f3.north);
\end{tikzpicture}
\legende{Schéma de la structure en deux sections pour la gestion différenciée des pieds de page (numérotation).}
\end{popfigure}

\begin{exoresolu}[Réalisation pas à pas du Rapport Annuel]
\textbf{Énoncé :} Produire le rapport annuel du club « Génération Numérique » conforme au cahier des charges (page de garde sans numéro, sommaire auto, corps structuré par styles, tableau formaté, image avec habillage, en-tête/pied de page différenciés par saut de section, justification écrite).
\tcblower
\textbf{Étape 1 : Configuration globale et Saut de section}
\begin{itemize}
    \item Ouvrir Word (ou LibreOffice Writer). \texttt{Disposition > Marges > Normales}.
    \item \texttt{Accueil > Police par défaut : Calibri 11, Interligne 1,15, Alignement Justifié}.
    \item Se placer en bas de la future page de garde. \texttt{Disposition > Sauts > Page suivante (Section)}. \emph{Ce saut crée la Section 2, permettant un pied de page différent de la Section 1 (page de garde).}
\end{itemize}

\textbf{Étape 2 : Page de garde (Section 1)}
\begin{itemize}
    \item Centrage vertical : \texttt{Disposition > Lanceur boîte dialogue Mise en page > Onglet Disposition > Centrage vertical : Centré > OK}.
    \item \texttt{Insertion > Images > Ce périphérique > logo\_club.png}. Clic image > \texttt{Format d'image > Habillage > Devant le texte}. Redimensionner (poignée d'angle + \texttt{Maj}) à \textbf{3 cm hauteur}. Centrer (\texttt{Ctrl+E}).
    \item Saisir titre principal. Sélectionner. \texttt{Accueil > Styles > Titre}. Clic droit \texttt{Titre} > \texttt{Modifier} > Police Arial 26, Couleur RVB(0,51,102), \texttt{Format > Paragraphe} > Espacement Avant 18 pt / Après 18 pt. OK.
    \item Saisir sous-titre. Appliquer style \texttt{Sous-titre} (modifier : Calibri 14, Italique).
    \item Saisir infos bas de page (Année, Président, Secrétaire, Encadreur).
    \item \textbf{Vérification pied de page :} Double-cliquer bas de page > Onglet \texttt{Création (En-tête/Pied de page)} > Décocher \texttt{Lier au précédent} > \texttt{Pied de page > Supprimer le pied de page}.
\end{itemize}

\textbf{Étape 3 : Sommaire et Structure (Début Section 2 - Page 2)}
\begin{itemize}
    \item Se placer en haut page 2. Taper \og SOMMAIRE \fg. \textbf{Pas de style Titre 1}. Mise en forme manuelle : Arial 16, Gras, Centré (\texttt{Ctrl+E}).
    \item \texttt{Références > Table des matières > Automatique 1}. La TdM est vide (texte à venir).
    \item \texttt{Insertion > Saut de page} pour commencer Chapitre 1 page 3.
\end{itemize}

\textbf{Étape 4 : Personnalisation des styles Titre 1 et Titre 2 (Corps du document)}
\begin{itemize}
    \item \texttt{Accueil > Volet Styles (Alt+Ctrl+Maj+S)}. Survoler \texttt{Titre 1} > Flèche > \texttt{Modifier}.
    \item Police : Arial, 16 pt, Couleur Bleu foncé (Accent 1), Gras. \texttt{Format > Paragraphe} : Avant 24 pt, Après 6 pt, \textbf{Niveau de plan : Niveau 1}. OK.
    \item Idem pour \texttt{Titre 2} : Arial, 13 pt, Bleu foncé, Gras. Niveau de plan : Niveau 2. Avant 12 pt, Après 6 pt.
    \item \emph{Justification :} Le \cle{Niveau de plan} garantit que la TdM automatique reconnaîtra la hiérarchie (1, 1.1, 2, 2.1...).
\end{itemize}

\textbf{Étape 5 : Rédaction Chapitre 1 et Insertion Image}
\begin{itemize}
    \item Taper \og 1. Bilan des activités \fg (Style \texttt{Titre 1}). Taper \og 1.1 Ateliers de formation \fg (Style \texttt{Titre 2}).
    \item Rédiger paragraphe. Insérer \texttt{atelier.jpg} (\texttt{Insertion > Images}). Clic image > \texttt{Format d'image > Habillage > Carré}. Glisser l'image à droite du paragraphe Python. Redimensionner à 5 cm largeur. Le texte contourne l'image.
\end{itemize}

\textbf{Étape 6 : Tableau des membres (Chapitre 2.1)}
\begin{itemize}
    \item \texttt{Insertion > Tableau > Insérer un tableau > 4 colonnes, 6 lignes}.
    \item Remplir en-tête : \textbf{Nom | Classe | Poste | Cotisation (FCFA)}. Remplir 5 lignes données fictives (ex. : Ngono Marie, 3ème A, Présidente, 2000).
    \item Clic dans tableau > \texttt{Création (Tableau) > Styles > Moyen 1 - Accent 1} (bleu).
    \item \texttt{Disposition (Tableau) > Données > Trier} : Trier par \texttt{Poste} ou \texttt{Nom}.
    \item \textbf{Crucial :} Sélectionner ligne d'en-tête > \texttt{Disposition > Répéter les lignes d'en-tête}.
    \item Centrer tableau : \texttt{Disposition > Propriétés > Tableau > Alignement : Centré}.
\end{itemize}

\textbf{Étape 7 : En-têtes, Pieds de page et Numérotation (Section 2)}
\begin{itemize}
    \item Double-cliquer en-tête page 2 (Sommaire). Onglet \texttt{Création} : Cliquer \texttt{Lier au précédent} pour le \textbf{DÉSACTIVER} (bouton non surligné).
    \item Saisir : \og Lycée Bilingue de Nkolbisson — Club Informatique \fg. Police Calibri 9, Italique. \texttt{Disposition > Bordures > Bordure inférieure}.
    \item Double-cliquer pied de page page 2. \texttt{Lier au précédent} \textbf{DÉSACTIVÉ}.
    \item \texttt{Pied de page > Numéro de page > Bas de page > Simple 2 (centré)}.
    \item \texttt{Numéro de page > Format de numéro de page} : Démarrer à \textbf{1}. Format : 1, 2, 3...
    \item Fermer l'en-tête/pied de page.
    \item \textbf{Test :} Page 1 (Garde) : pied vide. Page 2 (Sommaire) : \og 1 \fg. Page 3 (Chap 1) : \og 2 \fg. \emph{Conforme.}
\end{itemize}

\textbf{Étape 8 : Mise à jour Table des Matières}
\begin{itemize}
    \item Clic dans la TdM (page 2). Onglet \texttt{Références} (ou clic droit) > \texttt{Mettre à jour la table} > \texttt{Mettre à jour toute la table} > OK.
    \item Vérifier : Numéros de page corrects (Chap 1 p.2, Chap 2 p.3...). Titres 1.1, 2.1 en retrait (niveau 2).
\end{itemize}

\textbf{Étape 9 : Justification écrite (Fin du document)}
\begin{itemize}
    \item Insérer \texttt{Saut de page} après Chapitre 3.
    \item Titre \og JUSTIFICATION DES CHOIX TECHNIQUES \fg (Style Titre 1).
    \item Rédiger (exemple) :
    \og \textbf{Styles vs Manuel :} Les styles Titre 1/2 permettent la génération auto de la TdM et une mise à jour instantanée si on modifie le texte (pagination dynamique). La mise en forme manuelle aurait nécessité de refaire la TdM à la main à chaque modification, source d'erreurs.
    \textbf{Habillage Carré :} Il intègre l'image dans le flux de lecture sans casser la colonne de texte, optimisant l'espace et la lisibilité (texte à gauche, image à droite), contrairement à l'habillage \og En ligne \fg{} qui crée un grand blanc vertical.
    \textbf{Saut de section :} Indispensable pour dissocier la Section 1 (page de garde sans numéro) de la Section 2 (corps numéroté à partir de 1). Sans lui, la numérotation serait continue ou l'en-tête de la page de garde identique aux autres. \fg
\end{itemize}

\textbf{Étape 10 : Finalisation et remise}
\begin{itemize}
    \item \texttt{Fichier > Enregistrer sous > Rapport\_Club\_Binome\_NomPrenom.docx}.
    \item \texttt{Fichier > Info > Vérifier les problèmes > Inspecter le document} (supprimer propriétés personnelles).
    \item Exporter en PDF pour archive : \texttt{Fichier > Exporter > Créer un document PDF/XPS}.
\end{itemize}
\end{exoresolu}

\begin{attention}[Erreurs fréquentes observées — Points de vigilance]
\begin{itemize}
    \item \textbf{Oubli du saut de section :} Numérotation présente sur page de garde. \rightarrow \texttt{Disposition > Sauts > Page suivante} \textbf{AVANT} de configurer le pied de page Section 2.
    \item \textbf{\texttt{Lier au précédent} resté actif :} Modification en-tête/pied page 2 modifie page 1. \rightarrow Bien cliquer sur le bouton \texttt{Lier au précédent} (onglet \texttt{Création}) pour le désactiver (bouton non surligné).
    \item \textbf{TdM vide ou incomplète :} Textes mis en forme manuellement (Gras + Taille) au lieu d'appliquer le Style \texttt{Titre 1}. \rightarrow Sélectionner le texte > Clic sur \texttt{Titre 1} dans le ruban \texttt{Accueil}.
    \item \textbf{Image « En ligne » qui pousse le texte :} Habillage par défaut. \rightarrow Changer pour \texttt{Carré} ou \texttt{Serré} via \texttt{Format d'image > Habillage}.
    \item \textbf{Tableau coupé en bas de page sans en-tête :} Oubli de \texttt{Répéter les lignes d'en-tête} (onglet \texttt{Disposition} contextuel Tableau).
    \item \textbf{Numérotation ne démarre pas à 1 :} Oubli de \texttt{Format de numéro de page > Démarrer à : 1} dans la Section 2.
\end{itemize}
\end{attention}

\begin{aretenir}[Compétences certifiantes BEPC — Unité Traitement de texte]
À l'issue de cette leçon et de cette activité, tu dois être capable, en autonomie, de :
\begin{enumerate}
    \item Créer et structurer un document multipage avec \textbf{styles de titres} (Titre 1, Titre 2).
    \item Configurer la \textbf{mise en page} (marges, orientation, sauts de page/section).
    \item Gérer \textbf{en-têtes/pieds de page différents} selon les sections (ex: pas de numéro page de garde).
    \item Insérer et mettre en forme un \textbf{tableau de données} (style, tri, répétition en-tête).
    \item Insérer une \textbf{image} et gérer son \textbf{habillage} (Carré, Serré, Devant/Derrière le texte).
    \item Générer et \textbf{mettre à jour une table des matières automatique}.
    \item Enregistrer en \texttt{.docx} et exporter en \texttt{.pdf}.
\end{enumerate}
\end{aretenir}

\begin{savaistu}[Le format .docx : une archive ZIP standardisée]
Savais-tu qu'un fichier \texttt{.docx} (Office Open XML) est en réalité une \textbf{archive ZIP} ? Si tu changes l'extension en \texttt{.zip} et l'ouvres, tu y trouves une structure de dossiers standardisée : \texttt{/word/document.xml} (contenu texte), \texttt{/word/styles.xml} (styles), \texttt{/word/media/} (images), \texttt{/docProps/} (métadonnées). C'est pour cela qu'on peut récupérer des images d'un document Word sans même l'ouvrir, et que la corruption d'un seul fichier XML peut rendre le document illisible. La norme ISO/IEC 29500 garantit l'interopérabilité entre Word, LibreOffice, OnlyOffice, Google Docs...
\end{savaistu}

\newpage

\coursec{Méthodes et exercices résolus}

\coursec{Utilisation d'un logiciel de traitement de texte}

\courssub{Rappels : interface et manipulation de base}

\begin{definition}[Interface utilisateur (Word / LibreOffice Writer)]
L'écran de travail d'un traitement de texte moderne (Microsoft Word, LibreOffice Writer) s'organise autour de quatre zones principales :
\begin{itemize}
  \item Le \cle{Ruban} (ou Barre d'outils) : regroupement des commandes par \textbf{onglets} (\textbf{Accueil}, \textbf{Insertion}, \textbf{Disposition}, \textbf{Références}, \textbf{Révision}, \textbf{Affichage}...).
  \item La \cle{Zone de saisie} : grande surface blanche (page virtuelle) où l'on tape le texte ; le \textbf{curseur} (barre clignotante) indique la position d'insertion.
  \item La \cle{Barre d'état} (bas de fenêtre) : affiche le numéro de page, le nombre de mots, la langue, le zoom, le mode d'affichage.
  \item La \cle{Barre d'accès rapide} (haut gauche) : raccourcis permanents (Enregistrer, Annuler, Rétablir).
\end{itemize}
\end{definition}

\begin{aretenir}[Manipulations de base indispensables]
\begin{center}
\begin{tabularx}{\linewidth}{|l|X|X|}\hline
\rowcolor{popblueL} \textbf{Action} & \textbf{Souris} & \textbf{Raccourcis clavier (Word FR / LibreOffice)} \\\hline
Sélectionner un mot & Double-clic sur le mot & --- \\\hline
Sélectionner un paragraphe & Triple-clic dans le paragraphe & --- \\\hline
Sélectionner tout le document & Ctrl + clic dans la marge gauche & \texttt{Ctrl+A} (universel) \\\hline
Copier / Couper / Coller & Clic droit $\rightarrow$ Copier/Couper/Coller & \texttt{Ctrl+C} / \texttt{Ctrl+X} / \texttt{Ctrl+V} \\\hline
Annuler / Rétablir & Boutons barre d'accès rapide & \texttt{Ctrl+Z} / \texttt{Ctrl+Y} \\\hline
Enregistrer & Icône disquette / Fichier $\rightarrow$ Enregistrer & \texttt{Ctrl+S} \\\hline
Imprimer / Aperçu & Fichier $\rightarrow$ Imprimer & \texttt{Ctrl+P} \\\hline
Correction orthographique & Onglet Révision $\rightarrow$ Orthographe & \texttt{F7} \\\hline
\textbf{Gras} & Onglet Accueil $\rightarrow$ \textbf{G} & \texttt{Ctrl+G} (Word FR) / \texttt{Ctrl+B} (LibreOffice) \\\hline
\textbf{Italique} & Onglet Accueil $\rightarrow$ \textbf{I} & \texttt{Ctrl+I} (universel) \\\hline
\textbf{Souligné} & Onglet Accueil $\rightarrow$ \textbf{S} & \texttt{Ctrl+U} (universel) \\\hline
\textbf{Justifier} & Onglet Accueil $\rightarrow$ Icône justifier & \texttt{Ctrl+J} (universel) \\\hline
\end{tabularx}
\end{center}
\textbf{Attention :} Dans un établissement camerounais, les postes tournent souvent sous \textbf{LibreOffice} (gratuit, open source). Les raccourcis \texttt{Ctrl+B}, \texttt{Ctrl+I}, \texttt{Ctrl+U} sont universels ; \texttt{Ctrl+G} ne fonctionne \emph{que} dans Word version française. Privilégie les raccourcis universels.
\end{aretenir}

\courssub{Mise en forme des caractères et des paragraphes}

\begin{methode}[Bien appliquer une mise en forme (Caractères \& Paragraphes)]
La règle d'or : \textbf{Sélectionner d'abord, formater ensuite}.
\begin{enumerate}
  \item \cle{Sélectionner} la portion de texte concernée (mot, phrase, plusieurs paragraphes). Sans sélection, la mise en forme ne s'applique qu'au prochain caractère tapé.
  \item Dans l'onglet \textbf{Accueil}, groupe \textbf{Police} : choisir \textbf{Police} (ex. \texttt{Arial}, \texttt{Times New Roman}), \textbf{Taille} (ex. 12 pt), \textbf{Couleur}, \textbf{Mise en majuscules} (Aa $\rightarrow$ MAJUSCULES).
  \item Toujours dans \textbf{Accueil}, groupe \textbf{Police} : activer \textbf{Gras} (\texttt{Ctrl+B}), \textbf{Italique} (\texttt{Ctrl+I}), \textbf{Souligné} (\texttt{Ctrl+U}), \textbf{Barré}, \textbf{Exposant/Indice} ($x^2$, $H_2O$).
  \item Pour le \cle{paragraphe} (alignement, interligne, retraits) : cliquer \emph{n'importe où} dans le paragraphe (pas besoin de tout sélectionner). Onglet \textbf{Accueil}, groupe \textbf{Paragraphe} :
  \begin{itemize}
    \item \textbf{Alignement} : Gauche (défaut), Centré, Droite, \textbf{Justifié} (\texttt{Ctrl+J}) — aligné des deux côtés.
    \item \textbf{Interligne} : 1,0 (simple), 1,15 (défaut Word), 1,5, Double.
    \item \textbf{Retraits} : Augmenter/Diminuer le retrait ; ou clic droit $\rightarrow$ \textbf{Paragraphe} $\rightarrow$ \textbf{Retrait : Première ligne} (pour l'alinéa classique).
  \end{itemize}
  \item Vérifier visuellement, puis \textbf{Enregistrer} (\texttt{Ctrl+S}).
\end{enumerate}
\end{methode}

\begin{exemplebox}[Raccourcis de sélection rapides]
\begin{itemize}
  \item \texttt{Clic + Glisser} : sélection continue.
  \item \texttt{Ctrl + Clic} : sélectionner une phrase.
  \item \texttt{Shift + Flèches} : étendre la sélection caractère par caractère.
  \item \texttt{Ctrl+Shift+Flèches} : étendre mot par mot / paragraphe par paragraphe.
\end{itemize}
\end{exemplebox}

\begin{exoresolu}[Mise en forme d'une affiche scolaire]
Le club informatique du collège rédige une annonce : « Réunion du club informatique, vendredi à 14 h, salle multimédia ». On demande de : (a) mettre « Réunion du club informatique » en gras, taille 16, centré ; (b) mettre « vendredi à 14 h » en italique et en rouge ; (c) justifier le paragraphe. Décris les manipulations.
\tcblower
\textbf{Solution détaillée.}
\begin{enumerate}
  \item \textbf{Partie (a).} On sélectionne le texte « Réunion du club informatique » (cliquer-glisser ou double-clic mot à mot + Shift). Onglet \textbf{Accueil} $\rightarrow$ \textbf{Gras} (\texttt{Ctrl+B}), liste \textbf{Taille} $\rightarrow$ \textbf{16}, puis icône \textbf{Centrer} (groupe Paragraphe). Le titre apparaît en gras, plus grand et centré.
  \item \textbf{Partie (b).} On sélectionne « vendredi à 14 h », \textbf{Italique} (\texttt{Ctrl+I}), puis flèche du bouton \textbf{Couleur de police} (A avec barre colorée) $\rightarrow$ \textbf{Rouge}.
  \item \textbf{Partie (c).} On clique n'importe où dans le paragraphe (inutile de tout sélectionner pour l'alignement), icône \textbf{Justifier} (\texttt{Ctrl+J}). Le texte est aligné à gauche et à droite.
  \item On enregistre avec \texttt{Ctrl+S}.
\end{enumerate}
\textbf{Conclusion :} La règle d'or est respectée : sélectionner d'abord, formater ensuite ; l'alignement s'applique au paragraphe entier sans sélection complète.
\end{exoresolu}

\courssub{Mise en page du document}

\begin{methode}[Régler la mise en page (Document entier)]
La \cle{mise en page} concerne la feuille physique. Elle se règle dans l'onglet \textbf{Disposition} (Word) ou \textbf{Format} $\rightarrow$ \textbf{Page} (LibreOffice) :
\begin{enumerate}
  \item \textbf{Format papier :} \textbf{Taille} $\rightarrow$ \textbf{A4} (21 $\times$ 29,7 cm), standard au Cameroun.
  \item \textbf{Orientation :} \textbf{Portrait} (vertical, lettres, rapports) ou \textbf{Paysage} (horizontal, tableaux larges, graphiques).
  \item \textbf{Marges :} \textbf{Marges} $\rightarrow$ \textbf{Marges personnalisées} (bas de liste). Saisir \textbf{Haut, Bas, Gauche, Droite} (ex. 2 cm ou 2,5 cm « Normales »). Onglet \textbf{Mise en page} de la boîte : \textbf{Gouttière} (marge supplémentaire pour reliure) si nécessaire.
  \item \textbf{Colonnes :} \textbf{Colonnes} $\rightarrow$ Deux/Trois (style journal).
  \item \textbf{Sauts de page :} \textbf{Sauts} $\rightarrow$ \textbf{Page} (ou \texttt{Ctrl+Entrée}) pour forcer le début d'une nouvelle page (ex. début de chapitre).
  \item \textbf{En-tête / Pied de page :} Onglet \textbf{Insertion} $\rightarrow$ \textbf{En-tête} / \textbf{Pied de page} $\rightarrow$ \textbf{Modifier l'en-tête/pied}. Y insérer : nom établissement, titre document, \textbf{Numéro de page} (Insertion $\rightarrow$ Numéro de page $\rightarrow$ Bas de page / Haut de page $\rightarrow$ style Centré).
  \item \textbf{Vérification finale :} \textbf{Aperçu avant impression} (\texttt{Ctrl+P}) : contrôler marges, numérotation, en-têtes, sauts de page.
\end{enumerate}
\end{methode}

\begin{exoresolu}[Préparer l'impression d'un exposé]
Un élève de 3ème doit imprimer un exposé de 6 pages. Exigences : format A4, orientation portrait, marges 2 cm, numérotation pages en bas au centre, nom du collège en haut de chaque page. Décris la démarche.
\tcblower
\textbf{Solution détaillée.}
\begin{enumerate}
  \item \textbf{Format/Orientation :} Onglet \textbf{Disposition} $\rightarrow$ \textbf{Taille} $\rightarrow$ \textbf{A4} ; \textbf{Orientation} $\rightarrow$ \textbf{Portrait}.
  \item \textbf{Marges :} \textbf{Disposition} $\rightarrow$ \textbf{Marges} $\rightarrow$ \textbf{Marges personnalisées} ; saisir 2 cm partout $\rightarrow$ \textbf{OK}.
  \item \textbf{En-tête :} Onglet \textbf{Insertion} $\rightarrow$ \textbf{En-tête} $\rightarrow$ \textbf{Modifier l'en-tête} ; taper « Collège Bilingue de Nkol-Eton » ; fermer (bouton \textbf{Fermer l'en-tête et le pied de page} ou double-clic dans le corps).
  \item \textbf{Pied de page / Numérotation :} \textbf{Insertion} $\rightarrow$ \textbf{Numéro de page} $\rightarrow$ \textbf{Bas de page} $\rightarrow$ \textbf{Style centré (Simple 2)}. Les pages 1 à 6 apparaissent automatiquement.
  \item \textbf{Vérification :} \texttt{Ctrl+P} (Aperçu). Contrôler chaque page : en-tête présent, numéro centré en bas, texte dans marges. Imprimer.
\end{enumerate}
\textbf{Conclusion :} La mise en page se règle une fois pour tout le document ; l'aperçu évite le gaspillage de papier.
\end{exoresolu}

\courssub{Insertion d'objets dans le document}

\begin{methode}[Insérer et gérer les objets (Images, Tableaux, Listes, Symboles)]
Toutes les insertions démarrent depuis l'onglet \textbf{Insertion} (ou barre latérale LibreOffice).
\begin{enumerate}
  \item \textbf{Positionner le curseur} exactement à l'endroit d'insertion.
  \item \textbf{Image (Photo, Logo, Schéma) :} \textbf{Images} $\rightarrow$ \textbf{Ce périphérique} (fichier local) $\rightarrow$ choisir fichier $\rightarrow$ \textbf{Insérer}.
  \begin{itemize}
    \item \textbf{Redimensionner :} Tirer une \textbf{poignée d'angle} (coin) pour garder les proportions (ratio verrouillé par défaut).
    \item \textbf{Habillage (Position relative au texte) :} Clic droit sur l'image $\rightarrow$ \textbf{Habillage} (Word) / \textbf{Ancrage} puis \textbf{Habillage} (LibreOffice) :
    \begin{itemize}
      \item \emph{En ligne avec le texte} (défaut) : l'image se comporte comme un gros caractère.
      \item \emph{Carré / Contour / Devant / Derrière le texte} : l'image flotte, le texte contourne.
    \end{itemize}
  \end{itemize}
  \item \textbf{Tableau :} \textbf{Tableau} $\rightarrow$ glisser dans la grille (lignes $\times$ colonnes) ou \textbf{Insérer un tableau} (saisir nb lignes/colonnes).
  \begin{itemize}
    \item Remplir les cellules (Tabulation = cellule suivante).
    \item \textbf{Outils de tableau} (onglets \textbf{Conception} / \textbf{Disposition}) : fusionner/séparer cellules, bordures, ombrage, alignement dans cellules.
  \end{itemize}
  \item \textbf{Listes :} Sélectionner les lignes $\rightarrow$ \textbf{Puces} (liste non ordonnée) ou \textbf{Numérotation} (liste ordonnée 1., 2. / a), b) / i, ii). Niveaux : \texttt{Tab} (sous-niveau), \texttt{Shift+Tab} (remonter).
  \item \textbf{Symboles / Caractères spéciaux :} \textbf{Symbole} $\rightarrow$ \textbf{Autres symboles} (©, ®, ™, œ, æ, ∑, √, flèches...).
  \item \textbf{Formules mathématiques :} \textbf{Équation} $\rightarrow$ \textbf{Insérer une nouvelle équation} (éditeur d'équations natif).
\end{enumerate}
\end{methode}

\begin{exoresolu}[Créer l'emploi du temps de la classe]
La censeure veut publier l'emploi du temps de la 3ème A : 5 jours (Lundi–Vendredi) et 4 plages horaires par jour. (a) Dimension du tableau ? (b) Insertion ? (c) Ajouter logo au-dessus ?
\tcblower
\textbf{Solution détaillée.}
\begin{enumerate}
  \item \textbf{(a) Dimension :} 1 colonne « Heures » + 5 colonnes jours = \textbf{6 colonnes}. 1 ligne d'en-tête + 4 lignes plages = \textbf{5 lignes}. Tableau \textbf{5 lignes $\times$ 6 colonnes}.
  \item \textbf{(b) Insertion :} Curseur à l'emplacement $\rightarrow$ \textbf{Insertion} $\rightarrow$ \textbf{Tableau} $\rightarrow$ sélectionner 6$\times$5. Remplir : Ligne 1 = Heures, Lundi... Vendredi. Lignes 2–5 = matières. Mettre ligne 1 en \textbf{Gras} + \textbf{Centré} (sélectionner ligne $\rightarrow$ Accueil).
  \item \textbf{(c) Logo :} Curseur \emph{avant} le tableau (Entrée pour faire de la place) $\rightarrow$ \textbf{Insertion} $\rightarrow$ \textbf{Images} $\rightarrow$ fichier logo $\rightarrow$ \textbf{Insérer}. Redimensionner par poignée d'angle $\rightarrow$ \textbf{Centrer} (paragraphe).
\end{enumerate}
\textbf{Conclusion :} Bien dimensionner le tableau dès le départ (5$\times$6) évite suppressions/ajouts de lignes fastidieux.
\end{exoresolu}

\courssub{Styles de titres et table des matières automatique}

\begin{methode}[Appliquer les styles de titres (Structure du document)]
Les \cle{styles} (Titre 1, Titre 2, Titre 3...) sont la \textbf{seule} façon pour le logiciel de comprendre la structure hiérarchique.
\begin{enumerate}
  \item Placer le curseur dans la ligne du titre (pas besoin de sélectionner).
  \item Onglet \textbf{Accueil}, galerie \textbf{Styles} : cliquer \textbf{Titre 1} (chapitre principal), \textbf{Titre 2} (section), \textbf{Titre 3} (sous-section).
  \item Répéter pour tous les titres en respectant la hiérarchie : \textbf{Jamais de Titre 3 sans Titre 2 parent}.
  \item \textbf{Modifier l'apparence globale :} Clic droit sur le style dans la galerie $\rightarrow$ \textbf{Modifier} $\rightarrow$ changer police, taille, couleur, espacement $\rightarrow$ \textbf{OK}. \emph{Tous} les titres de ce style se mettent à jour instantanément.
\end{enumerate}
\end{methode}

\begin{attention}[Piège classique : « Faux titres »]
Mettre un texte en \textbf{Gras, Taille 16, Centré} « à la main » (via groupe Police/Paragraphe) ne crée \textbf{PAS} un titre structurel.
\begin{itemize}
  \item Le logiciel ne le voit pas comme un titre.
  \item Il n'apparaîtra \textbf{jamais} dans la Table des matières automatique.
  \item Il ne sera pas navigable dans le \textbf{Volet de navigation} (Affichage $\rightarrow$ Volet de navigation).
\end{itemize}
\textbf{Règle :} Un titre = un Style (Titre 1, 2 ou 3).
\end{attention}

\begin{exoresolu}[Structurer un rapport de stage]
Rapport contenant : « Introduction », « Chapitre 1 : Présentation de l'entreprise », « 1.1 Historique », « 1.2 Organigramme », « Chapitre 2 : Activités réalisées », « Conclusion ». Indiquer le style pour chaque titre.
\tcblower
\textbf{Solution détaillée.}
La hiérarchie a deux niveaux : Parties principales (Niveau 1) et Sous-parties (Niveau 2).
\begin{center}
\begin{tabularx}{\linewidth}{|l|X|}\hline
\rowcolor{popblueL} \textbf{Titre du document} & \textbf{Style à appliquer} \\\hline
Introduction & Titre 1 \\\hline
Chapitre 1 : Présentation de l'entreprise & Titre 1 \\\hline
1.1 Historique & Titre 2 \\\hline
1.2 Organigramme & Titre 2 \\\hline
Chapitre 2 : Activités réalisées & Titre 1 \\\hline
Conclusion & Titre 1 \\\hline
\end{tabularx}
\end{center}
\textbf{Justification :} « Introduction », Chapitres, « Conclusion » = même niveau (Titre 1). « 1.1 » et « 1.2 » dépendent du Chapitre 1 = Titre 2.
\textbf{Conclusion :} Document structuré, prêt pour table des matières automatique.
\end{exoresolu}

\begin{methode}[Générer et mettre à jour la table des matières (TDMA)]
\begin{enumerate}
  \item \textbf{Prérequis obligatoire :} Tous les titres utilisent les styles Titre 1, 2, 3.
  \item Placer le curseur à l'emplacement de la table (généralement page 2, après la garde).
  \item Onglet \textbf{Références} $\rightarrow$ \textbf{Table des matières} $\rightarrow$ \textbf{Table automatique 1} (ou 2, ou « Table des matières personnalisée... » pour choisir les niveaux affichés).
  \item La table s'insère avec titres et numéros de page.
  \item \textbf{Mise à jour impérative :} Après \emph{toute} modification (ajout texte, image, saut de page) : \textbf{Clic droit sur la table} $\rightarrow$ \textbf{Mettre à jour les champs} $\rightarrow$ cocher \textbf{Mettre à jour toute la table} $\rightarrow$ \textbf{OK}.
\end{enumerate}
\textbf{Réflexe :} Générer la table \emph{en dernier}, quand le contenu est stable ; la mettre à jour \emph{juste avant} d'imprimer/exporter PDF.
\end{methode}

\begin{exoresolu}[Table des matières du rapport de stage]
Le rapport (12 pages) a une table tapée à la main. Après ajout d'un paragraphe, tous les numéros sont faux. (a) Erreur ? (b) Démarche correcte ?
\tcblower
\textbf{Solution détaillée.}
\begin{enumerate}
  \item \textbf{(a) Erreur :} Table \emph{manuelle}. Elle ne se met pas à jour ; tout décalage de pagination la rend obsolète.
  \item \textbf{(b) Démarche correcte :}
  \begin{enumerate}
    \item Vérifier styles Titre 1 / Titre 2 (fait précédemment).
    \item Curseur au début $\rightarrow$ \textbf{Références} $\rightarrow$ \textbf{Table des matières} $\rightarrow$ \textbf{Table automatique 1}.
    \item Après ajout du paragraphe : Clic droit sur table $\rightarrow$ \textbf{Mettre à jour les champs} $\rightarrow$ \textbf{Mettre à jour toute la table}.
  \end{enumerate}
\end{enumerate}
\textbf{Conclusion :} La TDMA fondée sur les styles garantit l'exactitude permanente sans retouche manuelle.
\end{exoresolu}

\courssub{Synthèse : Produire un document complet}

\begin{methode}[Conduire un projet de document de A à Z (Ordre logique)]
Pour tout document structuré (exposé, rapport, magazine, lettre officielle) :
\begin{enumerate}
  \item \textbf{Saisir} le texte brut (contenu) sans mise en forme.
  \item \textbf{Structurer} avec les \textbf{Styles de titres} (Titre 1, Titre 2).
  \item \textbf{Mettre en forme} caractères/paragraphes (Gras, Italique, Alignements, Listes, Polices).
  \item \textbf{Insérer} les objets (Images, Tableaux, Symboles, Équations).
  \item \textbf{Mettre en page} (Marges, Orientation, En-tête/Pied de page, Numérotation, Sauts de page).
  \item \textbf{Générer} la \textbf{Table des matières automatique}.
  \item \textbf{Finaliser :} Correction orthographique (\texttt{F7}), \textbf{Aperçu avant impression} (\texttt{Ctrl+P}), \textbf{Enregistrer} (\texttt{Ctrl+S}), Imprimer / Exporter en PDF.
\end{enumerate}
\textbf{Règle d'or :} \texttt{Ctrl+S} \emph{régulièrement} (toutes les 5–10 min) pour éviter la perte de données (coupure courant fréquente).
\end{methode}

\begin{exoresolu}[Organiser la production d'un magazine scolaire]
Club presse : magazine 8 pages (Garde, TDM, 3 articles+photos, Tableau résultats). Ranger dans l'ordre : Générer TDM ; Saisir textes ; Insérer photos/tableau ; Appliquer styles ; Régler marges/numérotation ; Relire/Imprimer.
\tcblower
\textbf{Solution détaillée.}
Ordre correct (Méthode de production) :
\begin{enumerate}
  \item \textbf{Saisir les textes} (matière première).
  \item \textbf{Appliquer les styles de titres} (structure $\rightarrow$ condition TDM).
  \item \textbf{Insérer photos et tableau} (intégration dans le flux).
  \item \textbf{Régler marges et numérotation} (mise en page globale, maintenant que le volume est connu).
  \item \textbf{Générer la table des matières} (numéros de page désormais définitifs).
  \item \textbf{Relire et imprimer} (\texttt{F7}, \texttt{Ctrl+P}, \texttt{Ctrl+S}).
\end{enumerate}
\textbf{Justification :} TDM avant la fin = mises à jour incessantes. Mise en page avant saisie = perturbations par le contenu.
\textbf{Conclusion :} Saisir $\rightarrow$ Structurer $\rightarrow$ Insérer $\rightarrow$ Mettre en page $\rightarrow$ TDM $\rightarrow$ Relire = efficacité maximale.
\end{exoresolu}

\begin{experience}[TP Guidé : « Mon rapport de visite au Musée National »]
\textbf{Objectif :} Produire un document complet de 3 pages appliquant toutes les compétences de la leçon.
\textbf{Logiciel :} LibreOffice Writer (ou Word).
\textbf{Consignes pas à pas :}
\begin{enumerate}
  \item \textbf{Nouveau document} $\rightarrow$ Enregistrer sous \texttt{Rapport\_Visite\_Musee\_PrenomNom.odt}.
  \item \textbf{Saisie :} Taper le texte brut (fourni par l'enseignant) : Titre principal, Introduction, 2 parties avec 2 sous-parties chacune, Conclusion. \textbf{Ne pas formater}.
  \item \textbf{Styles :} Appliquer \textbf{Titre 1} au titre principal, aux 2 parties et à la Conclusion. Appliquer \textbf{Titre 2} aux 4 sous-parties. Vérifier dans le \textbf{Volet de navigation} (F5).
  \item \textbf{Mise en forme :} Titre principal : Taille 18, Centré, Bleu foncé. Corps : Police \texttt{Calibri} 11, Justifié, Interligne 1,15. Mots-clés en \textbf{Gras}. Citations en \emph{Italique}.
  \item \textbf{Insertion :} Insérer une \textbf{image} (photo musée) sous l'introduction $\rightarrow$ Habillage « Carré », Centré. Insérer un \textbf{tableau} 3$\times$3 (Horaires, Tarifs, Accès) à la fin $\rightarrow$ Ligne 1 en gras + fond gris clair.
  \item \textbf{Mise en page :} A4, Portrait, Marges 2 cm. En-tête : « Collège X — 3ème A — Année 2026 ». Pied de page : Numéro de page centré.
  \item \textbf{Table des matières :} Page 2 (après garde), \textbf{Références} $\rightarrow$ Table automatique. Mettre à jour.
  \item \textbf{Finalisation :} \texttt{F7} (Correction), \texttt{Ctrl+P} (Aperçu : vérifier garde, TDM page 2, tableau page 3, numérotation continue), \texttt{Ctrl+S}.
  \item \textbf{Remise :} Fichier \texttt{.odt} (ou \texttt{.docx}) sur clé USB / plateforme.
\end{enumerate}
\textbf{Critères d'évaluation (sur 20) :} Styles/TDM (5), Mise en forme (4), Insertion/Image/Tableau (4), Mise en page/En-tête/Pied (4), Orthographe/Propreté (3).
\end{experience}

\begin{savaistu}[Le format OpenDocument (.odt) : standard camerounais]
Le format \textbf{ODT} (OpenDocument Text) est un standard international \textbf{ISO/IEC 26300}, natif de \textbf{LibreOffice} et \textbf{OpenOffice}. Contrairement au \texttt{.docx} (Microsoft, propriétaire bien que documenté), l'\texttt{.odt} est \textbf{entièrement ouvert}, sans redevance, et garanti lisible sur le long terme. L'administration camerounaise (MINESEC, MINFOPRA) recommande l'usage des formats ouverts pour la pérennité des archives numériques. Enregistrer en \texttt{.odt} assure que ton document sera lisible dans 20 ans, même si Microsoft Word n'existe plus !
\end{savaistu}

\begin{attention}[Les 5 erreurs qui coûtent des points au BEPC]
\begin{enumerate}
  \item \textbf{Formater sans sélectionner.} Cliquer « Gras » sans sélection $\rightarrow$ rien ne change (ou ne s'applique qu'au texte suivant). \textbf{Réflexe : Sélectionner $\rightarrow$ Formater.}
  \item \textbf{Confondre Mise en forme et Mise en page.} Mise en forme = Caractères/Paragraphes (Gras, Taille, Alignement). Mise en page = Document entier (Marges, Orientation, En-tête, Numérotation). Ne jamais mélanger dans une réponse d'examen.
  \item \textbf{Faire de « faux titres » à la main.} Gras + Taille 16 $\neq$ Titre structurel. Sans style \textbf{Titre 1} : absent de la TDM, absent du Volet de navigation. \textbf{Erreur n°1 au BEPC.}
  \item \textbf{Taper la table des matières à la main.} Devient fausse au moindre ajout. \textbf{Obligatoire :} \textbf{Références $\rightarrow$ Table des matières $\rightarrow$ Automatique} + \textbf{Mise à jour} avant impression.
  \item \textbf{Imprimer sans aperçu et sans enregistrer.} Oublier \texttt{Ctrl+P} = gaspillage papier/encre. Oublier \texttt{Ctrl+S} = perte totale en cas de coupure (fréquente au Cameroun). \textbf{Cite toujours ces deux réflexes en examen.}
\end{enumerate}
\end{attention}

\begin{exercice}[Entraînement BEPC — Sujet type]
On veut réaliser une lettre de motivation pour un stage au cybercafé « NetPlus » à Yaoundé.
\begin{enumerate}
  \item Cite les \textbf{3 onglets} du ruban que tu utiliseras pour : (a) mettre ton adresse en haut à droite, (b) écrire « Objet : Demande de stage » en gras et centré, (c) numéroter les pages en bas au centre.
  \item Explique pourquoi il faut appliquer le style \textbf{Titre 1} à « Objet : Demande de stage » si tu veux qu'il apparaisse dans une table des matières.
  \item Tu insères le logo « NetPlus.png ». Le texte ne l'entoure pas, il reste en ligne. Quel habillage choisir pour que le texte contourne l'image ?
  \item Donne le raccourci clavier universel (Word + LibreOffice) pour : Enregistrer, Justifier, Tout sélectionner, Correction orthographique.
\end{enumerate}
\end{exercice}

\begin{corrige}[Exercice Entraînement BEPC]
\begin{enumerate}
  \item (a) \textbf{Insertion} (En-tête) ou \textbf{Accueil} (Alignement Droite) + Tabulation. (b) \textbf{Accueil} (Gras, Centrer). (c) \textbf{Insertion} (Numéro de page) ou \textbf{Disposition} (Pied de page).
  \item Parce que la table des matières automatique ne repère \textbf{que} les paragraphes marqués par les styles de titres (Titre 1, 2, 3). Une mise en forme manuelle (Gras/Centré) est invisible pour la fonction TDM.
  \item \textbf{Habillage « Carré »} (ou « Contour » / « Haut et bas »). L'option par défaut « En ligne avec le texte » empêche le contournement.
  \item Enregistrer : \texttt{Ctrl+S} ; Justifier : \texttt{Ctrl+J} ; Tout sélectionner : \texttt{Ctrl+A} ; Correction : \texttt{F7}.
\end{enumerate}
\end{corrige}

\newpage

\coursec{Exercices}

\coursec{Leçon 05 : Utilisation d'un logiciel de traitement de texte}

\courssub{Je vérifie mes connaissances}

\begin{exercice}[QCM : les bons réflexes]
Pour chaque question, une seule réponse est correcte. Recopie la lettre correspondante.
\begin{enumerate}
\item Pour mettre un mot en \textbf{gras} dans un traitement de texte, on utilise le raccourci clavier :
\begin{enumerate}
\item[a)] Ctrl + G \quad b)] Ctrl + B \quad c)] Ctrl + S \quad d)] Ctrl + I
\end{enumerate}
\item Un texte \cle{justifié} est un texte :
\begin{enumerate}
\item[a)] aligné uniquement à gauche ;
\item[b)] aligné à la fois sur la marge gauche et la marge droite ;
\item[c)] centré sur la page ;
\item[d)] écrit en majuscules.
\end{enumerate}
\item Pour insérer une image dans un document, on utilise le menu (l'onglet) :
\begin{enumerate}
\item[a)] Accueil \quad b)] Affichage \quad c)] Insertion \quad d)] Mise en page
\end{enumerate}
\item La génération automatique d'une table des matières nécessite d'avoir appliqué au préalable :
\begin{enumerate}
\item[a)] des couleurs aux titres ;
\item[b)] des styles de titres (Titre 1, Titre 2, ...) ;
\item[c)] des bordures de page ;
\item[d)] des numéros de page manuels.
\end{enumerate}
\end{enumerate}
\end{exercice}

\begin{exercice}[Vrai ou faux ? Justifie ta réponse]
Réponds par Vrai ou Faux, puis justifie en une ou deux phrases.
\begin{enumerate}
\item Pour passer à la page suivante, il est conseillé d'appuyer plusieurs fois sur la touche Entrée.
\item Un \cle{saut de page} permet de commencer une nouvelle page proprement.
\item Une table des matières automatique peut être mise à jour après l'ajout d'un nouveau titre dans le document.
\item L'\cle{habillage} d'une image définit la manière dont le texte s'organise autour de cette image.
\item Il est impossible d'avoir, dans un même document, une page en orientation portrait et une autre en orientation paysage.
\end{enumerate}
\end{exercice}

\begin{exercice}[Définitions]
Définis clairement chacun des termes suivants en une ou deux phrases :
\begin{enumerate}
\item un \cle{style} (de titre) ;
\item la \cle{mise en forme} d'un texte ;
\item la \cle{mise en page} d'un document ;
\item un \cle{saut de section} ;
\item une \cle{table des matières}.
\end{enumerate}
\end{exercice}

\begin{exercice}[Questions de cours]
\begin{enumerate}
\item Cite quatre mises en forme que l'on peut appliquer à des caractères.
\item Cite trois mises en forme que l'on peut appliquer à un paragraphe.
\item Explique pourquoi il est déconseillé de taper une table des matières « à la main » dans un long document.
\item Donne deux objets que l'on peut insérer dans un document à l'aide de l'onglet Insertion.
\end{enumerate}
\end{exercice}

\courssub{Je m'entraîne}

\begin{exercice}[Identifier des mises en forme]
Le professeur d'informatique projette la capture d'écran d'un document contenant le texte suivant : le titre « \textbf{\large RENTRÉE SCOLAIRE 2026} » est centré, en gras, en taille 16 et souligné ; le paragraphe qui suit est justifié, en taille 12, avec un retrait de première ligne ; la dernière ligne « \emph{Le Proviseur} » est alignée à droite et en italique.
\begin{enumerate}
\item Pour chaque élément du document (titre, paragraphe, dernière ligne), identifie les mises en forme appliquées : alignement, style, taille.
\item Cite les boutons du ruban (onglet Accueil) qui ont permis d'obtenir chacune de ces mises en forme.
\end{enumerate}
\end{exercice}

\begin{exercice}[Reproduire une lettre administrative]
Le secrétariat du collège te remet une lettre tapée sur papier et te demande de la reproduire à l'identique sur l'ordinateur. La lettre comporte : un en-tête « Collège Bilingue de Mfoundi — BP 1234 Yaoundé » centré en haut ; la date alignée à droite ; un objet en gras et souligné ; un corps de texte justifié ; la signature « \textbf{Le Proviseur} » en gras, alignée à droite.
\begin{enumerate}
\item Règle les marges du document à 2,5 cm de chaque côté. Indique l'onglet et la commande utilisés.
\item Saisis le texte de la lettre et applique toutes les mises en forme demandées.
\item Décris, étape par étape, la démarche suivie pour obtenir l'en-tête centré et la date alignée à droite sur la même zone de la page.
\end{enumerate}
\end{exercice}

\begin{exercice}[Créer un tableau de notes]
On veut créer dans un document le tableau des notes de la séquence 1 d'une classe de 3ème, comportant 5 colonnes (N°, Nom et prénom, Maths, Français, Moyenne) et 8 lignes (1 ligne de titre, 1 ligne d'en-têtes de colonnes, 6 lignes d'élèves).
\begin{enumerate}
\item Indique la démarche pour insérer un tableau de 5 colonnes et 8 lignes.
\item Fusionne les cellules de la première ligne et tape le titre « Notes de la séquence 1 — Classe de 3ème A », centré et en gras. Décris la manipulation.
\item Applique des bordures à tout le tableau. Cite la commande utilisée.
\item Centre verticalement et horizontalement le contenu de la ligne d'en-têtes de colonnes.
\end{enumerate}
\end{exercice}

\begin{exercice}[Insérer et habiller une image]
Tu rédiges un exposé d'une page sur « Le Cameroun numérique » et tu veux y insérer la photo d'un centre de télécommunications.
\begin{enumerate}
\item Décris la démarche pour insérer une image enregistrée sur l'ordinateur.
\item Le texte doit s'écouler autour de l'image : quel réglage faut-il utiliser ? Donne son nom et cite deux types d'habillage possibles.
\item Explique comment redimensionner l'image sans la déformer.
\end{enumerate}
\end{exercice}

\courssub{Je me prépare au BEPC}

\begin{exercice}[Exercice type BEPC : analyse d'une capture d'écran]
À l'examen, on te présente la capture d'écran d'un document produit avec un logiciel de traitement de texte. Le document s'intitule « \textbf{\Large Charte de la vie scolaire} » : le titre est centré, en gras, en taille 18 et de couleur bleue ; le document contient deux paragraphes justifiés avec un interligne de 1,5 ; une liste à puces de quatre règles ; en bas de page figure le numéro de page « Page 1 / 3 ».
\begin{enumerate}
\item Définis les termes : mise en forme, mise en page.
\item Releve dans le document quatre mises en forme de caractères et trois mises en forme de paragraphes.
\item Pour chacune des mises en forme relevées, cite le bouton ou la commande du ruban qui permet de l'obtenir.
\item Explique comment le numéro de page « Page 1 / 3 » a été inséré automatiquement en bas de chaque page.
\item Le candidat affirme : « Pour centrer le titre, j'ai tapé des espaces devant le texte. » Explique pourquoi cette méthode est mauvaise et donne la méthode correcte.
\end{enumerate}
\end{exercice}

\begin{exercice}[Exercice type BEPC : table des matières automatique]
Pour les activités du club informatique, tu dois structurer un texte brut de 2 pages intitulé « Guide de la salle multimédia ». Le texte contient les titres suivants, actuellement tapés comme du texte ordinaire : « 1. Présentation de la salle », « 1.1 Les équipements », « 1.2 Les horaires d'ouverture », « 2. Règles d'utilisation », « 2.1 Avant la séance », « 2.2 Après la séance ».
\begin{enumerate}
\item Explique ce qu'est un style de titre et donne l'intérêt de l'utiliser.
\item Décris la démarche pour appliquer le style \cle{Titre 1} aux titres « 1. Présentation de la salle » et « 2. Règles d'utilisation », et le style \cle{Titre 2} aux quatre sous-titres.
\item Décris la démarche pour générer automatiquement la table des matières en première page du document. Cite l'onglet et la commande utilisés.
\item Tu ajoutes ensuite un nouveau paragraphe et un nouveau titre « 2.3 Sanctions ». Que faut-il faire pour que la table des matières tienne compte de cette modification ? Décris la manipulation.
\item Donne deux avantages d'une table des matières automatique par rapport à une table tapée à la main.
\end{enumerate}
\end{exercice}

\begin{exercice}[Situation d'intégration : le document du proviseur]
Le proviseur de ton collège prépare le rapport de fin d'année destiné au conseil d'administration. Le document comporte 6 pages de texte rédigé en orientation portrait. À la page 4, il souhaite insérer un grand tableau récapitulatif des effectifs par classe (12 colonnes) qui ne tient pas dans la largeur d'une page portrait : cette page 4 doit donc être en orientation \cle{paysage}, les pages 5 et 6 revenant en portrait. Il te confie la réalisation complète du document.
\begin{enumerate}
\item Explique pourquoi il est impossible de changer l'orientation d'une seule page en utilisant simplement la commande « Orientation » de l'onglet Mise en page.
\item Définis ce qu'est un saut de section et explique son rôle dans ce problème.
\item Décris, étape par étape, la démarche complète pour obtenir : pages 1 à 3 en portrait, page 4 en paysage, pages 5 et 6 en portrait. Précise le type de saut à insérer et les endroits où le placer.
\item Le proviseur demande en plus que le tableau de la page 4 ait : une première ligne fusionnée portant le titre « Effectifs 2026-2027 », des bordures sur toutes les cellules et la ligne d'en-têtes en gras et centrée. Décris les manipulations nécessaires.
\item Enfin, il veut une table des matières en page 1. Rédige la démarche complète à suivre, de l'application des styles jusqu'à la génération de la table, en justifiant chaque étape.
\end{enumerate}
\end{exercice}

\begin{aretenir}[Synthèse de la leçon : Produire un document complet]
Pour réaliser un document professionnel et bien structuré, on suit la démarche suivante :
\begin{enumerate}
\item \textbf{Préparer la mise en page} (onglet \texttt{Mise en page}) : définir les \textbf{marges}, l'\textbf{orientation} (portrait/paysage), la \textbf{taille du papier}. Pour des mises en page différentes dans un même document (ex. une page paysage au milieu de pages portrait), on insère des \textbf{sauts de section} (type \texttt{Page suivante}) avant et après la zone concernée, puis on modifie l'orientation \emph{uniquement} pour cette section.
\item \textbf{Structurer le contenu avec les styles} (onglet \texttt{Accueil}) : appliquer les styles \texttt{Titre 1}, \texttt{Titre 2}, etc. aux titres du document. C'est la condition \emph{indispensable} pour générer une table des matières automatique.
\item \textbf{Mettre en forme les caractères et paragraphes} (onglet \texttt{Accueil}) : police, taille, gras, italique, souligné, couleur (caractères) ; alignement, interligne, retraits, puces/numérotation, tabulations (paragraphes).
\item \textbf{Insérer les objets} (onglet \texttt{Insertion}) : tableaux (et les mettre en forme : fusion, bordures, alignement cellule), images (et gérer l'\textbf{habillage} : \texttt{Au travers}, \texttt{Carré}, \texttt{Derrière le texte}...), en-têtes/pieds de page (numérotation automatique : \texttt{Numéro de page} $\to$ \texttt{Bas de page}), sauts de page.
\item \textbf{Générer la table des matières} (onglet \texttt{Références} $\to$ \texttt{Table des matières}) : choisir un style automatique. En cas de modification du document : clic droit sur la table $\to$ \texttt{Mettre à jour les champs} $\to$ \texttt{Mettre à jour toute la table}.
\item \textbf{Finaliser et vérifier} : aperçu avant impression, vérification orthographique, enregistrement régulier (Ctrl + S).
\end{enumerate}
\end{aretenir}

\newpage

\coursec{Corrigés des exercices}

\begin{corrige}[Exercice 1 — QCM : les bons réflexes]
\begin{enumerate}
\item \textbf{Réponse a) : Ctrl + G.} Dans les logiciels de traitement de texte utilisés au Cameroun (Microsoft Word, LibreOffice Writer), le raccourci \texttt{Ctrl + G} applique ou retire le \cle{gras} à la sélection. (Dans les versions à interface anglaise, c'est \texttt{Ctrl + B}, de l'anglais \emph{Bold}.) \texttt{Ctrl + S} sert à \emph{enregistrer} (Sauvegarder) et \texttt{Ctrl + I} à mettre en \emph{italique}.
\item \textbf{Réponse b).} Un texte \cle{justifié} est aligné à la fois sur la marge gauche et sur la marge droite : le logiciel ajuste automatiquement les espaces entre les mots pour que toutes les lignes aient la même longueur.
\item \textbf{Réponse c) : Insertion.} L'onglet \texttt{Insertion} regroupe toutes les commandes permettant d'ajouter des objets au document : images, tableaux, formes, en-têtes et pieds de page, sauts de page, symboles, etc.
\item \textbf{Réponse b).} La table des matières automatique est construite par le logiciel à partir des \cle{styles de titres} (\texttt{Titre 1}, \texttt{Titre 2}, ...) appliqués dans le document. Sans styles, le logiciel ne peut pas repérer les titres et la génération est impossible.
\end{enumerate}
\end{corrige}

\begin{corrige}[Exercice 2 — Vrai ou faux ?]
\begin{enumerate}
\item \textbf{Faux.} Appuyer plusieurs fois sur \texttt{Entrée} crée des paragraphes vides : si l'on ajoute ou supprime du texte plus haut, ces lignes vides se déplacent et la mise en page est cassée. La méthode correcte est d'insérer un \cle{saut de page}.
\item \textbf{Vrai.} Le saut de page (\texttt{Insertion} $\to$ \texttt{Saut de page}, ou \texttt{Ctrl + Entrée}) force le début d'une nouvelle page à un endroit précis, et ce repère reste stable même si le texte précédent est modifié.
\item \textbf{Vrai.} Après l'ajout d'un titre (auquel on a appliqué un style de titre), il suffit de faire un clic droit sur la table des matières $\to$ \texttt{Mettre à jour les champs} $\to$ \texttt{Mettre à jour toute la table} : le nouveau titre et son numéro de page apparaissent automatiquement.
\item \textbf{Vrai.} L'\cle{habillage} (ou \emph{renvoi à la ligne automatique}) détermine comment le texte s'organise autour de l'image : \texttt{Carré} (le texte contourne l'image en rectangle), \texttt{Au travers} (le texte épouse les contours), \texttt{Derrière le texte}, \texttt{Devant le texte}, etc.
\item \textbf{Faux.} C'est possible grâce aux \cle{sauts de section} : on insère un saut de section (page suivante) avant la page concernée et un autre après, puis on règle l'orientation paysage uniquement pour cette section. C'est exactement la situation du rapport du proviseur (exercice 11).
\end{enumerate}
\end{corrige}

\begin{corrige}[Exercice 3 — Définitions]
\begin{enumerate}
\item Un \cle{style} (de titre) est un ensemble de mises en forme prédéfinies (police, taille, couleur, gras, espacements) portant un nom (\texttt{Titre 1}, \texttt{Titre 2}, ...) que l'on applique d'un seul clic à un titre. Il permet d'uniformiser la présentation et sert de repère au logiciel pour construire la table des matières.
\item La \cle{mise en forme} d'un texte est l'ensemble des modifications d'aspect appliquées aux \emph{caractères} (police, taille, gras, italique, souligné, couleur, surlignage) et aux \emph{paragraphes} (alignement, interligne, retraits, puces, numérotation).
\item La \cle{mise en page} d'un document concerne la présentation globale des pages : marges, orientation (portrait ou paysage), taille du papier, en-têtes et pieds de page, numérotation des pages, colonnes.
\item Un \cle{saut de section} est une marque invisible qui découpe le document en sections indépendantes, chacune pouvant avoir sa propre mise en page (orientation, marges, en-têtes/pieds de page, numérotation...). On l'insère par \texttt{Mise en page} $\to$ \texttt{Sauts de page} $\to$ \texttt{Page suivante} (dans la catégorie \emph{Sauts de section}).
\item Une \cle{table des matières} est la liste des titres d'un document accompagnés de leurs numéros de page, placée généralement au début. Générée automatiquement à partir des styles de titres, elle peut être mise à jour en un clic.
\end{enumerate}
\end{corrige}

\begin{corrige}[Exercice 4 — Questions de cours]
\begin{enumerate}
\item Quatre mises en forme de \textbf{caractères} (au choix) : le gras, l'italique, le souligné, la taille des caractères, la police, la couleur du texte, la couleur de surlignage, l'exposition, les petites capitales.
\item Trois mises en forme de \textbf{paragraphes} (au choix) : l'alignement (gauche, centré, droit, justifié), l'interligne (simple, 1,5, double), le retrait de première ligne, les puces ou la numérotation, l'espacement avant/après le paragraphe, le trait de bordure.
\item Taper une table des matières « à la main » est déconseillé car, dans un long document, les numéros de page changent dès que l'on ajoute ou supprime du texte : la table devient alors fausse et il faut tout refaire. De plus, on risque des fautes de frappe dans les titres ou les numéros. La table automatique, elle, se met à jour en un clic et reste toujours exacte.
\item Deux objets insérables via l'onglet \texttt{Insertion} (au choix) : une image, un tableau, une forme, un en-tête ou pied de page, un numéro de page, un saut de page, un symbole, un WordArt, une zone de texte.
\end{enumerate}
\end{corrige}

\begin{corrige}[Exercice 5 — Identifier des mises en forme]
\begin{enumerate}
\item Mises en forme relevées sur le document :
\begin{itemize}
\item \textbf{Titre « RENTRÉE SCOLAIRE 2026 »} : alignement centré ; gras ; souligné ; taille 16.
\item \textbf{Paragraphe} : alignement justifié ; taille 12 ; retrait de première ligne.
\item \textbf{Dernière ligne « Le Proviseur »} : alignement à droite ; italique.
\end{itemize}
\item Boutons du ruban (onglet \texttt{Accueil}) utilisés :
\begin{itemize}
\item Titre : bouton \texttt{Centré} (groupe Paragraphe), bouton \texttt{G} (Gras), bouton \texttt{S} (Souligné), liste \texttt{Taille de police} réglée sur 16 (groupe Police).
\item Paragraphe : bouton \texttt{Justifié}, liste \texttt{Taille de police} sur 12, et pour le retrait de première ligne : boîte de dialogue \texttt{Paragraphe} $\to$ \texttt{Retrait} $\to$ \texttt{Première ligne} (ou manipulation sur la règle horizontale).
\item Dernière ligne : bouton \texttt{Aligner à droite} et bouton \texttt{I} (Italique).
\end{itemize}
\end{enumerate}
\textbf{Point de vigilance :} le retrait de première ligne ne s'obtient pas en tapant des espaces, mais avec la commande de retrait du paragraphe (ou la règle), afin qu'il soit identique sur tout le document.
\end{corrige}

\begin{corrige}[Exercice 6 — Reproduire une lettre administrative]
\begin{enumerate}
\item \textbf{Réglage des marges :} onglet \texttt{Mise en page} $\to$ bouton \texttt{Marges} $\to$ \texttt{Marges personnalisées} ; on saisit 2,5 cm pour les marges haute, basse, gauche et droite, puis on valide par \texttt{OK}. (On peut aussi choisir le préréglage \texttt{Normal} qui vaut déjà 2,5 cm dans Word et LibreOffice.)
\item \textbf{Saisie et mises en forme :}
\begin{enumerate}
\item Taper l'en-tête, le sélectionner, cliquer sur \texttt{Centré}.
\item Taper la date, la sélectionner, cliquer sur \texttt{Aligner à droite}.
\item Taper l'objet, le sélectionner, cliquer sur \texttt{G} (Gras) puis \texttt{S} (Souligné).
\item Sélectionner le corps de la lettre, cliquer sur \texttt{Justifié}.
\item Taper « Le Proviseur », le sélectionner, cliquer sur \texttt{G} et \texttt{Aligner à droite}.
\end{enumerate}
\item \textbf{Démarche pour l'en-tête centré et la date à droite :} l'en-tête et la date sont deux \emph{paragraphes différents}, chacun ayant son propre alignement.
\begin{enumerate}
\item Taper l'en-tête « Collège Bilingue de Mfoundi — BP 1234 Yaoundé ».
\item Placer le curseur dans ce paragraphe (ou le sélectionner) et cliquer sur le bouton \texttt{Centré} de l'onglet \texttt{Accueil}.
\item Appuyer sur \texttt{Entrée} pour créer un nouveau paragraphe, taper la date « Yaoundé, le ... ».
\item Placer le curseur dans ce paragraphe et cliquer sur \texttt{Aligner à droite}.
\end{enumerate}
Chaque paragraphe garde son alignement propre : c'est pourquoi on ne doit pas utiliser d'espaces pour « pousser » la date vers la droite.
\end{enumerate}
\end{corrige}

\begin{corrige}[Exercice 7 — Créer un tableau de notes]
\begin{enumerate}
\item \textbf{Insertion du tableau :} onglet \texttt{Insertion} $\to$ bouton \texttt{Tableau} $\to$ \texttt{Insérer un tableau} ; dans la boîte de dialogue, saisir \textbf{5} pour le nombre de colonnes et \textbf{8} pour le nombre de lignes, puis valider. (On peut aussi utiliser la grille visuelle en sélectionnant 5 colonnes $\times$ 8 lignes.)
\item \textbf{Fusion de la première ligne pour le titre :}
\begin{enumerate}
\item Sélectionner les 5 cellules de la première ligne (cliquer-glisser dessus).
\item Onglet contextuel \texttt{Disposition} (Outils de tableau) $\to$ \texttt{Fusionner les cellules} (ou clic droit $\to$ \texttt{Fusionner les cellules}).
\item Taper le titre « Notes de la séquence 1 — Classe de 3ème A ».
\item Sélectionner ce texte et cliquer sur \texttt{Centré} et \texttt{G} (Gras) dans l'onglet \texttt{Accueil}.
\end{enumerate}
\item \textbf{Bordures :} sélectionner tout le tableau $\to$ onglet \texttt{Création} (Outils de tableau) $\to$ bouton \texttt{Bordures} $\to$ \texttt{Toutes les bordures}.
\item \textbf{Centrage de la ligne d'en-têtes :} sélectionner la ligne d'en-têtes (N°, Nom et prénom, Maths, Français, Moyenne) ; pour le centrage horizontal, cliquer sur \texttt{Centré} (onglet \texttt{Accueil}) ; pour le centrage vertical, onglet \texttt{Disposition} $\to$ groupe \texttt{Alignement} $\to$ choisir l'icône \texttt{Centrer} (alignement centré horizontal et vertical de la cellule).
\end{enumerate}
\end{corrige}

\begin{corrige}[Exercice 8 — Insérer et habiller une image]
\begin{enumerate}
\item \textbf{Insertion de l'image :} placer le curseur à l'endroit voulu $\to$ onglet \texttt{Insertion} $\to$ bouton \texttt{Images} (ou \texttt{Image à partir d'un fichier}) $\to$ dans la boîte de dialogue, choisir le dossier, sélectionner le fichier image $\to$ cliquer sur \texttt{Insérer}.
\item \textbf{Réglage de l'habillage :} cliquer sur l'image pour la sélectionner $\to$ onglet contextuel \texttt{Format} $\to$ bouton \texttt{Habillage du texte} (ou l'icône d'habillage qui apparaît près de l'image). Le réglage qui permet au texte de s'écouler autour de l'image s'appelle l'\cle{habillage} ; deux types principaux permettent cet écoulement : \texttt{Carré} (le texte contourne l'image en formant un rectangle) et \texttt{Au travers} (le texte épouse au plus près les contours de l'image). On peut aussi citer \texttt{Rapproché}, \texttt{Derrière le texte}, \texttt{Devant le texte}.
\item \textbf{Redimensionnement sans déformation :} sélectionner l'image, puis tirer l'une des \textbf{poignées d'angle} (les petits cercles aux quatre coins) tout en maintenant éventuellement la touche \texttt{Maj} enfoncée : les proportions (largeur/hauteur) sont conservées. Il ne faut \emph{jamais} tirer les poignées du milieu des côtés, qui étirent l'image dans un seul sens et la déforment.
\end{enumerate}
\end{corrige}

\begin{corrige}[Exercice 9 — Type BEPC : analyse d'une capture d'écran]
\begin{enumerate}
\item \textbf{Définitions.} La \cle{mise en forme} est l'ensemble des modifications d'aspect appliquées aux caractères (police, taille, gras, couleur...) et aux paragraphes (alignement, interligne, puces...). La \cle{mise en page} concerne la présentation globale des pages : marges, orientation, taille du papier, en-têtes et pieds de page, numérotation.
\item \textbf{Mises en forme relevées sur la capture.}
Quatre mises en forme de \emph{caractères} : gras (titre), taille 18 (titre), couleur bleue (titre), police du titre (ou taille 12 du corps de texte si visible).
Trois mises en forme de \emph{paragraphes} : titre centré ; paragraphes justifiés ; interligne 1,5 ; liste à puces.
\item \textbf{Boutons et commandes correspondants (onglet Accueil) :}
\begin{itemize}
\item Gras : bouton \texttt{G} ou \texttt{Ctrl + G}.
\item Taille 18 : liste \texttt{Taille de police}.
\item Couleur bleue : bouton \texttt{Couleur de police}.
\item Titre centré : bouton \texttt{Centré} (groupe Paragraphe).
\item Paragraphes justifiés : bouton \texttt{Justifié}.
\item Interligne 1,5 : bouton \texttt{Interligne} $\to$ \texttt{1,5}.
\item Liste à puces : bouton \texttt{Puces}.
\end{itemize}
\item \textbf{Numéro de page automatique :} onglet \texttt{Insertion} $\to$ \texttt{Pied de page} (ou \texttt{Numéro de page} $\to$ \texttt{Bas de page}) $\to$ choisir un modèle affichant le numéro, par exemple « Page 1 / 3 » (numéro de page sur nombre total de pages). Le logiciel insère un \emph{champ automatique} dans le pied de page : le numéro se calcule tout seul et s'affiche correctement sur chacune des 3 pages.
\item \textbf{Pourquoi les espaces sont une mauvaise méthode :} centrer avec des espaces est imprécis (le résultat dépend de la police et de la taille) et surtout fragile : dès que l'on modifie le titre, la taille ou les marges, le « centrage » est perdu. La méthode correcte : placer le curseur dans le titre et cliquer sur le bouton \texttt{Centré} de l'onglet \texttt{Accueil} (ou \texttt{Ctrl + E}) : le titre reste parfaitement centré quelles que soient les modifications.
\end{enumerate}
\textbf{Barème indicatif (10 pts) :} définitions 2 pts ; relevé des mises en forme 2 pts ; boutons/commandes 3 pts ; numéro de page 1,5 pt ; justification du centrage 1,5 pt.
\textbf{Points de vigilance :} bien distinguer mise en \emph{forme} (caractères/paragraphes) et mise en \emph{page} (marges, orientation) ; citer le nom exact des boutons ; pour la question 5, exiger la justification (fragilité de la méthode) et pas seulement la méthode correcte.
\end{corrige}

\begin{corrige}[Exercice 10 — Type BEPC : table des matières automatique]
\begin{enumerate}
\item Un \cle{style de titre} est un ensemble de mises en forme prédéfinies (police, taille, gras, couleur, espacements) appliqué d'un clic à un titre. Intérêts : les titres sont présentés de façon \emph{uniforme} dans tout le document ; on peut modifier l'aspect de tous les titres d'un coup en modifiant le style ; et surtout les styles permettent au logiciel de \emph{reconnaître} les titres pour générer automatiquement la table des matières.
\item \textbf{Application des styles :}
\begin{enumerate}
\item Sélectionner le texte « 1. Présentation de la salle ».
\item Dans l'onglet \texttt{Accueil}, groupe \texttt{Styles}, cliquer sur \texttt{Titre 1}.
\item Recommencer pour « 2. Règles d'utilisation » avec \texttt{Titre 1}.
\item Sélectionner chacun des quatre sous-titres (« 1.1 Les équipements », « 1.2 Les horaires d'ouverture », « 2.1 Avant la séance », « 2.2 Après la séance ») et leur appliquer le style \texttt{Titre 2}.
\end{enumerate}
Le style \texttt{Titre 2} indique un niveau inférieur : les sous-titres apparaîtront décalés (en retrait) dans la table des matières.
\item \textbf{Génération de la table des matières :}
\begin{enumerate}
\item Placer le curseur en première page, à l'endroit où la table doit apparaître.
\item Onglet \texttt{Références} $\to$ bouton \texttt{Table des matières}.
\item Choisir un modèle de table automatique proposé.
\end{enumerate}
Le logiciel parcourt le document, repère tous les textes en styles \texttt{Titre 1} et \texttt{Titre 2}, et construit la table avec les numéros de page.
\item \textbf{Mise à jour après ajout du titre « 2.3 Sanctions » :} appliquer d'abord le style \texttt{Titre 2} au nouveau titre, puis cliquer sur la table des matières $\to$ bouton \texttt{Mettre à jour la table} (ou clic droit $\to$ \texttt{Mettre à jour les champs}) $\to$ choisir \texttt{Mettre à jour toute la table} (pour prendre en compte le nouveau titre et pas seulement les numéros de page).
\item \textbf{Deux avantages de la table automatique :}
\begin{itemize}
\item Elle se met à jour en un clic après toute modification (titres ajoutés, pages déplacées), sans risque d'erreur de numérotation.
\item Elle est construite sans faute de frappe (titres recopiés exactement) et, dans les documents numériques, ses entrées sont cliquables pour naviguer directement vers les chapitres (Ctrl + Clic).
\end{itemize}
\end{enumerate}
\textbf{Barème indicatif (10 pts) :} définition et intérêt du style 2 pts ; application des styles 2 pts ; génération de la table 2 pts ; mise à jour 2 pts ; avantages 2 pts.
\textbf{Points de vigilance :} la condition indispensable à citer est l'application préalable des \emph{styles} de titres ; pour la mise à jour, préciser « Mettre à jour \emph{toute} la table » car un nouveau titre est ajouté (et pas seulement des pages modifiées).
\end{corrige}

\begin{corrige}[Exercice 11 — Situation d'intégration : le document du proviseur]
\begin{enumerate}
\item La commande \texttt{Orientation} de l'onglet \texttt{Mise en page} s'applique par défaut à \textbf{tout le document} : si le document ne contient qu'une seule section, passer en paysage transforme les 6 pages. Il est donc impossible de changer l'orientation d'une seule page tant que le document n'est pas découpé en sections.
\item Un \cle{saut de section} est une marque invisible qui divise le document en sections indépendantes, chacune pouvant posséder sa propre mise en page (orientation, marges, en-têtes/pieds de page, numérotation...). Ici, son rôle est d'isoler la page 4 dans une section à part, afin de pouvoir la passer en paysage sans toucher aux pages 1--3 et 5--6.
\item \textbf{Démarche complète :}
\begin{enumerate}
\item Placer le curseur \emph{à la fin de la page 3} (juste avant le début du contenu de la page 4).
\item Onglet \texttt{Mise en page} $\to$ \texttt{Sauts de page} $\to$ \texttt{Sauts de section} $\to$ \texttt{Page suivante}. La page 4 commence alors une nouvelle section (section 2).
\item Placer le curseur \emph{à la fin de la page 4} (avant le contenu de la page 5) et insérer un second saut de section \texttt{Page suivante} : les pages 5--6 forment la section 3.
\item Placer le curseur n'importe où dans la page 4 (section 2) $\to$ onglet \texttt{Mise en page} $\to$ \texttt{Orientation} $\to$ \texttt{Paysage}. Seule la section 2 change d'orientation.
\item Vérifier dans l'aperçu avant impression : pages 1 à 3 en portrait, page 4 en paysage, pages 5 et 6 en portrait.
\end{enumerate}
\item \textbf{Mise en forme du tableau de la page 4 :}
\begin{enumerate}
\item Insérer le tableau (\texttt{Insertion} $\to$ \texttt{Tableau}) avec 12 colonnes et le nombre de lignes voulu.
\item Sélectionner toutes les cellules de la première ligne $\to$ onglet \texttt{Disposition} $\to$ \texttt{Fusionner les cellules} $\to$ taper le titre « Effectifs 2026-2027 » $\to$ le centrer et le mettre en gras.
\item Sélectionner tout le tableau $\to$ onglet \texttt{Création} $\to$ \texttt{Bordures} $\to$ \texttt{Toutes les bordures}.
\item Sélectionner la ligne d'en-têtes des colonnes $\to$ appliquer \texttt{G} (Gras) et \texttt{Centré} ; pour un centrage parfait dans les cellules, utiliser l'alignement centré horizontal et vertical du groupe \texttt{Alignement} (onglet \texttt{Disposition}).
\end{enumerate}
\item \textbf{Table des matières en page 1 :}
\begin{enumerate}
\item Parcourir le document et appliquer le style \texttt{Titre 1} aux grands titres du rapport et le style \texttt{Titre 2} aux sous-titres (\emph{justification : c'est la condition indispensable pour que le logiciel repère les titres}).
\item Placer le curseur en page 1, à l'emplacement prévu pour la table (éventuellement après un saut de page pour isoler la table).
\item Onglet \texttt{Références} $\to$ \texttt{Table des matières} $\to$ choisir une table automatique : la table est générée avec les numéros de page exacts.
\item Après toute modification ultérieure du rapport, mettre à jour la table : clic droit sur la table $\to$ \texttt{Mettre à jour les champs} $\to$ \texttt{Mettre à jour toute la table} (\emph{justification : les numéros de page évoluent avec le contenu}).
\end{enumerate}
\end{enumerate}
\textbf{Barème indicatif (20 pts) :} question 1 : 2 pts ; question 2 : 3 pts ; question 3 : 6 pts (2 pts par saut correctement placé, 2 pts pour l'orientation de la section 2) ; question 4 : 4 pts ; question 5 : 5 pts.
\textbf{Points de vigilance :} exiger la mention du type de saut (\texttt{Saut de section -- Page suivante}) et de ses deux emplacements (fin de page 3 et fin de page 4) ; rappeler que l'orientation se règle \emph{curseur placé dans la section concernée} ; pour la table des matières, la justification « les styles de titres sont indispensables » doit apparaître explicitement.
\end{corrige}

\newpage

\coursec{Fiche bilan}

\begin{popfigure}
\centering
\begin{tikzpicture}[
  mindmap,
  grow cyclic,
  every node/.style={concept, execute at begin node=\hskip0pt},
  concept color=popblue,
  level 1/.style={level distance=4.5cm, sibling angle=360/6},
  level 2/.style={level distance=3cm, sibling angle=45},
  text=white,
  font=\small\bfseries
]
\node[concept, scale=1.2] {Traitement de texte}
  child[concept color=poporange] { node[concept] {Interface \& Base} }
  child[concept color=popgreen] { node[concept] {Mise en forme} }
  child[concept color=poppurple] { node[concept] {Mise en page} }
  child[concept color=poppink] { node[concept] {Insertion d'objets} }
  child[concept color=popgold] { node[concept] {Styles \& TDM auto} }
  child[concept color=popdark] { node[concept] {Fichiers \& Impression} };
\end{tikzpicture}
\legende{Carte mentale : les 6 piliers de la leçon sur le traitement de texte (classe de 3ème).}
\end{popfigure}

\begin{bilan}

\textbf{Définitions clés}
\begin{itemize}
  \item \cle{Traitement de texte} : Logiciel permettant la création, la modification, la mise en forme et l'impression de documents textuels (ex. : LibreOffice Writer, MS Word).
  \item \cle{Mise en forme} : Modification de l'apparence des caractères (police, taille, gras, couleur) et des paragraphes (alignement, interligne, retraits, puces).
  \item \cle{Mise en page} : Définition de la structure physique de la page : marges, orientation (portrait/paysage), taille (A4), en-têtes/pieds de page, numérotation, sauts de page.
  \item \cle{Style} : Ensemble de formats (caractère + paragraphe) enregistré sous un nom (ex. : « Titre 1 », « Corps de texte ») pour uniformiser et automatiser la structure du document.
  \item \cle{Table des matières (TDM) automatique} : Liste générée dynamiquement à partir des styles de titres appliqués ; se met à jour sur commande (clic droit \rightarrow Mettre à jour).
\end{itemize}

\textbf{Raccourcis clavier indispensables (à connaître par cœur)}
\begin{center}
\begin{tabularx}{\linewidth}{|l|X|l|}
\hline
\rowcolor{popblueL}
\textbf{Action} & \textbf{Description} & \textbf{Raccourci (Win/Linux)} \\
\hline
Enregistrer & Sauvegarder le document & \texttt{Ctrl + S} \\
\hline
Annuler / Rétablir & Revenir en arrière / Refaire & \texttt{Ctrl + Z} / \texttt{Ctrl + Y} \\
\hline
Copier / Couper / Coller & Dupliquer / Déplacer la sélection & \texttt{Ctrl + C} / \texttt{Ctrl + X} / \texttt{Ctrl + V} \\
\hline
Sélectionner tout & Sélectionner tout le document & \texttt{Ctrl + A} \\
\hline
Gras / Italique / Souligné & Mise en forme rapide caractères & \texttt{Ctrl + B} / \texttt{Ctrl + I} / \texttt{Ctrl + U} \\
\hline
Alignement G / C / D / J & Paragraphe : Gauche, Centré, Droite, Justifié & \texttt{Ctrl + L} / \texttt{Ctrl + E} / \texttt{Ctrl + R} / \texttt{Ctrl + J} \\
\hline
Saut de page forcé & Nouvelle page propre (pas de retours ligne) & \texttt{Ctrl + Entrée} \\
\hline
Rechercher / Remplacer & Trouver un mot, le substituer & \texttt{Ctrl + F} / \texttt{Ctrl + H} \\
\hline
Mettre à jour TDM & Actualiser la table des matières & \texttt{F9} (ou clic droit \rightarrow Mettre à jour) \\
\hline
\end{tabularx}
\end{center}

\textbf{Méthodes express}
\begin{itemize}
  \item \textbf{Appliquer un style de titre :} Sélectionner la ligne \rightarrow Menu \texttt{Styles} (volet latéral ou barre) \rightarrow Choisir « Titre 1 », « Titre 2 », etc.
  \item \textbf{Insérer une image bien ancrée :} \texttt{Insertion \rightarrow Image} \rightarrow Clic droit sur l'image \rightarrow \texttt{Ancrage \rightarrow Au paragraphe} (ou « En caractère » pour déplacement fluide).
  \item \textbf{Créer une TDM auto :} 1. Appliquer styles « Titre 1/2/3 » sur tous les titres. 2. Placer curseur au début. 3. \texttt{Insertion \rightarrow Table des matières et index \rightarrow Table des matières} \rightarrow OK.
  \item \textbf{Numérotation pages (sans page 1) :} \texttt{Insertion \rightarrow En-tête et pied de page \rightarrow Pied de page} \rightarrow Cocher « Première page différente » \rightarrow Insérer numéro page.
\end{itemize}

\end{bilan}

\begin{aretenir}[Checklist avant le BEPC]
$\square$ Je sais identifier les zones de l'interface : barre de menus, barre d'outils, règle, barre d'état, zone de saisie.
$\square$ Je maîtrise les raccourcis \texttt{Ctrl+S}, \texttt{Ctrl+Z}, \texttt{Ctrl+C/V/X}, \texttt{Ctrl+A}, \texttt{Ctrl+B/I/U}.
$\square$ Je sais mettre en forme des caractères (police, taille, couleur, exposant/indice) et des paragraphes (alignement, interligne 1,5, retraits, puces/numérotation).
$\square$ Je sais configurer la page : marges (normales/étroites), orientation (portrait/paysage), format A4, colonnes.
$\square$ Je sais insérer et paramétrer : image (redimensionner, ancrage, habillage), tableau (créer, fusionner/séparer cellules, bordures), forme, zone de texte, saut de page (\texttt{Ctrl+Entrée}).
$\square$ Je sais créer et modifier un en-tête / pied de page avec numérotation de pages (y compris « première page différente »).
$\square$ J'applique les styles de titres (Titre 1, Titre 2, Titre 3) au lieu de formater manuellement.
$\square$ Je sais générer une Table des Matières automatique et la mettre à jour (\texttt{F9} ou clic droit) après modification.
$\square$ Je connais la différence entre « Enregistrer » (écrase) et « Enregistrer sous » (nouveau nom/emplacement/format .odt, .docx, .pdf).
$\square$ Je sais exporter en PDF pour figer la mise en page avant impression ou envoi.
$\square$ Je repère les erreurs fréquentes : retours à la ligne multiples au lieu de saut de page, espaces multiples pour aligner au lieu de tabulations/alignement, formatage direct au lieu de styles.
$\square$ Je sais utiliser « Rechercher / Remplacer » (\texttt{Ctrl+H}) pour corriger un mot répété dans tout le document.
\end{aretenir}

\findecours

\end{document}
